\documentclass[11pt]{article}
\usepackage[margin=.88in]{geometry}
\usepackage[T1]{fontenc}
\usepackage[utf8]{inputenc}
\usepackage{lmodern,microtype}
\usepackage{amsmath,amssymb,amsthm,mathtools}
\usepackage{booktabs,longtable,array,enumitem}
\usepackage{graphicx,float}
\usepackage{algorithm,algpseudocode}
\usepackage{xcolor,hyperref,listings}
\hypersetup{
  colorlinks=true,
  linkcolor=blue!50!black,
  urlcolor=blue!50!black,
  citecolor=blue!50!black,
  pdftitle={Resolution Theory - Universal Representation, Nonlinear Transfer, Logic, and Mathematical Practice},
  pdfauthor={Komorebi P.}
}
\setlength{\parindent}{0pt}
\setlength{\parskip}{6pt}
\emergencystretch=2em
\newcolumntype{L}[1]{>{\raggedright\arraybackslash}p{#1}}
\newtheorem{theorem}{Theorem}[section]
\newtheorem{proposition}[theorem]{Proposition}
\newtheorem{corollary}[theorem]{Corollary}
\newtheorem{lemma}[theorem]{Lemma}
\theoremstyle{definition}
\newtheorem{definition}[theorem]{Definition}
\newtheorem{example}[theorem]{Example}
\theoremstyle{remark}
\newtheorem{remark}[theorem]{Remark}
\newtheorem{conjecture}[theorem]{Conjecture}
\newcommand{\Path}{\Pi}
\newcommand{\Dfield}{\mathcal D}
\newcommand{\G}{\mathcal G}
\newcommand{\Lset}{\mathcal L}
\newcommand{\I}{\mathcal I}
\newcommand{\Pspace}{\mathcal P}
\newcommand{\Mspace}{\mathcal M}
\newcommand{\Qspace}{\mathcal Q}
\newcommand{\Nspace}{\mathcal N}
\newcommand{\R}{\mathbb R}
\newcommand{\N}{\mathbb N}
\DeclareMathOperator{\Comp}{Comp}
\DeclareMathOperator{\Lost}{Lost}
\DeclareMathOperator{\Bridge}{Bridge}
\DeclareMathOperator{\Depth}{Depth}
\DeclareMathOperator{\Cost}{Cost}
\DeclareMathOperator{\Noise}{Noise}
\DeclareMathOperator{\Var}{Var}
\lstset{
  basicstyle=\ttfamily\scriptsize,
  breaklines=true,
  showstringspaces=false,
  frame=single,
  columns=fullflexible,
  tabsize=4
}
\title{Appendix: Resolution Theory\\
\large Universal Representation, Nonlinear Transfer, Logic, and Mathematical Practice}
\author{Komorebi P.}
\date{}
\begin{document}
\maketitle
\begin{center}
\Large\itshape
What do mathematics, computer science, statistics, or even technology mean to you?
\end{center}

\vspace{0.4em}

\begin{abstract}
Resolution Theory begins from a practical observation: mathematical work is never carried out without a representation.  We choose coordinates, bases, partitions, scales, models, proof forms, algorithms, and instruments.  Those choices reveal some structure and hide other structure.  This appendix keeps the choice of resolution inside the mathematics and follows that decision through four connected layers.  Resolution geometry studies representations, refinement, compression, and reconstruction.  Nonlinear Resolution Theory records interaction through jets, hyperedges, and composition trees.  Resolution Logic carries truth, error, cost, failure probability, and evidence in one judgment.  Proof-path methodology studies the human life of these objects: how ideas are discovered, compressed, taught, reviewed, archived, and reopened.  The final section brings the four layers together and asks what mathematics, computer science, statistics, and technology become when representation, residuals, and evidence are no longer treated as invisible scaffolding.  The finite theory is given through axioms, proofs, worked examples, Python implementations, numerical checks, and automated tests.  The ordinary abstract-to-concrete representation problem is settled here for unital resolved $C^*$-systems.  Countable and unbounded categorical completion, general topos semantics, and monoidal realization without an assumed fiber functor remain open programs.
\end{abstract}
\tableofcontents
\clearpage
\section*{A question before the formalism}
\addcontentsline{toc}{section}{A question before the formalism}

What do mathematics, computer science, statistics, or even technology mean to you?  The question sounds personal, but it points to a precise structural problem.  Mathematics is often described as the study of patterns or proof.  Computer science is described as the study of computation.  Statistics is described as inference under uncertainty.  Technology is described as the practical use of knowledge.  Those descriptions are reasonable, but they leave out the common act that precedes all four: someone decides how the situation will be represented.

A theorem is written in one language rather than another.  A program stores one data structure rather than another.  A statistical model groups observations in one way rather than another.  A device measures one collection of variables rather than another.  These are not cosmetic decisions.  They determine which relations are visible, which operations are cheap, which errors can be measured, and which conclusions can be justified.

The proposal developed here is that a mature theory of representation should keep five things visible:
\[
\boxed{
\text{resolution}
+\text{transfer}
+\text{residual}
+\text{evidence}
+\text{cost}.
}
\]
The first three belong to the mathematical object and its approximation.  The last two belong to the way a conclusion is reached and used.  The purpose of the appendix is not to turn every branch of knowledge into the same subject.  It is to give those branches a common language for saying what was represented, what moved, what was omitted, what supports the conclusion, and what the representation cost.

\begin{figure}[H]
\centering
\includegraphics[width=.92\linewidth]{figures/synthesis/01_four_fields_one_architecture.png}
\caption{Mathematics, computer science, statistics, and technology ask different questions, but each depends on a chosen representation, a transfer rule, an evidence standard, and a cost.}
\end{figure}

\clearpage
\section{Resolution Theory: when representation becomes part of the mathematics}
\label{app:resolution-theory}

\subsection{The question behind the theory}

Mathematics is almost never practiced without first deciding how an object will be seen.  We choose coordinates, a basis, a scale, a cover, a filtration, a partition, or a family of observables.  Once the argument is finished, that choice is often treated as disposable scaffolding.  Sometimes it really is harmless.  Sometimes it is where the structure was hiding.

The proposal of Resolution Theory is simple: do not throw the resolution away after it has done its work.  Record it as part of the object being studied.  A resolution says how the object has been divided into packets.  A transfer field says how a map moves information between those packets.  Refinement opens one packet into smaller ones.  Compression keeps only selected transfers and records what was left out.  Reconstruction returns the resolved description to the original object.

The finite-dimensional model is intentionally plain.  Let
\[
\mathcal H=\bigoplus_{i\in I}\mathcal H_i,
\qquad
P_iP_j=\delta_{ij}P_i,
\qquad
\sum_iP_i=I_\mathcal H.
\]
For a second space
\[
\mathcal K=\bigoplus_{a\in A}\mathcal K_a
\]
with projectors \(Q_a\), every linear map \(T:\mathcal H\to\mathcal K\) has mixed blocks
\[
T_{ai}=Q_aTP_i.
\]
The formula is familiar.  What is less familiar is treating it as the beginning of the theory rather than the last line of a calculation.  Once the resolution is kept, several natural questions become unavoidable:
\begin{enumerate}[label=\arabic*.]
\item What abstract structure is carried by all such block systems?
\item How do two different resolutions of the same object communicate?
\item Which quantities increase, decrease, or remain invariant under refinement?
\item When does a family of operators admit one common exact resolution?
\item How should approximate or compressed resolutions carry certified error?
\item Can local resolutions glue when no global resolution exists?
\end{enumerate}

These questions lead naturally to operator algebra, double categories, refinement lattices, sheaf-like gluing, and a modest logic of certified approximation.  Proof theory fits inside this picture, but it is only one example.  The larger subject is the mathematics of choosing a representation, moving between representations, and keeping track of what those choices reveal or hide.

\begin{figure}[H]
\centering
\includegraphics[width=.91\linewidth]{figures/resolution/01_resolution_architecture.png}
\caption{The resolution is not an annotation placed beside the object.  It participates in the same diagram as the object, the transfer field, and the reconstruction map.}
\end{figure}

\subsection{What is proved here, and what is still a program}

Before going further, it helps to draw the boundary plainly.  The finite projector theory is proved exactly.  The compression and error estimates are proved in the norms named in their statements.  The finite Boolean and quantitative Resolution Logic is developed later in this appendix.  What remains open is the countable, unbounded, and topos-like completion of the theory.  The ordinary unital $C^*$-algebraic representation theorem is proved in a dedicated section below.

\begin{enumerate}[label=\textbf{\Alph*.}]
\item Exact finite-dimensional projector theorems.
\item Exact normed error and compression theorems under their stated hypotheses.
\item Broader categorical and infinite-dimensional extensions that are proposed, not claimed as complete.
\end{enumerate}

\begin{longtable}{L{2.05in}L{0.80in}L{2.95in}}
\toprule
Result & Status & Scope\\
\midrule
\endfirsthead
\toprule
Result & Status & Scope\\
\midrule
\endhead
mixed block reconstruction & proved & finite-dimensional exact \\
block composition and dagger & proved & finite-dimensional exact \\
refinement double category & proved & finite resolutions \\
best-k block approximation & proved & Hilbert-Schmidt norm \\
transfer entropy monotonicity & proved & orthogonal refinement \\
tensor entropy additivity & proved & nonzero finite operators \\
commutant obstruction & proved & exact reducing projectors \\
commutant Laplacian kernel & proved & finite-dimensional \\
resolution-indexed error logic & proved sound & normed linear and bilinear maps \\
finite Resolution Logic & proved sound and finite-complete & Boolean fibers, norm certificates, and finite refinement maps \\
general resolution topos & program & requires a site, sheaf axioms, and infinite completeness \\
countable and unbounded categorical representation & open extension & domain, topology, and convergence control \\
\bottomrule
\end{longtable}

That distinction is not a formality.  The finite theory already says something complete and useful.  The larger categorical program should grow from that solid base rather than being used to blur the difference between a theorem and an ambition.

\subsection{Axioms for a finite resolved transfer system}

The axioms below are not meant to hide the concrete picture behind formal language.  They simply collect the features that keep returning: packets, transfers, reconstruction, refinement, adjoints, tensor products, and a norm compatible with composition.  After stating them once, we return immediately to projector corners and prove that they realize the axioms exactly.

\begin{definition}[Finite resolved transfer system]
A finite resolved transfer system consists of the following data.
\begin{enumerate}[label=\textup{(R\arabic*)}]
\item Each object \(X\) has a finite packet set \(I_X\) and packet spaces \(X_i\).
\item For each pair \(X,Y\), the transfer space decomposes as
\[
\mathfrak T(X,Y)
\cong
\bigoplus_{a\in I_Y,\,i\in I_X}
\mathfrak T_{ai}(X,Y).
\]
\item Reconstruction is faithful: an element of \(\mathfrak T(X,Y)\) is uniquely determined by all of its packet transfers.
\item Composition is matrix composition:
\[
(S\circ T)_{ci}
=
\sum_{a\in I_Y}S_{ca}T_{ai}.
\]
\item Identity transfers are diagonal:
\[
(1_X)_{ji}=\delta_{ji}1_{X_i}.
\]
\item There is a dagger satisfying
\[
(T^\dagger)_{ia}=(T_{ai})^\dagger.
\]
\item Tensor products distribute over packets:
\[
(X\otimes Y)_{(i,a)}=X_i\otimes Y_a.
\]
\item A refinement replaces a packet by finitely many orthogonal descendants, and coarse transfers are sums of the corresponding fine transfers.
\item Transfer spaces are normed, and the chosen norm makes composition continuous.
\end{enumerate}
\end{definition}

These axioms describe the linear skeleton, not the whole body of the theory.  Nonlinear maps enter through projector-resolved jets and interaction hyperedges developed elsewhere in the calculus.  The point here is to make the underlying architecture clear enough that both the linear and nonlinear pieces have somewhere precise to live.

\subsection{Concrete representation by projector corners}

The axioms earn their keep only if the familiar projector construction realizes them without approximation.  In finite dimensions it does, and the proof is nothing more than inserting the two resolutions of the identity in the right places.

\begin{theorem}[Mixed block representation]
Let
\[
\mathcal H=\bigoplus_{i\in I}\mathcal H_i,
\qquad
\mathcal K=\bigoplus_{a\in A}\mathcal K_a,
\]
with orthogonal projectors \(P_i\) and \(Q_a\).  The map
\[
\Phi_{\mathcal Q\leftarrow\mathcal P}:
\mathcal B(\mathcal H,\mathcal K)
\longrightarrow
\bigoplus_{a,i}\mathcal B(\mathcal H_i,\mathcal K_a),
\qquad
T\longmapsto(Q_aTP_i)_{a,i},
\]
is a linear bijection.  Its inverse is
\[
\mathcal R((T_{ai}))=\sum_{a,i}T_{ai}.
\]
For a third resolved space \(\mathcal L=\bigoplus_c\mathcal L_c\),
\[
\Phi_{\mathcal R\leftarrow\mathcal P}(ST)_{ci}
=
\sum_a
\Phi_{\mathcal R\leftarrow\mathcal Q}(S)_{ca}
\Phi_{\mathcal Q\leftarrow\mathcal P}(T)_{ai}.
\]
Moreover,
\[
\Phi(T^*)_{ia}=\bigl(\Phi(T)_{ai}\bigr)^*.
\]
\end{theorem}

\begin{proof}
Insert the two resolutions of the identity on either side of \(T\):
\[
T
=I_\mathcal KTI_\mathcal H
=\left(\sum_aQ_a\right)T\left(\sum_iP_i\right)
=\sum_{a,i}Q_aTP_i.
\]
This is the reconstruction formula.  It also proves faithfulness: if every corner is zero, then their sum is zero and therefore \(T=0\).  Conversely, any family of maps \(T_{ai}:\mathcal H_i\to\mathcal K_a\) determines a unique map between the two direct sums by that same finite sum.

For composition, insert the intermediate resolution exactly where the two maps meet:
\[
R_cSTP_i
=R_cS\left(\sum_aQ_a\right)TP_i
=\sum_a(R_cSQ_a)(Q_aTP_i).
\]
The adjoint law is equally direct:
\[
(Q_aTP_i)^*=P_iT^*Q_a.
\]
Nothing is lost in passing to blocks, and nothing extra has been assumed.
\end{proof}

\begin{corollary}[Linking algebra realization]
For one resolved Hilbert space \(\mathcal H=\bigoplus_i\mathcal H_i\), the block algebra
\[
\bigoplus_{i,j}\mathcal B(\mathcal H_j,\mathcal H_i)
\]
with matrix composition and blockwise adjoint is a unital \(*\)-algebra canonically \(*\)-isomorphic to \(\mathcal B(\mathcal H)\).
\end{corollary}

This corollary explains why the categorical language is natural rather than decorative.  The packet spaces behave as objects, the corners behave as morphisms, and the full operator algebra is recovered as their linking algebra.

\subsection{Hilbert--Schmidt orthogonality and exact accounting}

\begin{theorem}[Packet Parseval identity]
If \(T:\mathcal H\to\mathcal K\) is Hilbert--Schmidt, then
\[
\boxed{
\|T\|_{\mathrm{HS}}^2
=
\sum_{a,i}\|Q_aTP_i\|_{\mathrm{HS}}^2.
}
\]
Distinct packet blocks are orthogonal in the Hilbert--Schmidt inner product.
\end{theorem}

\begin{proof}
For two blocks,
\[
\langle Q_aTP_i,Q_bTP_j\rangle_{\mathrm{HS}}
=
\operatorname{Tr}
\left(P_iT^*Q_aQ_bTP_j\right).
\]
If \(a\ne b\), then \(Q_aQ_b=0\).  If \(i\ne j\), cyclicity of the trace places \(P_jP_i=0\) together.  Thus distinct blocks are orthogonal.  Applying the Pythagorean theorem to the exact finite reconstruction proves the identity.
\end{proof}

This gives a literal meaning to transfer energy.  The squared Hilbert--Schmidt norm is not merely compared with the blocks; it is divided exactly among them.  Every block contributes one orthogonal share, and the shares add back to the whole.

\subsection{Refinement and coarsening}

Now suppose a coarse packet \(P_i\) is opened into finer packets:
\[
P_i=\sum_{\alpha\in C(i)}P_{i\alpha}.
\]
The fine labels remember distinctions that the coarse label suppresses.  Coarsening does not undo the mathematics; it simply forgets those descendant names and sums their transfers back together.

\begin{theorem}[Exact refinement law]
For every operator \(T\),
\[
\boxed{
Q_aTP_i
=
\sum_{\beta\in D(a)}
\sum_{\alpha\in C(i)}
Q_{a\beta}TP_{i\alpha}.
}
\]
If \(S\) and \(T\) are composable, coarsening their fine block product gives the product of their coarse block arrays.
\end{theorem}

\begin{proof}
Insert the two finite refinement identities on the left and right of \(T\).  For composition, begin with the fine block formula, sum over all descendants of the coarse source, intermediate, and target packets, and regroup the finite sums by their coarse parents.  The intermediate descendants sum to the coarse intermediate identity, leaving ordinary coarse block multiplication.
\end{proof}

This is the central coherence law of the finite theory:
\[
\boxed{
\text{refine, compose, and coarsen}
=
\text{compose at the coarse level}.
}
\]
The formula says that refinement adds detail without changing the underlying transformation.

\begin{figure}[H]
\centering
\includegraphics[width=.90\linewidth]{figures/resolution/15_refinement_tree.png}
\caption{A recursive resolution can be read in two ways at once: as a sequence of finer descriptions and as a Bratteli-style diagram recording how packets split.}
\end{figure}

\subsection{The double category of resolved systems}

The preceding theorem reveals two genuinely different kinds of movement.  One applies a mathematical transformation.  The other changes the resolution in which that transformation is being viewed.  Treating both as ordinary arrows would hide the distinction, which is why a double category is the right bookkeeping device.

\begin{theorem}[Resolution double category]
Finite resolved Hilbert spaces form a thin strict double category with:
\begin{itemize}
\item resolved Hilbert spaces as objects;
\item linear maps as horizontal arrows;
\item refinement/coarsening maps on one underlying Hilbert space as vertical arrows;
\item the block-coarsening identity as the unique compatible square.
\end{itemize}
Horizontal composition is operator composition, vertical composition is composition of refinement maps, and the interchange law follows from the exact refinement theorem.
\end{theorem}

\begin{proof}
Horizontal arrows are ordinary linear maps, so their associativity and identities are inherited from operator composition.  Vertical arrows record packet ancestry, and their composition is just repeated refinement or coarsening.  A compatible square says that each coarse block is the sum of its fine descendants.

The only point that needs checking is interchange.  Composing vertically means summing descendants; composing horizontally means multiplying blocks and summing over intermediate packets.  Because all sums are finite and multiplication distributes over addition, either order produces the same result.  That equality is exactly the interchange law.
\end{proof}

\begin{figure}[H]
\centering
\includegraphics[width=.82\linewidth]{figures/resolution/11_resolution_double_category.png}
\caption{The square keeps one transformation fixed while the vertical arrows change how its source and target are resolved.}
\end{figure}

The double-category language earns its place here because it keeps two ideas from being confused.  Refining a description is not another physical evolution, and changing coordinates is not another step of the dynamics.  The diagram remembers which kind of movement has occurred.

\subsection{Dagger and monoidal structure}

\begin{proposition}[Tensor compatibility]
For resolved systems \((\mathcal H,\{P_i\})\) and \((\mathcal K,\{Q_a\})\), the product projectors
\[
P_i\otimes Q_a
\]
form a resolution of \(\mathcal H\otimes\mathcal K\).  For operators \(S\) and \(T\),
\[
\bigl(P_i\otimes Q_a\bigr)(S\otimes T)
\bigl(P_j\otimes Q_b\bigr)
=
(P_iSP_j)\otimes(Q_aTQ_b).
\]
\end{proposition}

\begin{proof}
Orthogonality and completeness follow from the corresponding identities in each factor.  The block identity follows from the elementary tensor-product multiplication law.
\end{proof}

So the finite theory is naturally dagger symmetric monoidal.  Two resolved systems can be combined without abandoning the packet description: their packets simply become ordered pairs, and their transfers factor exactly when the operator does.

\begin{theorem}[Additivity of transfer entropy under tensor products]
Let \(S\) and \(T\) be nonzero Hilbert--Schmidt operators.  Normalize their packet energies to probability distributions
\[
p_{ij}^S
=
\frac{\|P_iSP_j\|_{\mathrm{HS}}^2}{\|S\|_{\mathrm{HS}}^2},
\qquad
p_{ab}^T
=
\frac{\|Q_aTQ_b\|_{\mathrm{HS}}^2}{\|T\|_{\mathrm{HS}}^2}.
\]
Then the product resolution satisfies
\[
\boxed{
H_{\mathcal P\otimes\mathcal Q}(S\otimes T)
=
H_\mathcal P(S)+H_\mathcal Q(T).
}
\]
\end{theorem}

\begin{proof}
The Hilbert--Schmidt norm is multiplicative under tensor products, and each product-block energy is the product of the two factor energies.  The normalized distribution is therefore \(p^S\otimes p^T\).  Shannon entropy is additive on product distributions.
\end{proof}

The numerical example gives
\[
H(S)=1.875561,
\qquad
H(T)=1.166399,
\qquad
H(S\otimes T)=3.041960.
\]

\begin{figure}[H]
\centering
\includegraphics[width=.76\linewidth]{figures/resolution/12_tensor_transfer_factorization.png}
\caption{Normalized transfer energy factorizes in a tensor-product resolution, which makes transfer entropy additive.}
\end{figure}

\subsection{Compression as orthogonal projection}

Compression is often described as deletion: keep the pieces that look important and throw away the rest.  In a fixed orthogonal resolution, the situation is cleaner.  The block-diagonal part is an orthogonal projection in operator space, and every omitted block remains visible as an exact residual.

For one resolution \(\mathcal P\), define the block-diagonal map
\[
\mathbb E_\mathcal P(T)
=
\sum_iP_iTP_i.
\]

\begin{theorem}[Conditional expectation and best diagonal approximation]
The map \(\mathbb E_\mathcal P\) is unital, completely positive, trace preserving, adjoint preserving, and idempotent.  In the Hilbert--Schmidt inner product it is the orthogonal projection onto the block-diagonal subalgebra
\[
\mathcal D_\mathcal P
=
\bigoplus_iP_i\mathcal B(\mathcal H)P_i.
\]
Consequently,
\[
\|T-\mathbb E_\mathcal P(T)\|_{\mathrm{HS}}
\le
\|T-D\|_{\mathrm{HS}}
\]
for every \(D\in\mathcal D_\mathcal P\).
\end{theorem}

\begin{proof}
The projectors \(P_i\) are Kraus operators, so the map is completely positive.  Completeness gives unitality and trace preservation.  Idempotence follows from orthogonality.  The packet Parseval identity shows that diagonal and off-diagonal block families are mutually orthogonal Hilbert--Schmidt subspaces.  Hence \(\mathbb E_\mathcal P\) is the orthogonal projection onto the diagonal subspace, and the best-approximation inequality follows.
\end{proof}

More generally, for an edge set \(E\subseteq A\times I\), define
\[
C_E(T)=\sum_{(a,i)\in E}Q_aTP_i.
\]
The selected blocks span an orthogonal subspace, and \(C_E\) is the corresponding projection.  What is left out is therefore not a vague notion of lost detail.  It is an explicit orthogonal residual whose norm can be computed exactly.

\begin{theorem}[Exact selective-compression error]
\[
\boxed{
\|T-C_E(T)\|_{\mathrm{HS}}^2
=
\sum_{(a,i)\notin E}
\|Q_aTP_i\|_{\mathrm{HS}}^2.
}
\]
Among all approximations using at most \(k\) blocks of one fixed resolution, the optimal Hilbert--Schmidt approximation retains the \(k\) blocks with largest squared norms.
\end{theorem}

\begin{proof}
The first statement is Pythagoras.  For the second, every admissible \(k\)-block approximation omits all but \(k\) entries from one nonnegative list of block energies.  The omitted sum is minimized by removing the \(k\) largest entries from that sum, equivalently by retaining them.
\end{proof}

\begin{figure}[H]
\centering
\includegraphics[width=.79\linewidth]{figures/resolution/05_best_k_error_curve.png}
\caption{The best-\(k\) curve is exact for a fixed resolution because the block residuals are orthogonal.}
\end{figure}

\subsection{Resolution complexity}

The best-\(k\) theorem immediately suggests a practical question: how many transfers must be kept before the approximation reaches a chosen tolerance?  The answer leads to a finite-resolution complexity function.

\begin{definition}[Resolution complexity]
For a tolerance \(\varepsilon\ge0\), define
\[
d_{\mathcal Q\leftarrow\mathcal P}(T;\varepsilon)
=
\min\left\{|E|:
\|T-C_E(T)\|_{\mathrm{HS}}\le\varepsilon
\right\}.
\]
\end{definition}

\begin{corollary}[Closed formula]
If the block energies are sorted as
\[
e_1\ge e_2\ge\cdots\ge e_m\ge0,
\]
then
\[
\boxed{
d(T;\varepsilon)
=
\min\left\{k:\sum_{r>k}e_r\le\varepsilon^2
\right\}.
}
\]
\end{corollary}

This quantity belongs to the resolved pair \((T,\mathcal Q\leftarrow\mathcal P)\), not to \(T\) in isolation.  A different resolution may make the same operator look far simpler or far more complicated.  To obtain a representation-independent quantity, one must first say which resolutions are admissible and then minimize over that class.

\begin{figure}[H]
\centering
\includegraphics[width=.79\linewidth]{figures/resolution/06_resolution_complexity.png}
\caption{Resolution complexity answers a concrete question: how many transfer blocks must be retained to meet a chosen error tolerance?}
\end{figure}

\subsection{Transfer entropy and what refinement reveals}

Once block energies have been normalized, they form a probability distribution.  Its entropy measures how broadly the operator's Hilbert--Schmidt energy is spread across the chosen transfer field.  For nonzero Hilbert--Schmidt \(T\), define
\[
p_{ai}
=
\frac{\|Q_aTP_i\|_{\mathrm{HS}}^2}{\|T\|_{\mathrm{HS}}^2}
\]
and
\[
H_{\mathcal Q\leftarrow\mathcal P}(T)
=-\sum_{a,i}p_{ai}\log p_{ai}.
\]

\begin{theorem}[Entropy refinement law]
Let fine resolutions refine coarse resolutions.  Then
\[
\boxed{
H_{\mathrm{fine}}(T)
=
H_{\mathrm{coarse}}(T)
+
\sum_{a,i}p_{ai}
H(\text{fine transfers inside coarse block }(a,i)).
}
\]
In particular,
\[
H_{\mathrm{fine}}(T)
\ge
H_{\mathrm{coarse}}(T).
\]
\end{theorem}

\begin{proof}
The coarse block is the orthogonal sum of its fine descendants, so its squared norm is the sum of the descendant squared norms.  The fine probability distribution is therefore a refinement of the coarse distribution.  The ordinary chain rule for Shannon entropy yields the identity.
\end{proof}

Refinement cannot lower this entropy when the coarse blocks are merely split into orthogonal descendants.  It may show that a coarse transfer was almost entirely concentrated in one descendant, or it may reveal that the same coarse transfer was hiding several competing channels.

\begin{figure}[H]
\centering
\includegraphics[width=.79\linewidth]{figures/resolution/04_entropy_under_refinement.png}
\caption{Splitting a packet exposes distinctions that were previously hidden inside one coarse transfer, so the resolved transfer entropy increases.}
\end{figure}

There is no contradiction with the next example.  Refining one fixed resolution can only reveal more internal detail, but replacing a poor resolution with a better adapted one may sharply reduce the apparent transfer spread.  For the same operator,
\[
\mathcal E_{\mathrm{off}}^{\mathrm{adapted}}
=0.000e+00,
\qquad
\mathcal E_{\mathrm{off}}^{\mathrm{unadapted}}
=21.378173.
\]
Its transfer entropy changes from
\[
2.141799
\quad\text{to}\quad
1.089108.
\]

\begin{figure}[H]
\centering
\includegraphics[width=.70\linewidth]{figures/resolution/02_adapted_transfer_energy.png}
\caption{The adapted resolution exposes the exact reducing sectors of a hidden block operator.}
\end{figure}

\begin{figure}[H]
\centering
\includegraphics[width=.70\linewidth]{figures/resolution/03_unadapted_transfer_energy.png}
\caption{An unrelated resolution spreads the same operator across many transfer blocks.}
\end{figure}

\subsection{Why a resolution is not a free speedup}

A useful resolution can make structure obvious, but the act of drawing block boundaries does not by itself save memory or arithmetic.

\begin{proposition}[Storage conservation under complete blockization]
Let \(\dim\mathcal H=n\) and \(\dim\mathcal K=m\).  If every mixed block is stored explicitly, then
\[
\sum_{a,i}
\dim\mathcal K_a\,\dim\mathcal H_i
=
mn.
\]
Thus complete blockization stores exactly as many scalar entries as the original dense map.
\end{proposition}

\begin{proof}
Factor the two finite sums:
\[
\sum_{a,i}\dim\mathcal K_a\dim\mathcal H_i
=
\left(\sum_a\dim\mathcal K_a\right)
\left(\sum_i\dim\mathcal H_i\right)
=mn.
\]
\end{proof}

The same warning applies to generic dense composition.  If every block triple is active and each one is multiplied densely, the scalar work has merely been rearranged.  Resolution becomes computationally valuable only when it exposes something real: zero transfers, invariant sectors, repeated blocks, low rank, local coupling, transform diagonalization, selective outputs, or enough reuse to amortize the setup cost.

\subsection{A cost-aware transfer planner}

Suppose block \((a,i)\) has modeled execution cost \(c_{ai}\) and omission penalty proportional to its energy \(e_{ai}\).  For an edge set \(E\), consider
\[
J_\lambda(E)
=
\sum_{(a,i)\in E}c_{ai}
+
\lambda
\sum_{(a,i)\notin E}e_{ai}.
\]

\begin{theorem}[Optimal separable planner]
The minimizer retains exactly those edges satisfying
\[
\boxed{c_{ai}<\lambda e_{ai}.}
\]
Ties may be resolved arbitrarily.
\end{theorem}

\begin{proof}
The objective is a sum of independent two-choice terms.  For one edge, retaining it contributes \(c_{ai}\); omitting it contributes \(\lambda e_{ai}\).  Choose the smaller term independently for every edge.
\end{proof}

The theorem is elementary, and that is part of its value: the decision boundary is visible.  Real planners are harder because packing, launch, memory, and dependency costs couple the choices.  The simple result is therefore a baseline against which more realistic planners can be judged.

\begin{figure}[H]
\centering
\includegraphics[width=.79\linewidth]{figures/resolution/07_planner_frontier.png}
\caption{Changing the value placed on error moves the planner along a visible cost--accuracy frontier.}
\end{figure}

\subsection{Commuting resolutions and exact common refinement}

Two resolutions may describe the same space faithfully and still refuse to fit inside one common orthogonal refinement.

\begin{theorem}[Common-refinement criterion]
Two finite orthogonal resolutions \(\{P_i\}\) and \(\{Q_a\}\) on the same Hilbert space admit a common orthogonal refinement if and only if
\[
[P_i,Q_a]=0
\qquad\text{for every }i,a.
\]
When they commute, the nonzero products \(P_iQ_a\) form a common refinement.
\end{theorem}

\begin{proof}
If the resolutions commute, each product \(P_iQ_a\) is a self-adjoint idempotent.  Products with different index pairs are orthogonal, and
\[
\sum_{i,a}P_iQ_a=I.
\]
Summing over \(a\) recovers \(P_i\), while summing over \(i\) recovers \(Q_a\).

Conversely, suppose \(\{R_\gamma\}\) refines both resolutions.  Each \(P_i\) and \(Q_a\) is a sum of members of one mutually orthogonal family.  Sums of members of one orthogonal family commute.  Hence every \(P_i\) commutes with every \(Q_a\).
\end{proof}

This theorem gives the boundary exactly.  Commuting resolutions can be merged into one finer chart.  Noncommuting resolutions must be related as an atlas of valid charts rather than forced into a false common one.

\subsection{Projector atlases for noncommuting resolutions}

When two resolutions do not commute, the products \(P_iQ_a\) are generally not projectors.  Pretending otherwise does not simplify the problem; it changes it.  The honest replacement is an atlas of valid resolutions connected by exact mixed corners.

Let \(\mathcal P=\{P_i\}\) and \(\mathcal Q=\{Q_a\}\).  For any map \(T\), the atlas transition data are
\[
Q_aTP_i.
\]
No commutation assumption is needed, and composition remains exact.  The atlas does not manufacture a global chart.  It keeps the two legitimate local descriptions and records the transition data between them.

\begin{figure}[H]
\centering
\includegraphics[width=.82\linewidth]{figures/resolution/09_noncommuting_projector_atlas.png}
\caption{The Pauli \(Z\) and \(X\) eigenprojectors have no nontrivial common refinement.  Their mixed corners remain perfectly well defined.}
\end{figure}

\subsection{Exact common structure and the commutant obstruction}

Let \(\mathcal F=\{A_1,\ldots,A_r\}\) be a finite operator family.

\begin{theorem}[Common reducing projector criterion]
The family \(\mathcal F\) admits a nontrivial exact two-packet reducing resolution if and only if its commutant
\[
\mathcal F'
=
\{X:[X,A_k]=0\text{ for every }k\}
\]
contains a projection other than \(0\) and \(I\).
\end{theorem}

\begin{proof}
If \(P\) is a reducing projector, then \([P,A_k]=0\) for every \(k\), so \(P\in\mathcal F'\).  Conversely, a nontrivial projection in the commutant gives the exact resolution \(\{P,I-P\}\), and each operator is block diagonal in that resolution.
\end{proof}

Irreducibility is therefore not a failure of the method.  It is information.  It says that the operator family has no nontrivial exact common packetization.

\subsection{The commutant Laplacian}

Even when the exact commutant is trivial, one may still look for nearly shared structure---provided the resulting packets are treated as approximate and their leakage is measured.  For self-adjoint \(A_k\), define
\[
\mathcal L_\mathcal F(X)
=
\sum_k[A_k,[A_k,X]].
\]

\begin{theorem}[Commutant-Laplacian identity]
With the Hilbert--Schmidt inner product,
\[
\boxed{
\langle X,\mathcal L_\mathcal F(X)\rangle_{\mathrm{HS}}
=
\sum_k\|[A_k,X]\|_{\mathrm{HS}}^2.
}
\]
Consequently, \(\mathcal L_\mathcal F\) is positive semidefinite and
\[
\ker\mathcal L_\mathcal F=\mathcal F'.
\]
\end{theorem}

\begin{proof}
For one self-adjoint \(A\), the commutator map \(\delta_A(X)=[A,X]\) is self-adjoint as an operator on Hilbert--Schmidt space up to the standard sign convention, and direct expansion gives
\[
\langle X,[A,[A,X]]\rangle
=
\|[A,X]\|_{\mathrm{HS}}^2.
\]
Sum over \(k\).  The quadratic form vanishes exactly when every commutator vanishes.
\end{proof}

Small nonzero eigenvalues of \(\mathcal L_\mathcal F\) point toward matrices that almost commute with the family.  Their spectral projectors are natural candidates for approximate packets.  They are candidates only: the commutator and leakage estimates still have to be carried into any later theorem.

\begin{figure}[H]
\centering
\includegraphics[width=.79\linewidth]{figures/resolution/08_commutant_laplacian_spectrum.png}
\caption{Exact common symmetries appear at zero.  A perturbation lifts them into a small positive cluster, which reveals approximate rather than exact structure.}
\end{figure}

\subsection{Refinement order and the partition lattice}

If all resolutions are required to refine one fixed atomic commuting resolution, they are ordinary set partitions of the atom labels.  Refinement becomes the standard partition order.

\begin{theorem}[Finite lattice of coordinate resolutions]
Partitions of a finite atomic packet set form a lattice.  The meet is the common refinement formed by all nonempty intersections.  The join is the coarsest partition containing both, obtained from the equivalence relation generated by belonging to a common block in either partition.
\end{theorem}

\begin{proof}
Every nonempty intersection lies inside a block of each original partition and is therefore a common refinement.  Any other common refinement must refine every such intersection, making it the greatest lower bound.  For the join, the transitive closure of the two block relations is an equivalence relation containing both.  Every common coarsening must contain that generated relation, so it is the least upper bound.
\end{proof}

\begin{figure}[H]
\centering
\includegraphics[width=.84\linewidth]{figures/resolution/10_partition_lattice.png}
\caption{Commuting coordinate resolutions inherit the familiar partition lattice.  Noncommuting projector systems generally do not.}
\end{figure}

This finite lattice is a natural starting point for a theory of infinite refinement.  Directed chains can be understood through their finite stages.  Once noncommuting or genuinely infinite resolutions are allowed, however, ordinary partition order is no longer enough; one needs an atlas or a richer domain structure.

\subsection{Local resolutions and gluing}

A single global projector family may fail to exist even when every small region carries a perfectly good local one.  This is a familiar geometric situation: the local descriptions are not wrong; they are twisted together.  The natural object is again an atlas.

\begin{theorem}[Finite unitary gluing]
Let \(\{U_a\}\) cover a parameter space.  Suppose each \(U_a\) carries a local packet model, and unitary transition maps \(g_{ab}\) on overlaps satisfy
\[
g_{aa}=I,
\qquad
g_{ab}=g_{ba}^{-1},
\qquad
g_{ab}g_{bc}=g_{ac}
\]
on triple overlaps.  Then the local models glue to a global Hermitian vector bundle.  Compatible local sections glue uniquely to global sections.
\end{theorem}

\begin{proof}
Take the disjoint union of the local trivial bundles and identify \((x,v)_b\) with \((x,g_{ab}(x)v)_a\) on overlaps.  The cocycle law makes the equivalence relation transitive.  Unitarity preserves the Hermitian structure.  Local sections that satisfy the transition law define one equivalence class at every point and therefore one global section.  Uniqueness follows because a global section is fixed by its local representatives.
\end{proof}

\begin{figure}[H]
\centering
\includegraphics[width=.84\linewidth]{figures/resolution/14_local_to_global_gluing.png}
\caption{Local packet charts glue into a global bundle when their overlap maps satisfy the usual cocycle law.}
\end{figure}

The theorem changes the meaning of failure.  The absence of one global packetization need not mean that the object has no resolved structure.  It may mean that the structure is genuinely global and can be seen only through compatible local charts.

\subsection{The normed seed of Resolution Logic}

The complete finite Resolution Logic developed later begins with the two rules below.  They are the normed seed of the larger system and remain useful on their own.  It begins with an ordinary promise: an approximation is within a stated distance of the object it represents.

Write
\[
\widehat x\Vdash x[\varepsilon]
\]
to mean
\[
\|x-\widehat x\|\le\varepsilon.
\]
The following rules are sound.

\begin{theorem}[Linear and bilinear error transport]
If \(\widehat x\Vdash x[\varepsilon]\) and \(A\) is bounded, then
\[
A\widehat x\Vdash Ax[\|A\|\varepsilon].
\]
If \(B:X\times Y\to Z\) is a bounded bilinear map with norm \(\|B\|\), and
\[
\widehat x\Vdash x[\varepsilon],
\qquad
\widehat y\Vdash y[\eta],
\]
then
\[
B(\widehat x,\widehat y)
\Vdash
B(x,y)
\left[
\|B\|
\bigl(
\|\widehat x\|\eta
+
\|\widehat y\|\varepsilon
+
\varepsilon\eta
\bigr)
\right].
\]
\end{theorem}

\begin{proof}
The linear rule is the definition of the operator norm.  For the bilinear rule, write
\[
x=\widehat x+e,
\qquad
y=\widehat y+f.
\]
Then
\[
B(x,y)-B(\widehat x,\widehat y)
=
B(e,\widehat y)
+B(\widehat x,f)
+B(e,f).
\]
Apply the bilinear norm bound and \(\|e\|\le\varepsilon\), \(\|f\|\le\eta\).
\end{proof}

The judgment becomes resolution-indexed when \(\widehat x\) is produced by a named packet compression and \(\varepsilon\) is the residual certified by that compression.  Refinement may lower the residual; coarsening may enlarge it.  A complete statement therefore records both the approximation and the resolution in which the error claim is valid.

\begin{figure}[H]
\centering
\includegraphics[width=.79\linewidth]{figures/resolution/13_resolution_error_logic.png}
\caption{Across the numerical sweep, the certified bound stays above the observed error, as the theorem requires.}
\end{figure}

\subsection{Compression, proof objects, and representation}

Proof compression is one application of the theory, but Resolution Theory is not merely proof theory under another name.  A proof certificate can be viewed as one resolved representation of a larger discovery object.  The steps, examples, failed routes, and neighboring lemmas that were omitted form a residual.  That residual may be irrelevant for verification and indispensable for a later theorem.

The operator-theoretic results suggest several disciplined questions:
\begin{enumerate}[label=\arabic*.]
\item What is the admissible class of proof resolutions?
\item Which morphisms survive compression?
\item Is the compressed certificate sufficient for verification only, or also for reuse?
\item What is the resolution complexity of reconstructing a bridge-critical argument?
\item Do two proof representations admit a common refinement, or must they be kept as an atlas?
\end{enumerate}

The practical conclusion is not that every proof should be expanded forever.  It is that compression should have a declared purpose, a visible residual, and some account of how the omitted material can be recovered when the purpose changes.

\subsection{Worked example: a hidden reducing decomposition}

Begin with a Hermitian operator that is block diagonal in a basis we do not yet know.  In the adapted resolution, every off-diagonal transfer vanishes.  In a coordinate resolution of the same block ranks, the blocks are generally dense.  The example establishes four separate facts:

\begin{enumerate}[label=\arabic*.]
\item Block structure is representation dependent.
\item The existence of an adapted resolution is an invariant statement: it is equivalent to nontrivial projections in the commutant.
\item Once the resolution is known, exact compression to diagonal blocks loses nothing.
\item Discovering the resolution may be harder than applying it.
\end{enumerate}

The reference experiment recovered zero adapted off-diagonal energy to floating-point precision while the unrelated coordinate resolution had off-diagonal energy
\[
21.378173.
\]

\subsection{Worked example: incompatible quantum charts}

Let
\[
P_\pm^Z=\frac12(I\pm\sigma_z),
\qquad
P_\pm^X=\frac12(I\pm\sigma_x).
\]
The two resolutions do not commute.  By the common-refinement theorem they have no common nontrivial orthogonal refinement.  Nevertheless,
\[
P_a^XTP_i^Z
\]
is well defined for every \(T\), and mixed block composition remains exact.

The example is only two-dimensional, but it contains the entire atlas idea.  Two incompatible descriptions do not have to be forced into one artificial coordinate system.  They can remain distinct and still be linked exactly.

\subsection{Worked example: tensor-resolved systems}

Suppose one system is resolved by energy scale and another by symmetry charge.  Their product resolution labels packets by pairs
\[
(\text{scale},\text{charge}).
\]
If the operator factors, its transfer distribution factors and the transfer entropy is additive.  If an interaction is introduced, new cross-packet structure appears.  The difference between the interacting entropy and the sum of factor entropies becomes one measure of how the chosen interaction spreads transfer across the product resolution.

\subsection{What the code actually checks}

The Python implementation is intentionally plain.  It uses dense projectors and favors readability over speed, so it should be read as a reference implementation and an audit trail, not as a production kernel.  Whenever a theorem has a finite numerical form, the tests check that form directly.

The automated suite verifies:
\begin{itemize}
\item exact mixed reconstruction;
\item block composition and dagger laws;
\item tensor resolution;
\item refinement/coarsening coherence;
\item packet Parseval identity;
\item optimal block-diagonal projection;
\item entropy monotonicity and tensor additivity;
\item best-\(k\) optimality by brute force;
\item resolution-complexity tails;
\item the separable planner rule;
\item commuting and noncommuting resolution behavior;
\item commutant-Laplacian recovery;
\item partition meet and join;
\item certified linear and bilinear error transport;
\item packetwise-unitary invariance of transfer entropy;
\item unitary cocycle consistency;
\item permutation-invariant resolution distance;
\item transfer mutual information;
\item budget--tolerance generalized inverse laws;
\item commutant-gap distance certificates;
\item approximate cocycle and path-disagreement bounds;
\item directed refinement consistency;
\item effective resolution-dimension recovery.
\end{itemize}

The complete run passed
\[
\boxed{28/28\text{ tests}}.
\]

\begin{center}
\small
\begin{tabular}{rrrr}
\toprule
Dimension & Packets & Median split/reconstruct time (ms) & Reconstruction error\\
\midrule
24 & 4 & 0.1717 & 0.00e+00 \\
48 & 6 & 2.1411 & 0.00e+00 \\
72 & 8 & 7.9590 & 0.00e+00 \\
96 & 8 & 12.8127 & 0.00e+00 \\
\bottomrule
\end{tabular}
\end{center}

\begin{figure}[H]
\centering
\includegraphics[width=.78\linewidth]{figures/resolution/16_reference_runtime.png}
\caption{The reference code is deliberately literal.  A production implementation would replace dense projectors with index maps, sparse blocks, or native packet storage.}
\end{figure}

\subsection{A metric geometry on the space of resolutions}

Up to this point a resolution has been treated mainly as an algebraic choice.  There is also a useful geometric question: how far apart are two ways of resolving the same space?  That question matters whenever a resolution moves continuously, is learned from data, or changes under perturbation.  A theory that can only say whether two resolutions are equal is too coarse for those situations.

For orthogonal projectors $P$ and $Q$ on the same finite-dimensional Hilbert space, define
\[
 d_{\mathrm{ch}}(P,Q)
 =\frac{1}{\sqrt2}\|P-Q\|_{\mathrm{HS}}.
\]
When $P$ and $Q$ have the same rank, this is the Euclidean norm of the sines of their principal angles.  It is zero exactly when the two subspaces agree.

\begin{proposition}[Chordal projector metric]
The function $d_{\mathrm{ch}}$ is a metric on the set of orthogonal projectors on a fixed finite-dimensional Hilbert space.
\end{proposition}

\begin{proof}
Symmetry and nonnegativity are immediate.  If $d_{\mathrm{ch}}(P,Q)=0$, then $P=Q$.  The triangle inequality is inherited from the Hilbert--Schmidt norm:
\[
 d_{\mathrm{ch}}(P,R)
 \le d_{\mathrm{ch}}(P,Q)+d_{\mathrm{ch}}(Q,R).
\]
Nothing more is needed.  The normalization by $\sqrt2$ is conventional and makes the distance agree with the principal-angle formula for equal ranks.
\end{proof}

A resolution is an unordered family of projectors, so packet labels should not contribute to distance.  For two resolutions with the same number of packets, define
\[
 d_{\mathrm{res}}(\mathcal P,\mathcal Q)^2
 =
 \min_{\sigma}
 \sum_i d_{\mathrm{ch}}(P_i,Q_{\sigma(i)})^2,
\]
where the minimum runs over all permutations of the packet labels.

\begin{theorem}[Permutation-invariant resolution metric]
On the space of finite resolutions with a fixed packet count, modulo relabeling, $d_{\mathrm{res}}$ is a metric.
\end{theorem}

\begin{proof}
Nonnegativity and symmetry follow from the projector metric.  The value vanishes precisely when the two families agree after a permutation.  For the triangle inequality, choose optimal matchings from $\mathcal P$ to $\mathcal Q$ and from $\mathcal Q$ to $\mathcal R$.  Their composition gives an admissible matching from $\mathcal P$ to $\mathcal R$.  Applying Minkowski's inequality to the corresponding vector of projector distances gives
\[
 d_{\mathrm{res}}(\mathcal P,\mathcal R)
 \le
 d_{\mathrm{res}}(\mathcal P,\mathcal Q)
 +
 d_{\mathrm{res}}(\mathcal Q,\mathcal R).
\]
The minimum over all direct matchings can only be smaller.
\end{proof}

This metric turns a moving packet family into an ordinary path in a metric space.  Its length measures how much the representation itself has moved, independently of the physical operator being represented.

\begin{figure}[H]
\centering
\includegraphics[width=.78\linewidth]{figures/resolution/17_resolution_metric_path.png}
\caption{A two-packet resolution rotated through a continuous family.  The distance is insensitive to packet relabeling but records genuine motion of the underlying subspaces.}
\end{figure}

The distinction between physical transfer and representational motion now has a geometric expression.  If $T(t)$ and $\mathcal P(t)$ both vary, one may separately measure
\[
 \|\dot T(t)\|
 \qquad\text{and}\qquad
 \operatorname{Length}(\mathcal P|_{[t,t+\Delta t]}).
\]
A large change in the block matrix may come from either source.  The metric prevents those effects from being silently conflated.

\subsection{Budget and tolerance as generalized inverses}

Compression is often described informally as choosing a budget and keeping the most important pieces.  There are actually two dual questions:
\begin{enumerate}[label=\arabic*.]
\item Given a budget of $k$ blocks, what is the smallest possible residual?
\item Given an allowed residual $\varepsilon$, how many blocks are necessary?
\end{enumerate}

Let the block energies be arranged in decreasing order,
\[
 e_1\ge e_2\ge\cdots\ge e_M\ge0.
\]
The best-$k$ theorem gives
\[
 E(k)
 =
 \left(\sum_{j>k}e_j\right)^{1/2}.
\]
Define the resolution complexity
\[
 C(\varepsilon)
 =
 \min\{k:E(k)\le\varepsilon\}.
\]

\begin{theorem}[Budget--tolerance duality]
For every $k$ and every $\varepsilon\ge0$,
\[
 \boxed{
 C(\varepsilon)\le k
 \quad\Longleftrightarrow\quad
 E(k)\le\varepsilon.
 }
\]
Thus $E$ and $C$ are monotone generalized inverses.
\end{theorem}

\begin{proof}
If $C(\varepsilon)\le k$, then some budget no larger than $k$ achieves residual at most $\varepsilon$.  Since $E(k)$ is the optimal residual among all $k$-block approximations and is nonincreasing in $k$, it follows that $E(k)\le\varepsilon$.  Conversely, if $E(k)\le\varepsilon$, then $k$ is an admissible value in the defining minimum for $C(\varepsilon)$, so $C(\varepsilon)\le k$.
\end{proof}

The theorem is elementary, but it is the exact logical form behind a practical planner.  A tolerance query and a resource query are not separate optimization problems; they are two views of the same sorted energy profile.

\begin{figure}[H]
\centering
\includegraphics[width=.94\linewidth]{figures/resolution/18_compression_galois_curve.png}
\caption{The best-$k$ residual and the minimum budget required for a prescribed tolerance.  Each curve determines the other.}
\end{figure}

This duality also clarifies what a compression certificate should contain.  A useful certificate does not merely list the retained blocks.  It records the entire residual curve, or at least enough of it to answer later budget and tolerance questions without recomputing the decomposition.

\subsection{Directed refinement and finite approximants}

Recursive refinement suggests an infinite limit, but an infinite diagram is useful only when finite stages approximate it in a declared topology.  The Hilbert--Schmidt class gives a clean setting because packet blocks are orthogonal there.

Let
\[
 \mathcal P_1\preceq\mathcal P_2\preceq\cdots
\]
be an increasing sequence of finite resolutions.  Write $\mathcal A_n$ for the finite algebra generated by the projectors at level $n$, and let $\mathbb E_n$ be the Hilbert--Schmidt orthogonal projection onto the transfer blocks retained at that level.

\begin{theorem}[Directed Hilbert--Schmidt convergence]
Suppose the union of the packet subspaces generated by the refinement chain is dense and $T$ is Hilbert--Schmidt.  Then the exact finite packet reconstructions converge to $T$ in Hilbert--Schmidt norm.  More generally, if the retained edge sets increase to all packet pairs, then
\[
 \|T-\mathbb E_nT\|_{\mathrm{HS}}\longrightarrow0.
\]
\end{theorem}

\begin{proof}
Hilbert--Schmidt orthogonality gives
\[
 \|T\|_{\mathrm{HS}}^2
 =
 \sum_{i,j}\|P_iTP_j\|_{\mathrm{HS}}^2
\]
for the limiting countable resolution.  Every finite stage retains a finite subset of this nonnegative summable family.  The squared residual is the complementary tail.  Monotone convergence of the partial sums therefore forces the tail to zero.
\end{proof}

This theorem does not say that every finite stage is computationally useful.  It says something more basic: the finite stages really are approximants to one object, not unrelated truncations sharing notation.

\begin{figure}[H]
\centering
\includegraphics[width=.78\linewidth]{figures/resolution/19_directed_refinement_convergence.png}
\caption{A localized Hilbert--Schmidt transfer field approximated by increasingly fine local packet interactions.  Exact reconstruction is available at every level; the plotted curve concerns a deliberately local retained-edge model.}
\end{figure}

The same idea points toward a domain-theoretic formulation.  Finite resolutions behave like compact approximants, while the limiting resolved object is their directed supremum.  A complete theory would identify the class of transfer functors that preserve such suprema.  The finite theorem above is the first continuity statement that such a theory must recover.

\subsection{Descent, cocycles, and approximate gluing}

A projector atlas becomes genuinely geometric when local representations can be transported and compared on overlaps.  Let $g_{ab}$ transport coordinates from chart $b$ to chart $a$.  Exact gluing requires
\[
 g_{ab}g_{bc}=g_{ac}
\]
on triple overlaps.

When the equality is approximate, two paths between the same charts may disagree.  That disagreement is not a nuisance to be hidden; it is the measurable obstruction to global reconstruction.

\begin{proposition}[Path accumulation bound]
Assume every transition is unitary and every elementary cocycle defect along a chosen triangulation is at most $\delta$ in operator norm.  If two homotopic chart paths between the same endpoints differ by at most $N$ elementary triangle moves, then their transported vectors differ by at most
\[
 N\delta\|v\|.
\]
\end{proposition}

\begin{proof}
One elementary triangle move replaces $g_{ab}g_{bc}$ by $g_{ac}$.  Its effect on a vector is bounded by $\delta\|v\|$.  All other transitions in the path are unitary and therefore do not amplify the norm.  Telescoping over at most $N$ moves gives the result.
\end{proof}

The bound is intentionally linear and conservative.  It is enough to make a conceptual point: approximate descent comes with an audit trail.  One can state how path dependent the reconstructed object may be, rather than merely saying that the charts are ``almost compatible.''

\begin{figure}[H]
\centering
\includegraphics[width=.78\linewidth]{figures/resolution/20_approximate_descent.png}
\caption{A controlled cocycle perturbation produces controlled path disagreement.  Local-to-global failure grows continuously with the defect.}
\end{figure}

In a more developed sheaf or stack formulation, the cocycle defect would be replaced by a higher coherence object.  The finite estimate already shows what that object must control: not only local consistency, but the dependence of reconstruction on the path used to compare local models.

\subsection{Approximate common structure and the commutant gap}

The exact common-commutant theorem answers a yes-or-no question.  Real systems often require a quantitative version: how far is a candidate structure from the exact joint commutant?

Let
\[
 \mathcal L(X)
 =
 \sum_r[A_r^*,[A_r,X]]
\]
be the commutant Laplacian.  Its kernel is the joint commutant.  Let $\lambda_+$ be the smallest positive eigenvalue of $\mathcal L$ on the Hilbert--Schmidt space.

\begin{theorem}[Commutant-gap distance certificate]
Let $\Pi$ be the Hilbert--Schmidt orthogonal projection onto $\ker\mathcal L$.  Then
\[
 \boxed{
 \|X-\Pi X\|_{\mathrm{HS}}^2
 \le
 \frac{\langle X,\mathcal LX\rangle_{\mathrm{HS}}}{\lambda_+}.
 }
\]
\end{theorem}

\begin{proof}
Expand $X$ in an orthonormal eigenbasis of the positive semidefinite operator $\mathcal L$:
\[
 X=X_0+\sum_{j:\lambda_j>0}c_jE_j,
 \qquad
 X_0\in\ker\mathcal L.
\]
Then
\[
 \|X-\Pi X\|_{\mathrm{HS}}^2
 =\sum_{\lambda_j>0}|c_j|^2,
\]
while
\[
 \langle X,\mathcal LX\rangle
 =\sum_{\lambda_j>0}\lambda_j|c_j|^2
 \ge
 \lambda_+\sum_{\lambda_j>0}|c_j|^2.
\]
Rearranging proves the claim.
\end{proof}

The theorem turns commutator energy into a genuine distance estimate.  A small value of
\[
 \sum_r\|[A_r,X]\|_{\mathrm{HS}}^2
\]
means little by itself if the spectral gap of the commutant Laplacian is tiny.  The ratio with $\lambda_+$ is the correct stability quantity.

\begin{figure}[H]
\centering
\includegraphics[width=.78\linewidth]{figures/resolution/21_commutant_stability.png}
\caption{The spectral certificate remains above the true distance to the exact joint commutant.  The gap of the commutant Laplacian controls the conversion from residual energy to geometric distance.}
\end{figure}

This result gives a rigorous basis for approximate common-resolution discovery.  One first identifies low-energy eigenvectors of $\mathcal L$, then asks whether they contain nearly idempotent self-adjoint elements that can be rounded to projectors.  The spectral theorem above certifies the first step; the rounding problem requires its own perturbation estimate.

\subsection{New invariants of resolved operators}

A resolved operator carries more information than its total norm.  Two operators may have the same Hilbert--Schmidt energy while distributing that energy very differently across source and target packets.

Normalize the block energies to obtain
\[
 p_{ai}
 =
 \frac{\|Q_aTP_i\|_{\mathrm{HS}}^2}{\|T\|_{\mathrm{HS}}^2}.
\]
Let $p_a^{\mathrm{out}}=\sum_i p_{ai}$ and $p_i^{\mathrm{in}}=\sum_a p_{ai}$.

The transfer entropy is the joint entropy
\[
 H(A,I)=-\sum_{a,i}p_{ai}\log p_{ai}.
\]
A second invariant is the transfer mutual information
\[
 I_{\mathcal Q\leftarrow\mathcal P}(T)
 =
 \sum_{a,i}p_{ai}
 \log\frac{p_{ai}}{p_a^{\mathrm{out}}p_i^{\mathrm{in}}}.
\]

\begin{proposition}[Nonnegativity and unitary invariance]
Transfer mutual information is nonnegative and is invariant under unitaries acting separately inside each source and target packet.
\end{proposition}

\begin{proof}
The expression is the Kullback--Leibler divergence between the joint transfer distribution and the product of its marginals, so it is nonnegative.  Packetwise unitaries preserve every Hilbert--Schmidt block norm and therefore preserve the entire probability table $p_{ai}$.
\end{proof}

Entropy and mutual information answer different questions.  Entropy measures how widely transfer energy is spread.  Mutual information measures how strongly knowledge of the source packet predicts the target packet.  A nearly diagonal operator may have substantial entropy across many diagonal packets and still have high source--target dependence.

\begin{figure}[H]
\centering
\includegraphics[width=.78\linewidth]{figures/resolution/23_transfer_information.png}
\caption{As cross-packet coupling grows, the joint transfer distribution spreads.  Entropy and source--target mutual information respond differently because they measure different structural features.}
\end{figure}

A third invariant concerns how representation cost grows as tolerance shrinks.  Define
\[
 d_T(\varepsilon)
 =
 \min\{k:E_T(k)\le\varepsilon\}.
\]
When the limit exists, set
\[
 \dim_{\mathrm{res}}(T)
 =
 \limsup_{\varepsilon\downarrow0}
 \frac{\log d_T(\varepsilon)}{\log(1/\varepsilon)}.
\]
This is not the dimension of the underlying vector space.  It is the scaling exponent of the representation needed to achieve a prescribed error.  It behaves more like an approximation dimension.

\begin{proposition}[Power-law block profile]
Suppose the sorted block energies satisfy $e_j\asymp j^{-p}$ with $p>1$.  Then
\[
 E(k)\asymp k^{(1-p)/2},
 \qquad
 d_T(\varepsilon)\asymp
 \varepsilon^{-2/(p-1)},
\]
and therefore
\[
 \dim_{\mathrm{res}}(T)=\frac{2}{p-1}.
\]
\end{proposition}

\begin{proof}
The tail of a convergent power series satisfies
\[
 \sum_{j>k}j^{-p}\asymp k^{1-p}.
\]
Taking the square root gives the first estimate.  Solving
\[
 \varepsilon\asymp k^{(1-p)/2}
\]
for $k$ gives the second, and the logarithmic quotient yields the stated dimension.
\end{proof}

\begin{figure}[H]
\centering
\includegraphics[width=.78\linewidth]{figures/resolution/22_effective_resolution_dimension.png}
\caption{Synthetic power-law transfer profiles.  The slope of the log--log complexity curve recovers the predicted resolution dimension.}
\end{figure}

These invariants are task dependent because they depend on a resolution family and a norm.  That dependence is not a defect.  Resolution Theory is explicitly about what a representation reveals.  The important requirement is to state the representation and the comparison class rather than pretending the invariant is coordinate free when it is not.

\subsection{Compression as projection, inclusion, and named residual}

The best finite compression has a simple Hilbert-space structure.  Let $E$ be a selected set of packet edges, and let $\mathcal V_E$ be the span of those Hilbert--Schmidt block subspaces.  The compression map
\[
 C_E:T\mapsto\sum_{(a,i)\in E}Q_aTP_i
\]
is the orthogonal projection onto $\mathcal V_E$.  The inclusion
\[
 J_E:\mathcal V_E\hookrightarrow\mathcal B_2(\mathcal H,\mathcal K)
\]
is its Hilbert-space adjoint:
\[
 \langle C_ET,S\rangle_{\mathrm{HS}}
 =
 \langle T,J_ES\rangle_{\mathrm{HS}}.
\]

Moreover,
\[
 C_EJ_E=I_{\mathcal V_E},
\qquad
 J_EC_E\ne I
\]
in general.  The defect is not anonymous:
\[
 R_E(T)=T-J_EC_ET.
\]

\begin{figure}[H]
\centering
\includegraphics[width=.84\linewidth]{figures/resolution/24_compression_retraction.png}
\caption{Compression is safe when the omitted component remains an explicit residual object.  Exact recovery is obtained by adding it back; approximate work carries its norm or structure as a certificate.}
\end{figure}

This is the operator-theoretic form of a broader principle.  A compressed mathematical object should come with a map back into the expanded space and a named defect measuring what that map does not recover.  In proof theory the defect is a lost discovery neighborhood.  In numerical analysis it is a residual.  In data compression it is reconstruction error.  The common structure is projection, inclusion, and retained information about the complement.

\subsection{Interfaces with categorical probability and quantum logic}

Resolution Theory does not need to replace existing categorical treatments of probability or quantum theory.  It can add one layer that those theories do not ordinarily make primary: a mutable, recursively refinable packet geometry attached to states, effects, and channels.

For a Markov kernel $K:X\rightsquigarrow Y$, source and target resolutions give transfer pieces
\[
 K_{ai}=Q_aKP_i.
\]
A refinement asks whether conditioning and packetization commute.  They do under sufficient-statistic or lumpability hypotheses and fail in general.  The failure itself can be measured as a conditional residual.

For quantum effects, sharp projectors are only one class of predicates.  A projector atlas can coexist with POVMs through Naimark dilation or sequential measurement, but noncommuting sharp predicates do not acquire a fictitious common refinement.  Resolution logic should therefore distinguish:
\begin{itemize}
\item exact sharp refinement;
\item sequentially compatible observation;
\item approximate joint measurability;
\item chart changes between incompatible contexts.
\end{itemize}

The finite semantics developed later makes this indexed picture explicit: each resolution carries a Boolean predicate fiber and a quantitative certificate fiber, while refinement and chart change induce reindexing operations.  Infinite-dimensional completeness and a genuine topos semantics remain open.

\subsection{From a candidate universal property to a theorem}

The original finite theory ended with a question: are projector-resolved systems merely examples, or can every abstract system of the relevant analytic kind be represented faithfully by projector corners?  The next section answers that question for unital resolved $C^*$-systems.  It proves a universal extension law for compatible corner maps, constructs faithful regular and GNS representations, preserves every named packet and refinement, and states the exact boundary at which tensor and higher-categorical realization become separate problems.

\subsection{Where the present theory stops}

The finite axioms are strong enough to support the theorems above, but they do not settle every question suggested by them.  In particular, they do not prove the following stronger statements.

\begin{enumerate}[label=\arabic*.]
\item They do not guarantee that a useful resolution exists for every operator family.
\item They do not guarantee that refinement lowers computational cost.
\item They do not guarantee inverse closure for an arbitrary weighted transfer class.
\item They do not produce a topos until an admissible site, covers, pullbacks, and sheaf axioms are specified.
\item They do not make arbitrary infinite block sums norm convergent.
\item They do not turn noncommuting projectors into a common exact resolution.
\item They do not identify one universally optimal resolution independent of task and error tolerance.
\end{enumerate}

These are not footnotes added for caution.  They are the actual edge of the present theory.  Keeping that edge visible makes the next problems easier to state honestly.

\subsection{Where the theory should go next}

The next stage should not begin by naming a grander structure and hoping the details follow.  It should proceed theorem by theorem, in roughly the following order.

\paragraph{1. Countable and unbounded categorical representation.}
Extend the resolved $C^*$-algebraic theorem to countable object systems, strict or strong projector sums, unbounded transfers, and local linking algebras with declared convergence and domain control.

\paragraph{2. Resolution site and sheaf theory.}
Specify admissible local resolutions and covers, then prove when resolved models form a stack or topos-like category.

\paragraph{3. Spectral invariance.}
Determine geometric and weight conditions under which a weighted transfer algebra is inverse closed.

\paragraph{4. Learned resolution stability.}
For data-driven projectors, prove perturbation, generalization, and leakage bounds rather than relying only on empirical packet quality.

\paragraph{5. Monoidal and higher representation.}
Characterize when a resolved tensor category has a unitary fiber functor, and determine the higher linking data required beyond an ordinary $C^*$-algebra.

\paragraph{6. Infinite Resolution Logic.}
Extend the finite indexed logic proved later to measurable or countable resolutions, sheaf semantics, and a completeness theorem for the analytic certificate layer.

\subsection{Closing perspective}

Resolution Theory begins with an elementary fact: every operator can be written in blocks.  The subject begins only when the block structure is allowed to vary and the choice of resolution is kept inside the mathematics.

The linear finite theory can be summarized in one sentence:

\begin{quote}
A resolved mathematical system is an object together with a coherent family of packet spaces, transfer morphisms, refinement maps, compression projections, residuals, and reconstruction laws.
\end{quote}

Seen this way, several familiar ideas line up without being collapsed into one another.  Operator corners provide the transfer algebra.  Refinement supplies a second categorical direction.  Commutants tell us when exact common structure exists.  Atlases handle the cases in which it does not.  Conditional expectations make compression precise.  Transfer entropy records what a finer resolution reveals.  Certified residuals carry approximation honestly through later calculations.

The theory is not an instruction to divide every object into smaller and smaller pieces.  Its guiding question is more selective and more human: which resolution makes the structure we care about visible, what does that resolution leave out, and how can the omitted part be recovered, bounded, or kept available for another purpose?



\section{Universal representation: when an abstract resolution can be made concrete}
\label{sec:universal-resolution-representation}

The first unresolved question in the original finite theory was not whether projector corners work.  They plainly do.  The question was whether they are merely one convenient construction or whether every abstract finite resolution system of the relevant kind can be realized in that form.  Without a converse theorem, the calculus remains an example.  With one, projector corners become a concrete model of an abstract structure.

The answer is affirmative for resolved $C^*$-systems.  The result is a resolved form of the Gelfand--Naimark representation principle and of the standard linking-algebra picture for $C^*$-categories.  The underlying representation theorems are classical; what is added here is the explicit preservation of object packets, mixed transfer corners, recursive refinement, and the universal extension law that Resolution Theory needs.

\begin{figure}[H]
\centering
\includegraphics[width=.94\linewidth]{figures/universal_representation/01_representation_pipeline.png}
\caption{The abstract system is first assembled into one linking algebra.  A faithful representation then turns the abstract object and packet idempotents into genuine orthogonal projections.  The last arrow is not a change of meaning: it is a concrete realization of the same multiplication, dagger, and refinement laws.}
\end{figure}

\subsection{The precise abstract object}

The cleanest starting point is a unital $C^*$-algebra carrying two nested levels of chosen projections.  The first level identifies mathematical objects.  The second identifies packets inside each object.

\begin{definition}[Resolved linking system]
A \emph{resolved linking system} is a tuple
\[
\mathfrak A=
\left(
A,
\{p_x\}_{x\in X},
\{e_{x,i}\}_{i\in I_x},
\mathfrak F
\right)
\]
with the following properties.
\begin{enumerate}[label=\textup{(R\arabic*)}]
\item $A$ is a unital complex $C^*$-algebra.
\item The object projections are mutually orthogonal and complete:
\[
p_xp_y=\delta_{xy}p_x,
\qquad
\sum_{x\in X}p_x=1_A.
\]
\item For each object $x$, the packet projections are mutually orthogonal and resolve $p_x$:
\[
e_{x,i}e_{x,j}=\delta_{ij}e_{x,i},
\qquad
\sum_{i\in I_x}e_{x,i}=p_x.
\]
\item Every declared refinement in $\mathfrak F$ replaces one packet by a finite orthogonal family:
\[
e_{x,i}=\sum_{\alpha\in\mathcal C(x,i)}e_{x,i,\alpha}.
\]
\end{enumerate}
The transfer space from object $x$ to object $y$ is the corner
\[
\mathcal C(x,y)=p_yAp_x,
\]
and the packet transfer from $(x,i)$ to $(y,a)$ is
\[
\mathcal C((x,i),(y,a))
=
e_{y,a}Ae_{x,i}.
\]
\end{definition}

This definition is deliberately modest.  It does not assume that the object projections are central.  It does not require different packet charts to commute.  It asks only for one resolved linking algebra at a time.  Noncommuting charts may all live inside the same represented algebra, but they need not possess a common orthogonal refinement.

\begin{proposition}[The corner category]
A resolved linking system determines a finite or small $C^*$-category whose objects are the $p_x$ and whose morphism spaces are $p_yAp_x$.  Composition is multiplication in $A$ and the dagger is the $C^*$-adjoint.
\end{proposition}

\begin{proof}
If $a\in p_yAp_x$ and $b\in p_zAp_y$, then
\[
ba=p_zbp_y\,p_yap_x=p_z(ba)p_x\in p_zAp_x.
\]
Associativity is inherited from $A$.  The identity at $x$ is $p_x$.  Finally,
\[
a^*=(p_yap_x)^*=p_xa^*p_y\in p_xAp_y,
\]
so the dagger reverses source and target.  The norm and $C^*$-identity are inherited from $A$.
\end{proof}

The point of this step is that a many-object transfer theory and a one-object linking algebra contain the same finite corner data.  The category remembers the typing.  The algebra makes a faithful Hilbert-space representation available.

\begin{figure}[H]
\centering
\includegraphics[width=.72\linewidth]{figures/universal_representation/02_linking_algebra_corners.png}
\caption{A linking algebra is a matrix whose $(y,x)$ corner is the transfer space from object $x$ to object $y$.  The off-diagonal blocks are not auxiliary; they are the morphisms that connect the object algebras.}
\end{figure}

\subsection{The universal property of the linking algebra}

The word ``universal'' should be earned by an extension property, not by emphasis.  The relevant property is elementary and exact.

\begin{theorem}[Universal extension of compatible corner maps]
Let $\mathfrak A$ be a resolved linking system with finite object set $X$.  Let $K$ be a Hilbert space with mutually orthogonal projections $Q_x$ satisfying
\[
\sum_xQ_x=I_K.
\]
Suppose that for every pair $x,y$ there is a linear map
\[
F_{yx}:p_yAp_x\longrightarrow Q_y\mathcal B(K)Q_x
\]
such that
\begin{align*}
F_{zy}(b)F_{yx}(a)&=F_{zx}(ba),\\
F_{xy}(a^*)&=F_{yx}(a)^*,\\
F_{xx}(p_x)&=Q_x.
\end{align*}
Then there is a unique unital $*$-homomorphism
\[
\widetilde F:A\longrightarrow\mathcal B(K)
\]
whose restriction to every corner is $F_{yx}$.  It is given by
\[
\boxed{
\widetilde F(a)
=
\sum_{x,y}F_{yx}(p_yap_x).
}
\]
\end{theorem}

\begin{proof}
Because $\sum_xp_x=1_A$, every $a\in A$ has the unique corner expansion
\[
a=\sum_{x,y}p_yap_x.
\]
The formula therefore defines the only possible linear extension.  For $a,b\in A$,
\begin{align*}
\widetilde F(b)\widetilde F(a)
&=
\sum_{z,y}\sum_{y',x}
F_{zy}(p_zbp_y)
F_{y'x}(p_{y'}ap_x).
\end{align*}
The factors vanish unless $y=y'$ because their supports lie in orthogonal object projections.  Compatibility with corner composition then gives
\begin{align*}
\widetilde F(b)\widetilde F(a)
&=
\sum_{z,y,x}
F_{zx}(p_zbp_yap_x)\\
&=
\sum_{z,x}
F_{zx}(p_zbap_x)\\
&=
\widetilde F(ba).
\end{align*}
The dagger law follows by reversing the corners:
\[
\widetilde F(a)^*
=
\sum_{x,y}F_{xy}(p_xa^*p_y)
=
\widetilde F(a^*).
\]
Finally,
\[
\widetilde F(1_A)
=
\sum_xF_{xx}(p_x)
=
\sum_xQ_x
=I_K.
\]
Thus $\widetilde F$ is a unital $*$-homomorphism, and uniqueness was already forced by the corner decomposition.
\end{proof}

\begin{figure}[H]
\centering
\includegraphics[width=.88\linewidth]{figures/universal_representation/11_universal_extension.png}
\caption{The linking algebra is universal for compatible corner assignments.  Once the action on every typed transfer space is fixed and the actions respect composition and dagger, there is exactly one action of the whole algebra.}
\end{figure}

This theorem answers one half of the problem.  If a concrete corner model exists, its pieces glue uniquely.  The remaining half is existence.

\subsection{A canonical faithful representation}

For finite-dimensional systems there is a particularly transparent construction.  It is not always the smallest representation, but it is canonical enough to prove existence and easy enough to audit line by line.

\begin{lemma}[Faithful trace in finite dimension]
Every finite-dimensional $C^*$-algebra admits a faithful tracial state.
\end{lemma}

\begin{proof}
The finite-dimensional structure theorem gives
\[
A\cong\bigoplus_{r=1}^m M_{n_r}(\mathbb C).
\]
Choose positive weights $c_r>0$ with $\sum_rc_r=1$ and define
\[
\tau(a_1,\ldots,a_m)
=
\sum_{r=1}^m
c_r\frac{\operatorname{Tr}(a_r)}{n_r}.
\]
This is tracial and positive.  If $a\ne0$, then some $a_r^*a_r\ne0$, so
\[
\tau(a^*a)>0.
\]
Hence $\tau$ is faithful.
\end{proof}

\begin{theorem}[Faithful left-regular representation]
Let $A$ be a finite-dimensional unital $C^*$-algebra and let $\tau$ be a faithful tracial state.  Equip the vector space $A$ with
\[
\langle a,b\rangle_\tau
=
\tau(a^*b).
\]
For $x\in A$, define left multiplication
\[
\lambda(x)a=xa.
\]
Then
\[
\lambda:A\longrightarrow\mathcal B(L^2(A,\tau))
\]
is a faithful unital $*$-representation.
\end{theorem}

\begin{proof}
Associativity gives
\[
\lambda(xy)=\lambda(x)\lambda(y),
\qquad
\lambda(1_A)=I.
\]
Using the trace property,
\begin{align*}
\langle\lambda(x)a,b\rangle_\tau
&=
\tau((xa)^*b)\\
&=
\tau(a^*x^*b)\\
&=
\langle a,\lambda(x^*)b\rangle_\tau.
\end{align*}
Thus $\lambda(x)^*=\lambda(x^*)$.  If $\lambda(x)=0$, then
\[
x=\lambda(x)1_A=0,
\]
so the representation is faithful.
\end{proof}

The argument is short because the algebra already contains its own regular module.  The representation does not invent a Hilbert space from outside.  It completes the algebra in an inner product supplied by a faithful positive functional and lets each element act by multiplication.

\begin{figure}[H]
\centering
\includegraphics[width=.86\linewidth]{figures/universal_representation/06_faithfulness_singular_values.png}
\caption{In the finite cyclic examples, the represented basis operators remain linearly independent.  Their positive singular values are the numerical signature of faithfulness in the chosen coordinates.}
\end{figure}

\subsection{The universal projector-corner representation theorem}

The finite regular construction is enough for computation, but the conceptual statement is broader.  The ordinary Gelfand--Naimark theorem says that every $C^*$-algebra has a faithful representation on a Hilbert space.  Applying it without forgetting the chosen projections gives the resolved theorem.

\begin{theorem}[Universal projector-corner representation]
\label{thm:universal-projector-corner}
Let
\[
\mathfrak A=
\left(A,\{p_x\},\{e_{x,i}\},\mathfrak F\right)
\]
be a unital resolved linking system.  There exist a Hilbert space $\mathcal H$ and a faithful unital $*$-representation
\[
\pi:A\longrightarrow\mathcal B(\mathcal H)
\]
such that:
\begin{enumerate}[label=\textup{(\roman*)}]
\item the object elements become a complete orthogonal projector family,
\[
P_x=\pi(p_x),
\qquad
P_xP_y=\delta_{xy}P_x,
\qquad
\sum_xP_x=I_\mathcal H;
\]
\item the packet elements become orthogonal subprojectors,
\[
B_{x,i}=\pi(e_{x,i}),
\qquad
\sum_iB_{x,i}=P_x;
\]
\item every abstract transfer corner is represented exactly as an operator corner,
\[
\boxed{
\pi(p_yAp_x)
=
P_y\pi(A)P_x;
}
\]
\item every packet transfer is represented exactly,
\[
\boxed{
\pi(e_{y,a}Ae_{x,i})
=
B_{y,a}\pi(A)B_{x,i};
}
\]
\item every declared refinement is preserved,
\[
e_{x,i}=\sum_\alpha e_{x,i,\alpha}
\quad\Longrightarrow\quad
B_{x,i}=\sum_\alpha B_{x,i,\alpha};
\]
\item reconstruction, composition, and dagger are preserved:
\[
\pi(a)
=
\sum_{x,y,i,a}
B_{y,a}\pi(a)B_{x,i},
\]
\[
\pi(ba)=\pi(b)\pi(a),
\qquad
\pi(a^*)=\pi(a)^*.
\]
\end{enumerate}
If $A$ is finite-dimensional, $\pi$ may be chosen to be the left-regular representation above.
\end{theorem}

\begin{proof}
Choose any faithful unital $*$-representation of $A$; the universal GNS representation is one canonical choice.  Because $*$-homomorphisms preserve products, sums, adjoints, and the unit, every projection remains a projection and every orthogonality relation remains exact.  Thus
\[
P_xP_y
=
\pi(p_xp_y)
=
\delta_{xy}\pi(p_x),
\]
and
\[
\sum_xP_x
=
\pi\left(\sum_xp_x\right)
=
\pi(1_A)
=
I_\mathcal H.
\]
The same argument applies to the packets and their refinements.

For $a\in p_yAp_x$,
\[
\pi(a)
=
\pi(p_yap_x)
=
P_y\pi(a)P_x,
\]
so
\[
\pi(p_yAp_x)
\subseteq
P_y\pi(A)P_x.
\]
Conversely, for any $c\in A$,
\[
P_y\pi(c)P_x
=
\pi(p_ycp_x)
\in
\pi(p_yAp_x).
\]
Hence the two corners are equal.  Replacing $p_x,p_y$ by packet projections gives the packet-corner identity.

Finally,
\begin{align*}
\pi(a)
&=
\pi\left(
\sum_{x,y,i,a}
e_{y,a}ae_{x,i}
\right)\\
&=
\sum_{x,y,i,a}
B_{y,a}\pi(a)B_{x,i}.
\end{align*}
Faithfulness guarantees that no abstract equality or inequality is collapsed by the representation.  Composition and dagger are preserved because $\pi$ is a $*$-homomorphism.
\end{proof}

This is the first abstract representation theorem Resolution Theory needed.  It says that the projector-corner language is not confined to systems that happened to be presented as matrices at the beginning.  Every unital resolved $C^*$-system has a faithful concrete model of exactly that kind.

\begin{corollary}[Resolved $C^*$-categories]
Every finite resolved $C^*$-category whose linking algebra is formed in the standard way admits a faithful concrete realization by Hilbert spaces and projector corners.  Its selected object and packet idempotents become orthogonal projections, and its morphisms become the corresponding operator corners.
\end{corollary}

\begin{proof}
Apply Theorem~\ref{thm:universal-projector-corner} to the linking algebra.  The object identities become the object projections, and each categorical hom-space is the corresponding linking corner.
\end{proof}

This corollary sits beside established representation results for $C^*$- and $W^*$-categories.  The present contribution is not the discovery that such categories can be represented concretely; it is the resolved formulation that retains the chosen packet hierarchy and makes its universal corner law explicit.  Standard categorical representation results provide the surrounding context \cite{GLR1985,Heunen2009,FritzWesterbaan2020}.

\subsection{A recognition theorem}

The previous theorem begins with an abstract resolved algebra and produces a concrete projector system.  The converse direction is immediate but worth stating, because together the two directions identify the exact abstract structure being modeled.

\begin{theorem}[Finite recognition theorem]
Consider the following two constructions.
\begin{enumerate}[label=\textup{(\alph*)}]
\item From a resolved linking system, choose a faithful representation and form its object and packet projectors.
\item From a concrete projector-corner system
\[
\left(
\mathcal H,
A\subseteq\mathcal B(\mathcal H),
\{P_x\},
\{B_{x,i}\}
\right),
\]
forget the ambient Hilbert coordinates but retain the abstract $C^*$-algebra $A$, its object projections, packet projections, and refinement relations.
\end{enumerate}
Then every abstract resolved linking system is isomorphic to the system recovered from any faithful concrete realization.  Conversely, every concrete projector-corner system gives a resolved linking system.
\end{theorem}

\begin{proof}
The second construction plainly satisfies the resolved linking axioms.  For the first direction, let $\pi$ be faithful.  The map
\[
a\longmapsto\pi(a)
\]
is an isomorphism from $A$ onto the concrete $C^*$-algebra $\pi(A)$.  It sends every named object and packet projection to its concrete counterpart and preserves all declared refinement equations.  Thus forgetting the coordinates recovers an isomorphic resolved linking system.
\end{proof}

The theorem does not say that all faithful realizations are literally the same Hilbert-space model.  They need not have the same dimension or multiplicities.  It says that they realize the same abstract resolved algebra, and that every compatible concrete model factors through that algebra by the universal extension theorem.

\subsection{Refinement survives representation exactly}

A representation theorem would be too weak for Resolution Theory if it preserved the coarse algebra but blurred the chosen refinement hierarchy.  Faithful $*$-representations do not do that.  Orthogonal sums remain orthogonal sums.

\begin{proposition}[Refinement naturality]
Let
\[
e=\sum_{\alpha=1}^r e_\alpha
\]
be a finite orthogonal refinement in a resolved linking system.  Under every unital $*$-representation,
\[
\pi(e)=\sum_{\alpha=1}^r\pi(e_\alpha),
\qquad
\pi(e_\alpha)\pi(e_\beta)=\delta_{\alpha\beta}\pi(e_\alpha).
\]
If $a\in A$, then the coarse transfer block is exactly the sum of its fine transfer blocks:
\[
\pi(eae')
=
\sum_{\alpha,\beta}
\pi(e_\alpha a e'_\beta).
\]
\end{proposition}

\begin{proof}
Linearity and multiplicativity of $\pi$ preserve the sum and orthogonality relations.  Expanding $eae'$ by the two finite resolutions gives the last identity before or after applying $\pi$.
\end{proof}

\begin{figure}[H]
\centering
\includegraphics[width=.88\linewidth]{figures/universal_representation/10_refinement_preservation.png}
\caption{Representation and refinement commute.  A coarse idempotent and its family of fine idempotents become a coarse orthogonal projector and the corresponding fine projector family without changing the refinement equation.}
\end{figure}

\subsection{Minimal and regular realizations}

The existence theorem does not claim that the left-regular representation is computationally optimal.  In fact, it often is not.

\begin{theorem}[Minimal faithful dimension in finite dimension]
Let
\[
A\cong\bigoplus_{r=1}^mM_{n_r}(\mathbb C).
\]
The smallest Hilbert-space dimension of a faithful $*$-representation of $A$ is
\[
\boxed{
\dim\mathcal H_{\min}
=
\sum_{r=1}^m n_r.
}
\]
The left-regular representation acts on a space of dimension
\[
\boxed{
\dim L^2(A,\tau)
=
\sum_{r=1}^m n_r^2.
}
\]
\end{theorem}

\begin{proof}
Every irreducible representation of $M_{n_r}(\mathbb C)$ has dimension $n_r$.  A representation of the direct sum is faithful only if it is nonzero on every simple summand, so it must contain at least one copy of every defining irreducible representation.  Hence its dimension is at least $\sum_rn_r$.  The direct sum of the defining representations attains that bound.

The underlying vector space of $A$ has dimension $\sum_rn_r^2$, and the left-regular representation acts on its Hilbert completion, which has the same finite dimension.
\end{proof}

Define the regular-representation overhead by
\[
\kappa_{\mathrm{reg}}(A)
=
\frac{\sum_rn_r^2}{\sum_rn_r}.
\]
The regular model is valuable because it is direct and canonical.  A minimal model is valuable because it is smaller.  Resolution Theory should distinguish these purposes rather than treating one representation as universally preferable.

\begin{figure}[H]
\centering
\includegraphics[width=.78\linewidth]{figures/universal_representation/04_regular_vs_minimal_dimension.png}
\caption{For a single $M_n(\mathbb C)$ block, the standard faithful representation has dimension $n$, whereas the left-regular representation has dimension $n^2$.  The theorem proves existence; a later planner may compress the realization.}
\end{figure}

\begin{figure}[H]
\centering
\includegraphics[width=.78\linewidth]{figures/universal_representation/05_representation_overhead.png}
\caption{The regular-to-minimal dimension ratio grows linearly for one full matrix block.  This is not a defect in the theorem.  It is the price of choosing a canonical existence proof rather than an optimized realization.}
\end{figure}

\begin{proposition}[Universality does not imply runtime optimality]
For $A=M_n(\mathbb C)$, applying a represented element in the regular model is an operation on $n^2$ coordinates, whereas the minimal defining representation acts on $n$ coordinates.  Therefore the universal construction may be asymptotically and practically more expensive even though it preserves the same algebra exactly.
\end{proposition}

The benchmark included with this appendix makes the point deliberately.  It is a negative control against the temptation to identify categorical universality with computational efficiency.

\begin{figure}[H]
\centering
\includegraphics[width=.80\linewidth]{figures/universal_representation/09_representation_runtime.png}
\caption{Reference CPU timings for the same algebra element in minimal and left-regular representations.  The regular model is the cleanest existence construction; it is not the fastest representation of a full matrix algebra.}
\end{figure}

\subsection{Four worked examples}

\subsubsection*{A single matrix algebra}

Let $A=M_3(\mathbb C)$ with packet projections
\[
e_1=E_{11}+E_{22},
\qquad
e_2=E_{33}.
\]
The defining representation on $\mathbb C^3$ is minimal.  The left-regular representation acts on the nine-dimensional Hilbert-Schmidt space.  In both realizations,
\[
A
=
e_1Ae_1
\oplus
e_1Ae_2
\oplus
e_2Ae_1
\oplus
e_2Ae_2
\]
as a vector-space decomposition, and multiplication joins the corners by the ordinary intermediate-packet rule.

\subsubsection*{A two-object category}

Take object dimensions two and one.  The linking algebra is $M_3(\mathbb C)$ with
\[
p_X=
\begin{pmatrix}
1&0&0\\
0&1&0\\
0&0&0
\end{pmatrix},
\qquad
p_Y=
\begin{pmatrix}
0&0&0\\
0&0&0\\
0&0&1
\end{pmatrix}.
\]
Then
\[
p_YAp_X\cong M_{1\times2}(\mathbb C)
\]
is the space of transfers from $X$ to $Y$, and
\[
p_XAp_Y\cong M_{2\times1}(\mathbb C)
\]
is the reverse transfer space.  The category is not an analogy added to the matrix algebra.  It is already present in the corners.

\subsubsection*{A cyclic group algebra}

For $A=\mathbb C[C_m]$, left multiplication gives the regular representation by cyclic permutation matrices.  Fourier character projectors diagonalize the commuting group action.  The same abstract algebra therefore admits two especially useful concrete views:
\[
\text{group-coordinate packets}
\qquad\text{and}\qquad
\text{character packets}.
\]
They are related by a finite Fourier transform.  Universality does not select one of them as the only correct resolution.

\subsubsection*{A resolved algebra with incompatible charts}

Suppose $\{e_i\}$ and $\{f_a\}$ are two projection families in the same algebra and some commutators do not vanish.  The representation theorem realizes all of them as genuine projections on one Hilbert space.  It does not turn $e_if_a$ into a projection, and it does not manufacture a common orthogonal refinement.  The proper object remains a projector atlas with mixed corners
\[
\pi(f_a)\pi(A)\pi(e_i).
\]
This is an important success condition: the representation theorem preserves incompatibility instead of hiding it.

\subsection{Approximate representations and audit defects}

A numerical implementation may preserve the algebra only approximately.  Given basis elements $u_i$ and proposed matrices $R_i$, define
\begin{align*}
\delta_{\mathrm{mult}}
&=
\max_{i,j}
\left\|
R_iR_j-
\sum_kc_{ij}^kR_k
\right\|,\\
\delta_*
&=
\max_i
\left\|
R_i^*-
\sum_ks_i^kR_k
\right\|,\\
\delta_1
&=
\left\|
\sum_iu_i^{(1)}R_i-I
\right\|.
\end{align*}
The aggregate representation defect is
\[
\delta_{\mathrm{rep}}
=
\max\{\delta_{\mathrm{mult}},\delta_*,\delta_1\}.
\]

The finite theorem concerns exact representations.  The defect does not by itself prove that an approximate representation lies near an exact one; that is a separate stability problem.  It does provide an honest implementation audit.

\begin{figure}[H]
\centering
\includegraphics[width=.78\linewidth]{figures/universal_representation/08_representation_defect.png}
\caption{Perturbing the basis matrices produces measurable multiplication, dagger, and identity defects.  These diagnostics tell us how badly a numerical realization has departed from the abstract laws; they do not silently relabel the perturbed system as exact.}
\end{figure}

\subsection{What happens in infinite dimension}

The finite regular representation was included because it is constructive and transparent.  The resolved representation theorem itself is not confined to finite-dimensional algebras.

\begin{theorem}[Universal GNS realization]
Let $A$ be a unital $C^*$-algebra with a finite resolved projection family.  Let
\[
\pi_u
=
\bigoplus_{\omega\in S(A)}\pi_\omega
\]
be the direct sum of the GNS representations over all states of $A$.  Then $\pi_u$ is faithful, and Theorem~\ref{thm:universal-projector-corner} holds with $\pi=\pi_u$.
\end{theorem}

\begin{proof}
If $a\ne0$, then $a^*a\ne0$.  States separate the positive cone of a $C^*$-algebra, so there is a state $\omega$ with
\[
\omega(a^*a)>0.
\]
In the GNS representation of that state,
\[
\|\pi_\omega(a)[1]\|^2
=
\omega(a^*a)>0,
\]
so $\pi_\omega(a)\ne0$.  Hence the direct sum over all states is faithful.  The projection and corner conclusions follow exactly as in the proof of Theorem~\ref{thm:universal-projector-corner}.
\end{proof}

\begin{figure}[H]
\centering
\includegraphics[width=.88\linewidth]{figures/universal_representation/13_universal_gns.png}
\caption{The universal GNS representation separates algebra elements by allowing every state to contribute a Hilbert-space view.  The direct sum is faithful because a nonzero positive element is detected by at least one state.}
\end{figure}

This settles the unital $C^*$-algebraic existence problem.  It does not yet settle every infinite Resolution Theory question.  Countably many object projections may sum only strongly or strictly.  Unbounded operators require domains or quadratic forms.  Higher categorical resolutions require more than one linking algebra.  Those issues belong to the next stage rather than being hidden inside the theorem just proved.

\subsection{The tensor boundary}

A faithful corner representation preserves addition, multiplication, adjoint, norms, projections, and refinement.  It does not automatically preserve a separate monoidal product carried by an abstract category.

\begin{proposition}[Monoidal enhancement criterion]
Let $\mathcal C$ be a resolved monoidal $C^*$-category.  A concrete realization preserving the tensor product by unitary coherent isomorphisms is exactly a faithful unitary strong monoidal $*$-functor
\[
F:\mathcal C\longrightarrow\mathrm{Hilb}.
\]
Thus a monoidal corner realization exists whenever such a fiber functor exists; ordinary corner representability alone does not supply it.
\end{proposition}

\begin{proof}
A concrete monoidal realization assigns Hilbert spaces to objects, bounded operators to morphisms, and unitary coherence maps
\[
F(X)\otimes F(Y)\cong F(X\otimes Y)
\]
compatible with associators, units, dagger, and composition.  This is precisely a unitary strong monoidal $*$-functor.  Conversely, any such functor is a concrete monoidal realization by definition.
\end{proof}

\begin{figure}[H]
\centering
\includegraphics[width=.90\linewidth]{figures/universal_representation/12_monoidal_boundary.png}
\caption{Ordinary $C^*$-categorical representation is automatic under the stated hypotheses.  Tensor-preserving representation is a separate existence problem governed by a unitary fiber functor.}
\end{figure}

This boundary matters.  It prevents the universal theorem from being stretched into a claim that every abstract tensor category already has a Hilbert-space fiber functor.  The representation theorem solves the first foundational item while leaving a sharper second problem visible.

\subsection{Computational verification}

The accompanying Python module stores a finite $*$-algebra by multiplication and dagger structure constants.  It then constructs the left-regular representation, checks the algebraic laws, recovers coefficients from represented operators, constructs object and packet projections, verifies exact corner reconstruction, and compares regular with minimal representations.

The tests cover:
\begin{enumerate}[label=\arabic*.]
\item associativity, unit, dagger, and involution for matrix-unit algebras;
\item exactness and faithfulness of the left-regular representation;
\item cyclic group algebras;
\item direct sums of matrix blocks;
\item orthogonal object projections;
\item typed corner support;
\item refinement preservation;
\item exact packet reconstruction;
\item uniqueness of the universal linear extension;
\item coefficient recovery;
\item positive faithfulness singular values.
\end{enumerate}

\begin{figure}[H]
\centering
\includegraphics[width=.78\linewidth]{figures/universal_representation/07_corner_reconstruction_error.png}
\caption{The finite reference systems reconstruct random represented elements from all packet corners to machine precision.  The numerical check mirrors the exact identity proved above.}
\end{figure}

\begin{table}[H]
\centering
\small
\begin{tabular}{L{2.0in}L{1.2in}L{2.55in}}
\toprule
Statement & Status & Boundary\\
\midrule
Finite linking algebra construction & proved & finite selected object family\\
Universal extension of compatible corners & proved & exact compositional and dagger compatibility\\
Faithful regular representation & proved & finite-dimensional $C^*$-algebra\\
Universal GNS representation & proved & unital $C^*$-algebra\\
Packet and refinement preservation & proved & finite sums as stated\\
Minimal faithful dimension & proved & finite-dimensional $C^*$-algebras\\
Monoidal representation & conditional & requires a unitary fiber functor\\
Countable-object categorical completion & open extension & strict/strong convergence and multiplier control\\
Higher-categorical realization & open & requires genuinely higher linking data\\
\bottomrule
\end{tabular}
\caption{The representation theorem closes the ordinary abstract-to-concrete gap while preserving the exact boundary of the monoidal and higher-categorical questions.}
\end{table}

\subsection{What has now been resolved}

The original finite theory stopped with a candidate universal property.  The present section replaces that candidate with a theorem:
\[
\boxed{
\begin{minipage}{0.86\linewidth}
Every unital resolved $C^*$-system has a faithful Hilbert-space realization in which objects and packets are orthogonal projections, transfers are operator corners, refinements remain exact, and every compatible family of corner maps extends uniquely through the linking algebra.
\end{minipage}
}
\]

The theorem does not say that every mathematical theory is a $C^*$-system.  It does not make all packet charts commute.  It does not choose the most efficient representation.  It does not provide a tensor fiber functor automatically.  What it does say is decisive enough: within its stated analytic class, projector-resolved mathematics is not merely a concrete example.  It is a faithful realization theorem with a universal linking object.

\subsection{The next unresolved item}

The first item on the foundational list is now complete in the $C^*$-algebraic setting.  The next item is categorical and infinite: to determine when directed families of resolutions, unbounded transfers, countable object systems, and local projector atlases admit a complete representation theory with controlled limits and a sound internal logic.

The finite theorem gives that program a stable base.  Any future extension should reduce to the theorem here when the system is finite, unital, and bounded.

\begin{figure}[H]
\centering
\includegraphics[width=.94\linewidth]{figures/universal_representation/14_theorem_dependency.png}
\caption{The representation theorem is not one isolated statement.  It is the closure of a dependency chain: resolved axioms produce a linking algebra; the linking algebra admits a faithful representation; the named idempotents become projectors; and the universal extension and recognition results follow.}
\end{figure}

\subsection{Transition to nonlinear resolution}

The results above treat a transformation as a family of one-input transfers.  That is the right language for linear maps, but it cannot record the interaction of two or more packets inside a nonlinear map.  The next part keeps the same resolved objects and refinement maps while replacing ordinary edges by multilinear hyperedges.  Nothing in the linear theory is discarded; it becomes the first derivative layer of the nonlinear one.


\section{Nonlinear Resolution Theory: jets, hyperedges, and interaction trees}
\label{sec:nonlinear-resolution-theory}

\subsection{Why the linear calculus is not the end of the story}

The linear transfer identity
\[
        T=\sum_{i,j}Q_iTP_j
\]
answers a remarkably broad question: once the source and target have been resolved into packets, which source packet contributes to which target packet?  That question is enough for linear dynamics, linear filtering, matrix algorithms, quantum channels after linearization, and many local arguments.  It is not enough for a genuinely nonlinear map.

If
\[
        x=\sum_jP_jx,
\]
then one usually has
\[
        F(x)\ne\sum_jF(P_jx).
\]
The missing terms are not errors.  They are the interactions between packets.  A quadratic map sees pairs of packets, a cubic map sees triples, and an analytic map carries an entire hierarchy of interactions.  The correct nonlinear object is therefore not an ordinary directed graph.  It is a directed hypergraph whose edges may have several inputs.

For a smooth map $F:X\to Y$, with source projectors $P_j$ and target projectors $Q_i$, define
\[
F^{[m]}_{i;j_1,\ldots,j_m}(x)
=
Q_iD^mF(x)
\bigl[P_{j_1}\cdot,\ldots,P_{j_m}\cdot\bigr].
\]
At first order this is the familiar transfer edge $j\to i$.  At second order it is a hyperedge
\[
(j_1,j_2)\longrightarrow i,
\]
and higher derivatives continue the same pattern.

\begin{figure}[H]
\centering
\includegraphics[width=.88\linewidth]{figures/nonlinear/01_linear_to_nonlinear_hypergraph.png}
\caption{A linear transfer has one input packet.  A quadratic interaction joins two input packets before transferring their joint contribution to an output packet.}
\end{figure}

The purpose of this part is to extend Resolution Theory without changing its basic commitments.  We make the hierarchy exact, prove that it composes and refines coherently, and identify the cases in which the representation genuinely reduces a nonlinear calculation.  The result is a recursively refinable colored operad of packet maps, together with local inverse, implicit-function, flow, stability, and approximation theorems.

\subsection{Resolved Banach spaces}

Let $X$ and $Y$ be Banach spaces carrying finite bounded resolutions
\[
I_X=\sum_{j\in J}P_j,
\qquad
I_Y=\sum_{i\in I}Q_i.
\]
Write $X_j=P_jX$ and $Y_i=Q_iY$.  Positive weights $a_j$ and $b_i$ define packet maximum norms
\[
\|x\|_a=\max_j\frac{\|P_jx\|}{a_j},
\qquad
\|y\|_b=\max_i\frac{\|Q_iy\|}{b_i}.
\]
The weights can record expected amplitude, spatial scale, energy, uncertainty, or a desired decay law.  They are not decoration: a good choice can turn a formally dense nonlinear map into a contractive one on the physically relevant packet ball.

\begin{definition}[Resolved multilinear map]
For an $m$-linear map $A:X^m\to Y$, its packet block is
\[
A_{i;j_1,\ldots,j_m}
=
Q_iA(P_{j_1}\cdot,\ldots,P_{j_m}\cdot).
\]
The ordered tuple $(j_1,\ldots,j_m)\to i$ is the corresponding nonlinear transfer hyperedge.
\end{definition}

\begin{theorem}[Exact multilinear reconstruction]
Every bounded $m$-linear map satisfies
\[
A(x_1,\ldots,x_m)
=
\sum_{i,j_1,\ldots,j_m}
A_{i;j_1,\ldots,j_m}
(P_{j_1}x_1,\ldots,P_{j_m}x_m).
\]
For finite resolutions the sum is exact and finite.
\end{theorem}

\begin{proof}
Insert the source identities into every input slot and the target identity after applying $A$:
\[
A(x_1,\ldots,x_m)
=
\sum_iQ_iA\left(\sum_{j_1}P_{j_1}x_1,\ldots,\sum_{j_m}P_{j_m}x_m\right).
\]
Multilinearity expands the expression into the displayed sum.  No approximation has occurred.
\end{proof}

\begin{theorem}[Weighted packet bound]
Define
\[
\|A\|^{(m)}_{a\to b}
=
\max_i\frac{1}{b_i}
\sum_{j_1,\ldots,j_m}
\|A_{i;j_1,\ldots,j_m}\|
\prod_{r=1}^ma_{j_r}.
\]
Then
\[
\|A(x_1,\ldots,x_m)\|_b
\le
\|A\|^{(m)}_{a\to b}
\prod_{r=1}^m\|x_r\|_a.
\]
\end{theorem}

\begin{proof}
For every output packet $i$, apply the triangle inequality to the exact reconstruction and use
\[
\|P_jx_r\|\le a_j\|x_r\|_a.
\]
After dividing by $b_i$, the definition of the weighted block norm gives the result.  Taking the maximum over $i$ finishes the proof.
\end{proof}

\begin{figure}[H]
\centering
\includegraphics[width=.88\linewidth]{figures/nonlinear/02_tensor_block_energy.png}
\caption{A dense quadratic tensor can be highly uneven after packet resolution.  The useful object is the distribution of interaction energy across hyperedges, not merely the fact that the ambient tensor is dense.}
\end{figure}

\subsection{Jets and the nonlinear transfer hypergraph}

At a base point $x$, the $r$-jet of a smooth map is
\[
J_x^rF
=
\left(F(x),DF(x),\ldots,D^rF(x)\right).
\]
Resolving each derivative produces a hierarchy
\[
\left\{
F^{[m]}_{i;j_1,\ldots,j_m}(x)
\right\}_{0\le m\le r}.
\]
The constant term is an output packet, the first derivative is an ordinary transfer graph, and every higher derivative adds a layer of hyperedges.  This hierarchy is the nonlinear analogue of a block matrix.

If $D^mF(x)$ is symmetric, then permuting the input labels does not change the underlying multilinear map.  One may store ordered hyperedges, which is simple and direct, or quotient by the symmetric-group action, which reduces storage but introduces multiplicity factors.  The theory permits both conventions, provided the convention is fixed before composition.

\subsection{Composition and the colored symmetric operad}

For smooth maps
\[
X\xrightarrow{F}Y\xrightarrow{G}Z,
\]
the chain rule at first order is matrix-like:
\[
D(G\circ F)(x)=DG(F(x))DF(x).
\]
At second order one already sees the two ways an interaction may be formed:
\[
D^2(G\circ F)(x)[u,v]
=
DG(F(x))D^2F(x)[u,v]
+
D^2G(F(x))[DF(x)u,DF(x)v].
\]
The first term carries an inner quadratic interaction through a linear outer edge.  The second joins two inner linear edges inside one outer quadratic hyperedge.

The general law is Fa\`a di Bruno's formula.

\begin{theorem}[Packet-resolved Fa\`a di Bruno formula]
Let $\Pi_n$ denote the set of partitions of $\{1,\ldots,n\}$.  Then
\[
D^n(G\circ F)(x)[h_1,\ldots,h_n]
=
\sum_{\pi\in\Pi_n}
D^{|\pi|}G(F(x))
\left[
D^{|A_1|}F(x)[h_r]_{r\in A_1},\ldots,
D^{|A_{|\pi|}|}F(x)[h_r]_{r\in A_{|\pi|}}
\right].
\]
After inserting the source, intermediate, and target projectors, every term becomes a rooted tree of packet hyperedges, summed over all intermediate packet labels.
\end{theorem}

\begin{proof}
The ordinary Banach-space Fa\`a di Bruno formula gives the partition sum.  Insert the identity on every source argument before each derivative of $F$, insert the intermediate identity after each derivative of $F$, and insert the target identity after the derivative of $G$.  Multilinearity expands all identities into packet labels.  Each block of the partition tells which input leaves are first combined by an inner derivative.  The outer derivative then joins those inner outputs.  The resulting combinatorial object is exactly a rooted packet tree.
\end{proof}

\begin{figure}[H]
\centering
\includegraphics[width=.78\linewidth]{figures/nonlinear/03_faa_di_bruno_growth.png}
\caption{The number of set partitions grows rapidly with derivative order.  Packet sparsity matters because it can remove most rooted-tree substitutions before they are evaluated.}
\end{figure}

\begin{proposition}[Colored symmetric operad]
Take packet spaces as colors and continuous multilinear packet maps as operations.  Permutation of inputs gives the symmetric action, identity maps give operadic units, and rooted-tree substitution gives operadic composition.  These operations satisfy the colored symmetric operad axioms.
\end{proposition}

\begin{proof}
Associativity follows from associativity of substitution of multilinear maps, equivalently from associativity of composition of smooth maps and the uniqueness of their derivatives.  Identity operations are the packet identity maps.  Permutation compatibility is inherited from the natural symmetric-group action on multilinear inputs.  The packet labels merely type the allowed input and output spaces; they do not alter the standard operad identities.
\end{proof}

This statement is more than a change of vocabulary.  It identifies the nonlinear closure law: linear paths become rooted interaction trees, and composition is substitution of one tree into the leaves of another.

\subsection{Recursive refinement}

Suppose a source packet is refined as
\[
P_j=\sum_{\alpha\in\mathcal C(j)}P_{j,\alpha},
\]
and similarly for the target packets.  Every nonlinear block refines exactly:
\[
F^{[m]}_{i;j_1,\ldots,j_m}
=
\sum_{\gamma,\alpha_1,\ldots,\alpha_m}
F^{[m]}_{(i,\gamma);(j_1,\alpha_1),\ldots,(j_m,\alpha_m)}.
\]

\begin{theorem}[Refinement coherence]
Refining a nonlinear jet, composing the refined jets, and then coarsening gives the same jet as composing first and resolving at the coarse level.
\end{theorem}

\begin{proof}
Both procedures insert complete resolutions of the identity into the same Fa\`a di Bruno formula.  The refined calculation keeps the additional labels until the final coarsening sum.  The coarse calculation performs those sums immediately.  Finite distributivity makes the two expressions identical.
\end{proof}

\begin{figure}[H]
\centering
\includegraphics[width=.70\linewidth]{figures/nonlinear/04_refinement_energy_conservation.png}
\caption{Refinement creates more hyperedges but does not create nonlinear energy.  In a coordinate realization, the sum of squared block norms is conserved exactly.}
\end{figure}

\begin{proposition}[No-free-refinement principle]
If every refined hyperedge is active and every contribution must be evaluated, refinement does not reduce the exact arithmetic count.  A computational gain requires additional structure: vanished blocks, low rank, repeated coefficients, selective outputs, transform structure, or reuse.
\end{proposition}

\begin{proof}
A coarse tensor entry belongs to exactly one refined block.  Summing the entry counts of all refined blocks therefore reproduces the entry count of the original tensor.  Refinement changes organization, not the number of scalar coefficients.  Work decreases only when the refined organization allows coefficients or contractions to be omitted or represented more economically.
\end{proof}

\subsection{Weighted analytic packet spaces}

Let
\[
F(x)=\sum_{m\ge0}A_m[x,\ldots,x]
\]
be analytic on a packet ball, with $A_m$ symmetric $m$-linear.  Define the scalar majorant
\[
\mathfrak m_F(t)
=
\sum_{m\ge0}\|A_m\|^{(m)}t^m.
\]
The weighted multilinear bound immediately gives
\[
\|F(x)\|\le\mathfrak m_F(\|x\|).
\]

\begin{theorem}[Analytic composition bound]
If $F(0)=0$ and $\mathfrak m_F(t)$ lies inside the convergence interval of $\mathfrak m_G$, then
\[
\mathfrak m_{G\circ F}(t)
\le
\mathfrak m_G\!\left(\mathfrak m_F(t)\right).
\]
\end{theorem}

\begin{proof}
Apply the packet multilinear bound to every derivative block of $G$ evaluated on the analytic expansion of $F$.  Every coefficient of the resulting rooted-tree sum is nonnegative after taking norms.  Grouping terms by total degree gives exactly the scalar composition of the two majorant series, possibly with cancellations removed.  Removing cancellations can only enlarge the bound.
\end{proof}

\begin{figure}[H]
\centering
\includegraphics[width=.82\linewidth]{figures/nonlinear/05_majorant_composition.png}
\caption{The scalar majorant records the worst packetwise nonlinear growth while discarding signs and cancellations.  It is conservative, but it composes cleanly.}
\end{figure}

The resulting analytic packet space is complete under its natural coefficient norm.  It is closed under composition whenever the inner majorant remains in the radius of the outer map.  This is the nonlinear analogue of a weighted Banach algebra.

\subsection{A local nonlinear Wiener theorem}

The linear Wiener principle asks whether structured localization survives inversion.  Nonlinearly, the natural first question is local.  Let
\[
F=I+N,
\qquad
N(0)=0.
\]
Assume that on the packet ball $B_R$,
\[
\kappa
=
\sup_{x\in B_R}\|DN(x)\|<1.
\]

\begin{theorem}[Near-identity nonlinear inverse]
For every $y$ satisfying
\[
\|y\|\le(1-\kappa)R,
\]
there is a unique $x\in B_R$ with $F(x)=y$.  The inverse map $H=F^{-1}$ is analytic on the resulting neighborhood and
\[
DH(y)
=
\left[I+DN(H(y))\right]^{-1}
=
\sum_{q=0}^{\infty}
[-DN(H(y))]^q.
\]
If $DN$ belongs to an inverse-closed weighted packet algebra, then $DH$ belongs to the same algebra.
\end{theorem}

\begin{proof}
For fixed $y$, define
\[
\Psi_y(x)=y-N(x).
\]
The derivative bound makes $\Psi_y$ a contraction with constant $\kappa$.  Moreover,
\[
\|\Psi_y(x)\|
\le
\|y\|+\|N(x)-N(0)\|
\le
(1-\kappa)R+\kappa R
=R,
\]
so $\Psi_y$ maps $B_R$ into itself.  Banach's fixed-point theorem gives a unique solution.  The analytic inverse-function theorem gives analyticity.  Differentiating $F(H(y))=y$ gives the inverse derivative formula, and the Neumann series converges because $\|DN\|<1$.
\end{proof}

\begin{center}
\small
\begin{tabular}{rrr}
\toprule
target & iterations & final residual\\
\midrule
0.15 & 12 & 8.02e-15 \\
0.30 & 18 & 3.94e-15 \\
0.45 & 24 & 2.94e-15 \\
0.60 & 30 & 4.77e-15 \\
\bottomrule
\end{tabular}
\end{center}

\begin{figure}[H]
\centering
\includegraphics[width=.82\linewidth]{figures/nonlinear/06_near_identity_inverse_convergence.png}
\caption{The fixed-point inverse remains reliable inside the contractive region.  Convergence slows as the target approaches the boundary where the derivative bound loses margin.}
\end{figure}

The theorem is deliberately local.  It does not claim that every globally invertible nonlinear map has an inverse with the same packet sparsity.  That stronger statement requires global hypotheses such as properness, coercivity, monotonicity, or a uniform lower bound on the derivative.

\subsection{Inverse jets}

For a univariate normalized map
\[
F(x)=x+\sum_{m\ge2}a_mx^m,
\]
write
\[
H(y)=F^{-1}(y)=y+\sum_{n\ge2}b_ny^n.
\]
Substituting $H$ into $F$ and setting every coefficient beyond the linear term to zero determines $b_n$ recursively.  In the multivariate packet setting the same recursion is expressed through rooted trees.

\begin{center}
\small
\begin{tabular}{rrr}
\toprule
inverse degree & maximum residual & residual at largest target\\
\midrule
2 & 1.07e-02 & 1.07e-02 \\
4 & 9.11e-04 & 9.11e-04 \\
8 & 1.44e-05 & 1.44e-05 \\
12 & 3.31e-07 & 3.31e-07 \\
18 & 1.52e-09 & 1.52e-09 \\
\bottomrule
\end{tabular}
\end{center}

\begin{figure}[H]
\centering
\includegraphics[width=.82\linewidth]{figures/nonlinear/07_inverse_series_accuracy.png}
\caption{Higher inverse jets enlarge the useful local range, but no finite jet becomes a global inverse automatically.}
\end{figure}

\subsection{Implicit functions and branch continuation}

Suppose
\[
\mathcal F(x,y)=0,
\]
and $D_y\mathcal F(x_0,y_0)$ is invertible.  The implicit-function theorem gives a local branch $y=h(x)$ with
\[
Dh(x_0)
=
-
\left[D_y\mathcal F(x_0,y_0)\right]^{-1}
D_x\mathcal F(x_0,y_0).
\]
Every factor can be resolved into packet blocks.  The same inverse-closure issue reappears: the branch derivative remains localized when the $y$-Jacobian inverse remains localized.

\begin{figure}[H]
\centering
\includegraphics[width=.82\linewidth]{figures/nonlinear/09_implicit_function_branches.png}
\caption{The two branches of $y^2+x-1=0$ are locally regular until they meet at the fold.  A local implicit theorem cannot cross the point where $D_y\mathcal F$ vanishes.}
\end{figure}

This example is elementary, but it states the general boundary plainly.  Packet continuation is reliable only while the relevant Jacobian packet remains invertible and the inverse remains controlled.

\subsection{Newton--Kantorovich certificates}

Let $F:X\to X$ be differentiable near $x_0$, let $A=DF(x_0)^{-1}$, and define
\[
\beta=\|AF(x_0)\|.
\]
If $ADF$ is Lipschitz with constant $L$ and
\[
h=\beta L\le\frac12,
\]
then the Newton iteration has a certified local convergence region.  One standard radius is
\[
r=
\frac{1-\sqrt{1-2h}}{L}.
\]

\begin{figure}[H]
\centering
\includegraphics[width=.80\linewidth]{figures/nonlinear/08_kantorovich_certified_region.png}
\caption{The scalar condition $\beta L\le1/2$ is conservative but explicit.  Packet estimates can replace one global $L$ by local block bounds, often improving the certificate.}
\end{figure}

The packet form is obtained by bounding the inverse Jacobian and derivative variation in weighted block norms.  This is especially useful when most cross-packet derivative blocks are absent or weak.

\subsection{Global inversion: what can and cannot be inferred}

A local inverse theorem does not prove global injectivity.  Even a derivative that is nonsingular at every sampled point does not prove it is nonsingular everywhere.  The finite diagnostic below tracks the smallest singular value on a grid; it is useful for locating a likely loss of invertibility, but it is not a theorem about points that were not sampled.

\begin{figure}[H]
\centering
\includegraphics[width=.82\linewidth]{figures/nonlinear/21_global_inverse_diagnostic.png}
\caption{A sampled minimum singular value can warn that a global fold is approaching.  It remains a diagnostic until analytic bounds control the entire domain.}
\end{figure}

A genuine global theorem needs additional assumptions.  In finite dimensions, one may use a Hadamard-type condition: everywhere-invertible derivative together with properness.  In monotone problems, coercivity or strong monotonicity can replace properness.  Resolution Theory helps organize these hypotheses packetwise, but does not remove them.

\subsection{Sparse nonlinear algorithms}

A dense quadratic map on $\mathbb R^n$ stores $n^3$ coefficients.  A hypergraph with $E$ active quadratic interactions stores $O(E)$ coefficients and evaluates in $O(E)$ scalar products.  The representation is exact when the omitted tensor blocks are truly zero.

\begin{algorithm}[H]
\caption{Sparse packet-hyperedge evaluation}
\begin{algorithmic}[1]
\State split every input into packet coordinates
\State initialize every output packet to zero
\For{each active hyperedge $(j_1,\ldots,j_m)\to i$}
    \State evaluate its multilinear block on the selected input packets
    \State accumulate the contribution in output packet $i$
\EndFor
\State reconstruct the output
\end{algorithmic}
\end{algorithm}

\begin{center}
\scriptsize
\begin{tabular}{rrrrrr}
\toprule
$n$ & active edges & sparse ms & dense ms & speedup & relative error\\
\midrule
32 & 160 & 0.0733 & 0.0621 & 0.85$\times$ & 8.13e-17 \\
48 & 240 & 0.1117 & 0.2012 & 1.80$\times$ & 1.20e-16 \\
64 & 320 & 0.1667 & 0.4476 & 2.68$\times$ & 1.51e-16 \\
80 & 400 & 0.2125 & 1.1684 & 5.50$\times$ & 1.24e-16 \\
96 & 480 & 0.2376 & 1.9899 & 8.38$\times$ & 1.22e-16 \\
128 & 640 & 0.3705 & 4.8722 & 13.15$\times$ & 8.79e-17 \\
160 & 800 & 0.4945 & 9.2726 & 18.75$\times$ & 1.18e-16 \\
\bottomrule
\end{tabular}
\end{center}

\begin{figure}[H]
\centering
\includegraphics[width=.82\linewidth]{figures/nonlinear/18_sparse_quadratic_speedup.png}
\caption{The sparse representation loses at the smallest size because Python dispatch dominates.  Once the ambient tensor becomes large enough, the removed cubic work becomes decisive.}
\end{figure}

These are CPU reference timings for a direct implementation.  They establish the arithmetic crossover in the tested environment, not a universal hardware law.

\subsection{Pruning and residual certification}

When weak hyperedges are discarded, the omitted blocks form a residual nonlinear map.  If $F=F_E+R_E$, then
\[
\|F(x)-F_E(x)\|
\le
\sum_{m\ge1}
\|R_E^{[m]}\|\,\|x\|^m.
\]
For a quadratic tensor with orthogonal coordinate blocks, the square root of the omitted block-energy sum is an exact Hilbert--Schmidt residual for the tensor itself.  State error is then bounded by multiplying the residual by the relevant packet norms of the inputs.

\begin{center}
\small
\begin{tabular}{rrr}
\toprule
retained hyperedges & actual state error & block residual\\
\midrule
5 & 85 & 114 \\
10 & 85.5 & 112 \\
20 & 84.9 & 109 \\
40 & 74.4 & 102 \\
70 & 57.6 & 90.9 \\
120 & 43.8 & 71.2 \\
216 & 5.65e-14 & 0 \\
\bottomrule
\end{tabular}
\end{center}

\begin{figure}[H]
\centering
\includegraphics[width=.82\linewidth]{figures/nonlinear/19_hyperedge_pruning.png}
\caption{A residual certificate is intentionally conservative.  Its purpose is not to predict the exact error perfectly, but to prevent a visually sparse hypergraph from being mistaken for an exact model.}
\end{figure}

\subsection{Approximate projectors and adaptive packet discovery}

A useful packet system need not be exactly invariant.  For a Jacobian $J=DF(x)$ and projector $P$, the instantaneous cross-packet leakage is
\[
\ell(P,J)=\|(I-P)JP\|.
\]
If this quantity is small on the region of interest, the packet is approximately invariant under the linearized nonlinear dynamics.

\begin{figure}[H]
\centering
\includegraphics[width=.82\linewidth]{figures/nonlinear/20_approximate_projector_leakage.png}
\caption{Packet quality depends on alignment with the Jacobian field.  Rotating away from a nearly invariant direction increases cross-packet leakage.}
\end{figure}

This suggests an adaptive strategy: choose or update projectors to minimize a weighted combination of leakage, hyperedge count, approximation residual, and refinement cost.  The optimization problem is meaningful, but a general globally optimal discovery algorithm remains open.

\subsection{Quadratic PDEs and residue selection}

For periodic Burgers dynamics,
\[
u_t+uu_x=\nu u_{xx},
\]
the Fourier coefficients satisfy
\[
\dot{\widehat u}_k
=
-\nu k^2\widehat u_k
-
\frac{ik}{2}
\sum_{p+q=k}\widehat u_p\widehat u_q.
\]
The nonlinear interaction is the triad hyperedge
\[
(p,q)\longrightarrow k,
\qquad p+q=k.
\]
Under a residue decomposition modulo $M$,
\[
p\equiv r,
\qquad q\equiv s
\]
forces
\[
k\equiv r+s\pmod M.
\]
Therefore
\[
\mathcal N_r\times\mathcal N_s
\longrightarrow
\mathcal N_{r+s\bmod M}.
\]
This is an exact nonlinear polyphase selection law.

\begin{figure}[H]
\centering
\includegraphics[width=.66\linewidth]{figures/nonlinear/10_burgers_polyphase_selection.png}
\caption{Every pair of input residue packets has one permitted output residue.  The rule is addition in the cyclic grading group.}
\end{figure}

\begin{center}
\small
\begin{tabular}{rr}
\toprule
mode count & packet/full relative error\\
\midrule
17 & 1.29e-16 \\
25 & 1.78e-16 \\
33 & 2.31e-16 \\
49 & 2.58e-16 \\
65 & 3.84e-16 \\
\bottomrule
\end{tabular}
\end{center}

\begin{figure}[H]
\centering
\includegraphics[width=.78\linewidth]{figures/nonlinear/11_burgers_packet_exactness.png}
\caption{The residue-resolved evaluator agrees with the full convolution to floating-point precision.  The packet law reorganizes the triads; it does not approximate them.}
\end{figure}

The same idea extends to higher polynomial PDEs, where the grading law adds all input residues.  Differential operators act linearly inside the same packet system, while nonlinearities generate the hyperedges.

\subsection{Reaction networks}

For a mass-action system with stoichiometric matrix $S$ and reaction orders $\alpha_{ir}$,
\[
\dot x
=
S\,v(x),
\qquad
v_r(x)=k_r\prod_i x_i^{\alpha_{ir}},
\]
each reaction is a nonlinear hyperedge from its reactant multiset to the species whose concentrations it changes.

\begin{figure}[H]
\centering
\includegraphics[width=.82\linewidth]{figures/nonlinear/12_reaction_hypergraph.png}
\caption{A bimolecular reaction is literally a two-input hyperedge.  Reversible reactions appear as oppositely directed interaction patterns.}
\end{figure}

Conservation laws are left-null vectors of $S$.  A packet decomposition aligned with conserved moieties separates invariant combinations from reactive degrees of freedom.  This is a useful example of representation choice changing the apparent complexity without changing the chemistry.

\begin{figure}[H]
\centering
\includegraphics[width=.82\linewidth]{figures/nonlinear/13_reaction_dynamics.png}
\caption{The trajectories are ordinary mass-action solutions.  The packet language records which nonlinear reaction channels produce each transfer.}
\end{figure}

\subsection{Volterra systems and nonlinear memory}

A second-order causal Volterra system has the form
\[
y(t)
=
\int_0^t h_1(\tau)x(t-\tau)\,d\tau
+
\int_0^t\int_0^t
h_2(\tau_1,\tau_2)
 x(t-\tau_1)x(t-\tau_2)
\,d\tau_1d\tau_2.
\]
The first kernel is a linear transfer through time.  The second kernel is a continuum of two-input memory hyperedges.  Temporal packets may be lag windows, exponential modes, wavelets, or fractional-memory components.

\begin{figure}[H]
\centering
\includegraphics[width=.84\linewidth]{figures/nonlinear/14_volterra_memory_response.png}
\caption{The output depends on the present signal and on nonlinear interactions among recent lag packets.}
\end{figure}

A sum-of-exponentials approximation to the kernels converts long memory into a finite augmented state.  The resulting system can then be treated with the same state-space and Kalman machinery developed in the probability and memory appendix.

\subsection{Bayesian normalization as a nonlinear transfer}

Given prior packet weights $w_i$ and likelihoods $L_i$, Bayes' rule is
\[
p_i
=
\frac{w_iL_i}{\sum_jw_jL_j}.
\]
In log coordinates it is a softmax:
\[
p=\operatorname{softmax}(\log w+\log L).
\]
Its Jacobian is
\[
D\operatorname{softmax}(z)
=
\operatorname{diag}(p)-pp^T.
\]
Every packet competes with every other through the normalization denominator.  The map is nonlinear even when the unnormalized evidence multiplication is packetwise.

\begin{figure}[H]
\centering
\includegraphics[width=.82\linewidth]{figures/nonlinear/15_softmax_sensitivity.png}
\caption{The normalization map is least sensitive when one packet already dominates and most sensitive in the intermediate competition regime.}
\end{figure}

This makes clear why Bayesian packetization requires both local evidence and a global normalization channel.  Dropping the latter changes the theory.

\subsection{Nonlinear feedback and control}

For
\[
\dot x=f(x,u),
\qquad
u=\kappa(x),
\]
the closed-loop field is
\[
F_{\mathrm{cl}}(x)=f(x,\kappa(x)).
\]
Its jets are obtained by the same packet Fa\`a di Bruno law.  A controller can therefore create, cancel, or redirect nonlinear hyperedges.  A feedback design that appears local in one coordinate system may become dense in another.

A practical packet-control objective is
\[
\min_{\kappa,\mathcal R}
\left[
\text{closed-loop residual}
+
\lambda_1\text{ active hyperedges}
+
\lambda_2\text{ control effort}
+
\lambda_3\text{ packet leakage}
\right].
\]
The formula does not solve nonlinear control by itself.  It states the structural tradeoff that a sparse packet controller must manage.

\subsection{Normal forms and resonances}

Near an equilibrium, nonresonant polynomial terms can often be removed by a near-identity change of coordinates.  For frequencies $\omega_i$, a quadratic interaction $(i,j)\to k$ has detuning
\[
\Delta_{ijk}=\omega_i+\omega_j-\omega_k.
\]
Small $|\Delta_{ijk}|$ marks a resonant or nearly resonant hyperedge.  Large detuning produces rapid phase cancellation and is a candidate for normal-form elimination.

\begin{figure}[H]
\centering
\includegraphics[width=.82\linewidth]{figures/nonlinear/16_resonance_detuning.png}
\caption{Ordering triads by detuning exposes the small set that a normal-form calculation must retain carefully.}
\end{figure}

The homological equation that removes a nonresonant term divides by its detuning.  This is where small divisors enter.  Packet sparsity does not remove the small-divisor problem, but it can isolate exactly which interaction families carry it.

\subsection{Hamiltonian and symplectic packets}

Let $J$ be a symplectic matrix and let
\[
\dot z=J\nabla H(z).
\]
If the packet system is compatible with the symplectic splitting, Hamiltonian derivatives become structured hyperedges and Poisson brackets preserve the Hamiltonian algebra:
\[
\{F,G\}=\nabla F^TJ\nabla G.
\]
The bracket is bilinear, antisymmetric, satisfies the Leibniz rule, and obeys the Jacobi identity.

\begin{figure}[H]
\centering
\includegraphics[width=.66\linewidth]{figures/nonlinear/17_hamiltonian_phase_portrait.png}
\caption{The implicit midpoint rule preserves the quadratic oscillator energy to numerical precision, illustrating why the discretization must respect the native geometric structure.}
\end{figure}

A generic sparse truncation of Hamiltonian hyperedges need not remain Hamiltonian.  Structure-preserving pruning must retain terms in combinations that preserve the generating Hamiltonian or an explicitly controlled modified Hamiltonian.

\subsection{Semilinear evolution and Duhamel trees}

For
\[
\dot u=Lu+N(u),
\]
the variation-of-constants formula is
\[
u(t)
=
e^{tL}u_0
+
\int_0^te^{(t-s)L}N(u(s))\,ds.
\]
Iterating this identity produces Duhamel trees.  Linear edges propagate packets through $e^{tL}$; nonlinear vertices join them through derivatives of $N$.

\begin{figure}[H]
\centering
\includegraphics[width=.68\linewidth]{figures/nonlinear/22_semilinear_packet_flow.png}
\caption{A damped rotating linear flow is bent by a quadratic packet interaction.  The trajectory combines ordinary transfer paths with nonlinear vertices.}
\end{figure}

If a pruned nonlinearity $N_E$ satisfies
\[
\sup_{\|u\|\le R}\|N(u)-N_E(u)\|\le\varepsilon
\]
and $N$ is Lipschitz with constant $L_N$, a Gr\"onwall argument gives
\[
\|u(t)-u_E(t)\|
\le
M\varepsilon
\int_0^te^{(\omega+ML_N)(t-s)}\,ds
\]
whenever $\|e^{tL}\|\le Me^{\omega t}$.  This is the flow-level certificate corresponding to hyperedge pruning.

\subsection{Lie brackets and the order of nonlinear operations}

Two nonlinear flows need not commute.  Let $f$ and $g$ be vector fields with local flows $\Phi_t^f$ and $\Phi_s^g$.  The first obstruction to exchanging their order is the Lie bracket
\[
[f,g](x)
=
Dg(x)f(x)-Df(x)g(x).
\]
In packet form, $Df$ and $Dg$ are ordinary transfer graphs, while multiplication by the field values selects which edges are active at the current state.  The bracket is therefore a signed sum of two composed packet paths.

\begin{proposition}[Packet closure of Lie brackets]
If $f$ and $g$ are packet-analytic vector fields on a common packet ball and their Jacobians belong to the weighted transfer algebra there, then $[f,g]$ is packet analytic on the same ball.  Its majorant is bounded by
\[
\mathfrak m_{[f,g]}(t)
\le
\mathfrak m_{Dg}(t)\,\mathfrak m_f(t)
+
\mathfrak m_{Df}(t)\,\mathfrak m_g(t).
\]
\end{proposition}

\begin{proof}
The bracket is the difference of the two compositions $Dg\cdot f$ and $Df\cdot g$.  Analytic packet maps are closed under composition and multiplication by bounded linear packet maps.  Applying the triangle inequality to the two terms gives the stated majorant.
\end{proof}

\begin{figure}[H]
\centering
\includegraphics[width=.68\linewidth]{figures/nonlinear/23_lie_bracket_field.png}
\caption{The Lie bracket is itself a vector field.  It measures the leading discrepancy between following one nonlinear flow and then the other in the opposite order.}
\end{figure}

This matters whenever algorithms alternate nonlinear substeps.  If the bracket is small in the active packet region, operator splitting is accurate to higher order than a crude global estimate suggests.  If the bracket is concentrated in a few packets, one can refine only those packets.

\subsection{Stochastic nonlinear transfer and the It\^o correction}

For the stochastic differential equation
\[
dX_t=b(X_t)\,dt+\sigma(X_t)\,dW_t,
\]
a smooth observable $\varphi$ evolves according to the generator
\[
\mathcal L\varphi
=
\nabla\varphi\cdot b
+
\frac12\operatorname{tr}
\left(
\sigma\sigma^T D^2\varphi
\right).
\]
The second term is not optional.  It is the contribution of quadratic variation.  In the nonlinear packet calculus it is a contraction of a second-order observable jet with two diffusion transfer edges.

\begin{theorem}[Packet form of the It\^o generator]
Let $P_j$ resolve the state space and let $Q_i$ resolve the observable output.  Then
\[
Q_i\mathcal L\varphi
=
\sum_j Q_iD\varphi[P_jb]
+
\frac12
\sum_{j,k}
Q_iD^2\varphi
[P_j\sigma, P_k\sigma],
\]
with the diffusion indices contracted over the driving noise coordinates.
\end{theorem}

\begin{proof}
Insert the state resolution into the drift vector and into both input slots of the Hessian term.  The It\^o formula is linear in the drift and bilinear in the diffusion columns, so the packet expansion follows exactly by distributivity.
\end{proof}

\begin{figure}[H]
\centering
\includegraphics[width=.82\linewidth]{figures/nonlinear/24_ito_quadratic_variation.png}
\caption{For the observable $x^2$, the stochastic generator differs from the ordinary drift derivative by a constant second-order correction.  Omitting that correction changes the dynamics.}
\end{figure}

This is the precise point at which stochastic calculus enters the nonlinear operad.  Drift uses first-order edges; quadratic variation uses second-order hyperedges.  Higher-order stochastic expansions produce still larger interaction trees.

\subsection{Interaction-tree growth in Duhamel and Picard expansions}

For a quadratic nonlinearity, each additional Picard substitution replaces one leaf by a binary vertex.  The number of binary rooted trees at order $n$ is the Catalan number
\[
C_n=
\frac{1}{n+1}\binom{2n}{n}.
\]
The growth is already exponential up to a polynomial factor:
\[
C_n\sim\frac{4^n}{\sqrt{\pi}n^{3/2}}.
\]

\begin{figure}[H]
\centering
\includegraphics[width=.78\linewidth]{figures/nonlinear/25_duhamel_tree_growth.png}
\caption{The combinatorial growth of interaction trees explains why nonlinear perturbation theory becomes difficult long before the state dimension itself becomes large.}
\end{figure}

Packet selection can reduce this growth because many labelings of a tree may violate residue, symmetry, conservation, or locality rules.  The tree shape remains, but the number of admissible packet decorations can collapse dramatically.  This is one of the strongest reasons to combine operadic composition with selection laws.

\subsection{Block Newton methods and nonlinear preconditioning}

Full Newton iteration solves
\[
DF(x_n)\,\delta_n=F(x_n),
\qquad
x_{n+1}=x_n-\delta_n.
\]
When the Jacobian is nearly block diagonal in a packet system, a cheaper step uses only the diagonal packet blocks:
\[
\left(P_iDF(x_n)P_i\right)\delta_{n,i}
=
P_iF(x_n).
\]
The method is a nonlinear block-Jacobi preconditioner.

\begin{theorem}[Local convergence of block Newton under weak coupling]
Write
\[
DF(x)=D(x)+R(x),
\]
where $D(x)$ is the packet-diagonal part.  Suppose $D(x)$ is invertible on a neighborhood and
\[
\sup_x\|D(x)^{-1}R(x)\|\le\eta<1.
\]
If the usual Newton regularity assumptions hold, then the block-Newton correction is a contraction of the exact Newton error up to higher-order terms, with linear factor at most $\eta+o(1)$ near the solution.
\end{theorem}

\begin{proof}
Near a solution $x_*$, write the error $e_n=x_n-x_*$.  The block step gives
\[
e_{n+1}
=
e_n-D(x_n)^{-1}F(x_n).
\]
Taylor expansion yields
\[
F(x_n)=DF(x_*)e_n+O(\|e_n\|^2).
\]
Replacing $DF(x_*)$ by $D(x_*)+R(x_*)$ gives
\[
e_{n+1}
=
-D(x_*)^{-1}R(x_*)e_n+O(\|e_n\|^2),
\]
so the leading contraction factor is bounded by $\eta$.
\end{proof}

\begin{figure}[H]
\centering
\includegraphics[width=.82\linewidth]{figures/nonlinear/26_block_newton_convergence.png}
\caption{As cross-packet coupling grows, the block approximation to the Jacobian becomes less accurate and the iteration slows.  The method remains useful while the neglected coupling is a controlled perturbation.}
\end{figure}

This theorem gives a practical representation criterion: choose packets that make the Jacobian diagonal blocks easy to invert and the off-diagonal remainder small.

\subsection{Constrained systems and multiplier packets}

Many nonlinear systems are not free to move in the whole state space.  Equality constraints
\[
c(x)=0
\]
introduce Lagrange multipliers and a saddle-point Newton system
\[
\begin{pmatrix}
D_xr & Dc(x)^T\\
Dc(x) & 0
\end{pmatrix}
\begin{pmatrix}
\delta x\\
\delta\lambda
\end{pmatrix}
=
-
\begin{pmatrix}
r(x,\lambda)\\
c(x)
\end{pmatrix}.
\]
The multiplier variables form their own packet family.  Their transfer blocks encode which constraints couple which state packets.

\begin{figure}[H]
\centering
\includegraphics[width=.76\linewidth]{figures/nonlinear/28_constraint_newton.png}
\caption{A single linear constraint couples two state coordinates through one multiplier packet.  Larger constrained systems inherit the same saddle-point block structure.}
\end{figure}

Differential-algebraic equations fit the same pattern.  The dynamic packets carry time derivatives, while algebraic packets enforce compatibility.  One must preserve the constraint index and rank conditions; block notation alone does not guarantee a well-posed DAE.

\subsection{Invariant manifolds and graph equations}

Suppose a system splits into slow and fast packets,
\[
\dot x=Ax+f(x,y),
\qquad
\dot y=By+g(x,y).
\]
An invariant manifold written as $y=h(x)$ satisfies
\[
Dh(x)\left(Ax+f(x,h(x))\right)
=
Bh(x)+g(x,h(x)).
\]
This is a nonlinear packet equation for the graph map $h$.  Expanding $h$ into jets gives a recursive family of homological equations.  The linear operator to invert at each order contains combinations of the spectra of $A$ and $B$; resonances appear when those combinations approach zero.

The graph-transform method can also be written packetwise.  A spectral gap between the slow and fast packets makes the transform contractive.  Without that gap, the proposed manifold may fail to exist or fail to be smooth.

\subsection{Hyperedge entropy and the cost of resolution}

Given squared block energies $e_\gamma=\|A_\gamma\|^2$, normalize
\[
p_\gamma
=
\frac{e_\gamma}{\sum_\eta e_\eta}
\]
and define the hyperedge entropy
\[
H_{\mathrm{edge}}
=
-\sum_\gamma p_\gamma\log p_\gamma.
\]
Low entropy means that most nonlinear energy is concentrated in a small set of hyperedges.  High entropy means that the interaction is spread broadly and is less compressible by simple pruning.

\begin{figure}[H]
\centering
\includegraphics[width=.78\linewidth]{figures/nonlinear/27_hyperedge_entropy_refinement.png}
\caption{A generic dense tensor develops more distinguishable interaction channels under refinement.  Refinement becomes useful only when the newly exposed distribution is sufficiently uneven or structured.}
\end{figure}

The quantity is resolution-dependent.  That is a feature rather than a defect: it measures how complicated the same nonlinear map appears in the chosen representation.  Comparing entropies across admissible resolutions suggests a concrete packet-discovery objective.

\subsection{Countable packet systems}

The finite theory extends to countably many packets only with convergence hypotheses.  Let
\[
I=\sum_{j\ge0}P_j
\]
converge strongly.  For an $m$-linear map, require absolute summability of the weighted block norms:
\[
\sup_i\frac1{b_i}
\sum_{j_1,\ldots,j_m}
\|A_{i;j_1,\ldots,j_m}\|
\prod_ra_{j_r}
<\infty.
\]

\begin{theorem}[Countable multilinear reconstruction]
Under the summability condition above, the packet reconstruction converges absolutely in each output packet and defines a bounded multilinear map on the weighted packet space.  Finite truncation after packet index $N$ has error bounded by the corresponding tail of the block-norm sum.
\end{theorem}

\begin{proof}
The weighted block estimate makes the series in each output packet absolutely convergent.  Completeness of the Banach space gives a limit.  The same estimate applied to the omitted indices gives the truncation bound.  Summing the output packets through the strong resolution reconstructs the map on the declared packet domain.
\end{proof}

This is the correct infinite-dimensional statement.  Merely writing an infinite hypergraph does not establish convergence.

\subsection{Nonlinear packet atlases}

One packet system may be well adapted near one part of state space and poor elsewhere.  Let $\{U_a\}$ cover the state domain and let each chart carry projectors $P_j^{(a)}$.  On an overlap, a coordinate change
\[
\psi_{ba}:U_a\to U_b
\]
transports jets by the nonlinear chain rule.  The transition itself has a packet hypergraph, so changing resolution is a nonlinear operation rather than a relabeling.

\begin{proposition}[Atlas compatibility]
Suppose the chart transitions satisfy the cocycle law
\[
\psi_{ca}=\psi_{cb}\circ\psi_{ba}
\]
on triple overlaps.  Then the corresponding packet jets satisfy the same cocycle law under operadic substitution.
\end{proposition}

\begin{proof}
Take derivatives of the transition identity at every order.  Fa\`a di Bruno composition turns the right-hand side into rooted-tree substitution of the two transition jets.  Since the maps agree, their jets agree.  Inserting the local projectors gives the packet version.
\end{proof}

This is the nonlinear counterpart of the projector atlas developed in the core extensions.  It allows the theory to remain local when no single global packetization is stable.


\subsection{Neural layers and learned packet systems}

A feed-forward layer
\[
x\longmapsto W_2\sigma(W_1x+b_1)+b_2
\]
is a nonlinear transfer map.  Its Jacobian is
\[
W_2\operatorname{diag}(\sigma'(W_1x+b_1))W_1,
\]
and higher derivatives insert higher derivatives of $\sigma$ between the two linear maps.  A packet system may be imposed from architecture, learned from data, or adapted from Jacobian leakage.

The useful question is not whether a neural layer can be written in packet notation; it can.  The useful question is whether one can discover packets for which the active derivative hypergraph is smaller, more stable, or easier to audit than the original dense representation.

\subsection{Computational interpretation}

The nonlinear calculus suggests a compiler pipeline:
\[
\begin{aligned}
\text{map}
&\longrightarrow
\text{jet extraction}
\longrightarrow
\text{packet hypergraph}\\
&\longrightarrow
\text{pruning or low-rank compression}
\longrightarrow
\text{grouped tensor contractions}.
\end{aligned}
\]
A practical implementation must distinguish four costs:
\[
T_{\mathrm{total}}
=
T_{\mathrm{discover}}
+
T_{\mathrm{compile}}
+
R\,T_{\mathrm{apply}}
+
T_{\mathrm{materialize}}.
\]
The same representation may be poor for one use and excellent when reused many times.

GPU realization would group hyperedges by arity, packet dimensions, dtype, and output packet.  Equal-size groups map naturally to batched tensor contractions; irregular groups require grouped kernels or a persistent scheduler.  These are engineering consequences of the mathematical hypergraph, not additional theorems.

\subsection{Worked example: a two-packet quadratic map}

Consider
\[
F(x_1,x_2)
=
\begin{pmatrix}
x_1+\alpha x_1x_2\\
x_2+\beta x_1^2
\end{pmatrix},
\]
with coordinate packets
\[
P_1(x_1,x_2)=(x_1,0),
\qquad
P_2(x_1,x_2)=(0,x_2).
\]
The linear jet has only two diagonal edges:
\[
1\to1,
\qquad
2\to2.
\]
The quadratic jet has two nonzero hyperedges:
\[
(1,2)\to1
\]
with coefficient $\alpha$, and
\[
(1,1)\to2
\]
with coefficient $2\beta$ in derivative normalization.  The factor of two is not cosmetic: the Taylor term is
\[
\frac12D^2F(0)[x,x].
\]
This small example shows why coefficient tensors and derivative tensors must not be mixed without a normalization convention.

If the packets are merged, the map looks like one dense quadratic tensor.  If they are kept separate, the nonlinear graph contains only two active hyperedges.  The exact arithmetic reduction is immediate: evaluating the two monomials replaces a general two-output quadratic tensor contraction.

Now compose $F$ with
\[
G(y_1,y_2)
=
\begin{pmatrix}
y_1^2+y_2\\
y_1-y_2
\end{pmatrix}.
\]
The first output of $G\circ F$ contains both inherited inner hyperedges and new rooted trees obtained by feeding the linear outputs of $F$ into the quadratic vertex of $G$.  This is precisely the second-order Fa\`a di Bruno decomposition.  The accompanying polynomial-composition test verifies the equality coefficient by coefficient.

\subsection{Worked example: inverse jets near the identity}

Take
\[
F(x)=x+ax^2,
\qquad a>0.
\]
The exact inverse branch near zero is
\[
H(y)
=
\frac{-1+\sqrt{1+4ay}}{2a}.
\]
Its series begins
\[
H(y)
=
y-ay^2+2a^2y^3-5a^3y^4+14a^4y^5-\cdots.
\]
The coefficients are signed Catalan numbers.  The inverse-jet recursion in the Python implementation reproduces them without using the closed form.

For $a=1/4$, the nearest branch point occurs at
\[
y=-1.
\]
The Taylor series therefore has radius one around the origin.  The finite-degree experiment stays well inside that disk and shows monotone residual improvement as the inverse degree increases.  This example makes the locality of the nonlinear Wiener theorem visible: the inverse exists on a much larger real interval, but the Taylor expansion around one point is controlled by the nearest complex singularity.

\subsection{Worked example: block Newton with weak coupling}

Consider four equations grouped into two packets:
\[
\begin{aligned}
F_1(x)&=x_1^2-1+\varepsilon x_3,\\
F_2(x)&=x_2^2-4+\varepsilon x_4,\\
F_3(x)&=x_3^2-9+\varepsilon x_1,\\
F_4(x)&=x_4^2-16+\varepsilon x_2.
\end{aligned}
\]
The packet-diagonal Jacobian contains the four derivative terms $2x_i$.  The off-diagonal packet coupling is proportional to $\varepsilon$.  The block-Newton method ignores that coupling when solving each correction, then lets later iterations repair the residual.

At $\varepsilon=0$, the two packets are independent.  As $\varepsilon$ grows, the block step remains inexpensive but requires more iterations.  The experiment records exactly that transition.  The method is not superior because it is ``more block-like.''  It is superior only while the saved solve cost exceeds the extra iterations caused by the omitted coupling.

\subsection{Worked example: nonlinear Bayesian competition}

Let three hypotheses have prior weights
\[
w=(0.2,0.3,0.5)
\]
and likelihoods
\[
L=(0.5,1.5,0.8).
\]
The unnormalized evidence is
\[
w\odot L=(0.1,0.45,0.4),
\]
so the posterior is
\[
p
=
\frac{(0.1,0.45,0.4)}{0.95}
\approx
(0.1053,0.4737,0.4211).
\]
The multiplication is packetwise, but normalization couples every packet to every other packet through the denominator.  The Jacobian
\[
\operatorname{diag}(p)-pp^T
\]
has row and column sums zero, expressing conservation of total probability.  When one posterior packet approaches probability one, the Jacobian collapses and the update becomes locally insensitive.  When several packets compete, the map is more responsive.

This is a useful warning for approximate inference: pruning a weak hypothesis before normalization and pruning it after normalization are different operations.

\subsection{Worked example: a quadratic Volterra memory}

Let the memory window have two lags and write
\[
y_t
=
h_0x_t+h_1x_{t-1}
+
q_{00}x_t^2
+2q_{01}x_tx_{t-1}
+q_{11}x_{t-1}^2.
\]
There are two linear lag edges and three symmetric quadratic memory hyperedges.  The mixed term $(t,t-1)\to t$ is counted once in the symmetric representation and twice in the expanded polynomial.

If the kernel is long but approximately separable,
\[
h_2(\tau_1,\tau_2)
\approx
\sum_{r=1}^R
u_r(\tau_1)\nu_r(\tau_2),
\]
then the quadratic memory can be evaluated through $R$ filtered coordinates rather than a full two-dimensional lag tensor.  This is the nonlinear memory analogue of low-rank matrix compression.

\subsection{Worked example: an equality-constrained nonlinear solve}

The quadratic problem
\[
\min_x\frac12\|x\|^2
\qquad
\text{subject to}
\qquad
x_1+x_2=1
\]
has KKT equations
\[
\begin{pmatrix}
x_1+\lambda\\
x_2+\lambda\\
x_1+x_2-1
\end{pmatrix}
=0.
\]
Treating $x$ and $\lambda$ as separate packet types produces the saddle matrix
\[
\begin{pmatrix}
I&A^T\\
A&0
\end{pmatrix},
\qquad
A=(1,1).
\]
One Newton step from the origin gives
\[
x_1=x_2=\frac12,
\qquad
\lambda=-\frac12.
\]
The example is linear, but the same multiplier packet structure persists when the objective and constraints are nonlinear.

\subsection{Theorem dependency ledger}

The nonlinear theory has a clear dependency order.  The table below records the main logical chain so that later applications are not mistaken for independent assumptions.

\begin{longtable}{p{.24\linewidth}p{.28\linewidth}p{.38\linewidth}}
\toprule
Result & Depends on & What it supplies\\
\midrule
\endfirsthead
\toprule
Result & Depends on & What it supplies\\
\midrule
\endhead
Multilinear reconstruction & finite projector resolutions & exact hyperedge decomposition\\
Weighted block estimate & reconstruction and packet weights & norm control for one nonlinear order\\
Fa\`a di Bruno packets & smooth composition and multilinearity & rooted-tree composition law\\
Operad structure & Fa\`a di Bruno and symmetry actions & abstract nonlinear closure\\
Refinement coherence & complete child resolutions & consistency across scales\\
Analytic majorant theorem & weighted multilinear bounds & convergence under nonlinear composition\\
Near-identity inverse & derivative contraction & local existence, uniqueness, and inverse analyticity\\
Inverse-jet recursion & formal composition identity & finite Taylor approximations of the inverse\\
Implicit-function theorem & invertible partial Jacobian & local parameterized solution branches\\
Newton--Kantorovich & inverse derivative and Lipschitz control & explicit convergence radius\\
Block Newton theorem & invertible diagonal blocks and weak coupling & cheaper approximate Newton corrections\\
Residual pruning theorem & omitted block norm bounds & certified sparse approximation\\
Semilinear flow bound & semigroup estimate and Lipschitz nonlinearity & time-dependent trajectory error\\
Countable reconstruction & absolute weighted summability & infinite packet extension\\
Atlas compatibility & smooth transition cocycles & local nonlinear representations without one global chart\\
\bottomrule
\end{longtable}


\subsection{What has been proved}

The appendix establishes the following finite-dimensional and Banach-space results under the hypotheses stated in each theorem:
\begin{enumerate}[label=(\arabic*)]
\item exact multilinear packet reconstruction;
\item weighted packet norm control;
\item projector-resolved jets;
\item the packet Fa\`a di Bruno composition law;
\item the colored symmetric operad structure;
\item exact recursive refinement and coarsening coherence;
\item the no-free-refinement principle;
\item closure of weighted analytic packet maps under composition;
\item a local nonlinear Wiener inverse theorem for near-identity maps;
\item inverse-jet recursion;
\item packet forms of the implicit-function and Newton--Kantorovich mechanisms;
\item exact residue selection for quadratic Fourier interactions;
\item hypergraph representations of mass-action and Volterra systems;
\item residual-certified nonlinear pruning;
\item leakage diagnostics for approximate projectors;
\item Hamiltonian and semilinear-flow compatibility statements.
\end{enumerate}

\subsection{What remains open}

The strongest unresolved directions are now easier to state.

\paragraph{Global nonlinear Wiener theory.}
The local theorem uses a derivative contraction on a ball.  A global inverse-closed class needs hypotheses that prevent folding and escape at infinity.

\paragraph{Infinite packet systems.}
Countably infinite hypergraphs require unconditional convergence, topology choices, and careful domain control.

\paragraph{Sharp inverse-jet decay.}
Local membership follows under smallness.  Optimal decay exponents and constants remain open.

\paragraph{Singular nonlinear PDEs.}
When products are not classically defined, renormalization, paracontrolled calculus, or regularity-structure machinery must be added before the hypergraph is analytic rather than formal.

\paragraph{Automatic packet discovery.}
A general algorithm must balance leakage, hyperedge count, residual, refinement cost, and stability.

\paragraph{Structure-preserving hardware.}
A fused nonlinear packet kernel must be compared fairly with optimized dense and tensor-library baselines on identical maps and output contracts.

\subsection{Closing perspective}

The linear calculus closes transfer paths.  The nonlinear extension closes interaction trees.  Its basic object is the hierarchy
\[
Q_iD^mF(x)
\left[P_{j_1}\cdot,\ldots,P_{j_m}\cdot\right].
\]
These blocks reconstruct the local nonlinear map, refine without ambiguity, compose through rooted trees, and admit analytic control through scalar majorants.  Near the identity, the inverse is packet analytic and its derivative is a convergent sum of transfer paths.

Within Resolution Theory, the concise description is:
\[
\boxed{
\text{a recursively refinable colored operad of analytic packet maps,}
}
\]
\[
\boxed{
\text{equipped with local inverse, implicit, flow, and residual certificates.}
}
\]



\subsection{Transition to Resolution Logic}

The nonlinear operad tells us what the interaction tensors are and how they compose.  It does not by itself say what may be concluded when the inputs, charts, or retained hyperedges are only approximate.  That is a separate problem.  Resolution Logic adds judgments that carry error, cost, failure probability, and evidence through the same refinement and substitution diagrams.

\section{Resolution Logic: truth, error, and representation in one judgment}
\label{sec:resolution-logic}

\subsection{Why a separate logic is needed}

The first version of Resolution Theory already carried a small error calculus.  A vector was represented at a named resolution, an approximation came with a norm bound, and linear or bilinear maps transported that bound forward.  That core was useful, but it left several ordinary mathematical situations outside the formal language.  It did not say how a claim changes when one packet is opened into finer packets.  It did not distinguish a theorem from a conditional theorem or a finite computation.  It had no quantifiers for statements that become visible only after refinement.  Most importantly, it did not yet explain how the nonlinear packet calculus should carry proof obligations alongside values.

Resolution Logic begins from a plain observation.  In serious work, a mathematical statement is rarely just true or false in isolation.  It is true under hypotheses, in a particular representation, to a stated tolerance, at some computational or verification cost, and with a visible chain of evidence.  None of these qualifiers weakens mathematics.  Hiding them does.

A typical quantitative judgment will be written
\[
   \Gamma\vdash_{\mathcal R}
   \widehat x : x
   [\varepsilon,c,\eta;\mathcal E].
\]
It means:
\begin{itemize}
\item the assumptions in $\Gamma$ are in force;
\item $\widehat x$ is the available representative of the intended object $x$;
\item the representative is expressed in resolution $\mathcal R$;
\item the deterministic error is at most $\varepsilon$;
\item the declared cost is $c$;
\item the certificate may fail with probability at most $\eta$;
\item $\mathcal E$ records proof steps, assumptions, computations, and conjectural inputs without collapsing them into one status word.
\end{itemize}

The exact finite logic developed below has two layers.  The first is Boolean and asks which packets satisfy a predicate.  The second is quantitative and carries norm bounds, costs, failure probabilities, and evidence through linear and nonlinear maps.  The two layers meet when a numerical certificate decides, fails to decide, or refutes a predicate at a given resolution.

\begin{figure}[H]
\centering
\includegraphics[width=.90\linewidth]{figures/logic/01_logic_architecture.png}
\caption{Resolution Logic keeps the represented object, the chosen resolution, and the certificate in the same mathematical diagram.  Changing the resolution is therefore an inference step with its own residual and cost.}
\end{figure}

\subsection{The finite semantics of refinement}

Let $I_f$ be a finite set of fine packets and $I_c$ a finite set of coarse packets.  A refinement map is a surjection
\[
   r:I_f\longrightarrow I_c.
\]
The fiber $r^{-1}(j)$ lists the fine packets contained in coarse packet $j$.  Predicates are subsets of the relevant packet set.  Equivalently, they are Boolean-valued functions.

Three operations are forced on us by the map $r$.

\begin{definition}[Pullback, existential visibility, and universal persistence]
For $B\subseteq I_c$ and $A\subseteq I_f$, define
\[
\begin{aligned}
   r^*B
   &=\{i\in I_f:r(i)\in B\},\\
   \exists_r A
   &=\{j\in I_c:\text{some }i\in r^{-1}(j)\text{ lies in }A\},\\
   \forall_r A
   &=\{j\in I_c:r^{-1}(j)\subseteq A\}.
\end{aligned}
\]
\end{definition}

The meanings are concrete.  Pullback asks what a coarse statement says about every fine packet beneath it.  Existential visibility says that a coarse packet contains at least one fine witness.  Universal persistence says that every fine packet inside the coarse packet satisfies the statement.

\begin{theorem}[Refinement Galois connections]
For every $A\subseteq I_f$ and $B\subseteq I_c$,
\[
   \exists_r A\subseteq B
   \quad\Longleftrightarrow\quad
   A\subseteq r^*B,
\]
and
\[
   r^*B\subseteq A
   \quad\Longleftrightarrow\quad
   B\subseteq\forall_r A.
\]
Thus
\[
   \exists_r\dashv r^*\dashv\forall_r
\]
as monotone maps between the finite Boolean algebras of predicates.
\end{theorem}

\begin{proof}
The first equivalence follows packet by packet.  Suppose $\exists_r A\subseteq B$ and let $i\in A$.  Then $r(i)\in\exists_r A$, hence $r(i)\in B$, so $i\in r^*B$.  Conversely, if $A\subseteq r^*B$ and $j\in\exists_rA$, choose $i\in A$ with $r(i)=j$.  Since $i\in r^*B$, one has $j\in B$.

For the second equivalence, assume $r^*B\subseteq A$ and let $j\in B$.  Every $i\in r^{-1}(j)$ belongs to $r^*B$ and therefore to $A$, so $j\in\forall_rA$.  Conversely, if $B\subseteq\forall_rA$ and $i\in r^*B$, then $r(i)\in B\subseteq\forall_rA$, which says precisely that every point in the fiber of $r(i)$, including $i$, lies in $A$.
\end{proof}

This theorem supplies the first genuinely logical structure of refinement.  It is not an analogy with quantifiers.  On finite packet sets, the adjunctions are exact.

\begin{figure}[H]
\centering
\includegraphics[width=.88\linewidth]{figures/logic/10_refinement_quantifiers.png}
\caption{A refinement map carries three natural logical operations.  Pullback repeats a coarse predicate on every finer packet, while existential and universal pushforwards distinguish local visibility from uniform persistence.}
\end{figure}

\subsection{Completeness of the finite Boolean fiber}

At one finite resolution, every predicate is simply a union of packet atoms.  This elementary fact gives a complete logic for the Boolean fiber.

\begin{theorem}[Atomic completeness at a finite resolution]
Let $I$ be a finite packet set.  Every predicate $A\subseteq I$ is the disjunction of the atoms it contains:
\[
   A=\bigvee_{i\in A}\{i\}.
\]
For predicates $A_1,\ldots,A_m,B\subseteq I$,
\[
   A_1,\ldots,A_m\models B
\]
if and only if every atom contained in all $A_k$ is contained in $B$.
\end{theorem}

\begin{proof}
The first statement is the definition of a finite subset as the union of its singleton elements.  For entailment, the conjunction of the premises is the intersection $\bigcap_kA_k$.  It is contained in $B$ exactly when every singleton atom in that intersection belongs to $B$.
\end{proof}

This is a modest completeness result, but it matters because it fixes the boundary honestly.  The finite Boolean logic is complete.  The later quantitative and analytic layers are proved sound; their completeness in infinite-dimensional settings remains open.

\subsection{Base change and the Beck--Chevalley law}

Local computations often refine two different objects over the same coarse base.  The expected compatibility is that quantifying along one side of a pullback square and then reindexing along the other gives the same result as reindexing first and quantifying second.

Consider finite maps
\[
   f:X\to Z,
   \qquad
   g:Y\to Z,
\]
and their pullback
\[
   X\times_ZY
   =\{(x,y):f(x)=g(y)\}.
\]
Let $p_X$ and $p_Y$ denote the two projections.

\begin{theorem}[Finite Beck--Chevalley identity]
For every predicate $A\subseteq Y$,
\[
   f^*(\exists_g A)
   =
   \exists_{p_X}(p_Y^*A).
\]
The dual universal identity also holds.
\end{theorem}

\begin{proof}
A point $x\in X$ belongs to the left side exactly when $f(x)$ lies in $\exists_gA$.  This means that there exists $y\in A$ with $g(y)=f(x)$.  Equivalently, there exists a pullback point $(x,y)\in X\times_ZY$ whose $Y$-projection lies in $A$.  That is precisely membership in the right side.  The universal statement is proved by replacing existence with the assertion that every compatible $y$ lies in $A$.
\end{proof}

The theorem is the finite local-to-global consistency condition behind many packet diagrams.  It is also the first sign that Resolution Logic belongs naturally to indexed or fibred categorical logic: each resolution carries a predicate algebra, and refinement maps carry reindexing functors with adjoints.

\subsection{Quantitative judgments}

The Boolean layer records which packets satisfy a predicate.  Most calculations also need a quantitative statement.  Let $X$ be a normed space.

\begin{definition}[Certified approximation]
A certified approximation at resolution $\mathcal R$ is a tuple
\[
   (\widehat x,\mathcal R,\varepsilon,c,\eta,\mathcal E)
\]
with the intended semantics
\[
   \|x-\widehat x\|\le\varepsilon
\]
except on an event of probability at most $\eta$.  The number $c$ records cost, and $\mathcal E$ is an evidence ledger.
\end{definition}

The deterministic exact case has $\varepsilon=\eta=0$.  A finite floating-point computation may have a nonzero deterministic residual but zero statistical failure probability.  A randomized certificate may carry both.  A conditional theorem records its assumptions in $\mathcal E$ rather than quietly deleting them after composition.

Grades combine conservatively.  Sequential error and cost are added, and failure probabilities are combined by the union bound:
\[
   (\varepsilon,c,\eta)
   \oplus
   (\varepsilon',c',\eta')
   =
   (\varepsilon+\varepsilon',c+c',\min(1,\eta+\eta')).
\]
This rule is intentionally simple.  Sharper covariance-aware combinations can be added when independence or dependence information is available.

\begin{proposition}[Grade monoid]
The operation $\oplus$ is associative and commutative, with unit $(0,0,0)$.
\end{proposition}

\begin{proof}
Addition of the first two coordinates is associative and commutative.  The third coordinate is truncated addition on $[0,1]$; direct inspection shows
\[
   \min(1,\min(1,a+b)+c)
   =
   \min(1,a+b+c),
\]
and similarly under permutation.  The zero triple is the unit.
\end{proof}

\subsection{Weakening, refinement, and coarsening}

Every sound logic admits weakening: a stronger certificate may be replaced by a weaker one.

\begin{theorem}[Weakening]
If
\[
   \Gamma\vdash_{\mathcal R}\widehat x:x[\varepsilon,c,\eta;\mathcal E],
\]
then for every $\delta,d,\theta\ge0$,
\[
   \Gamma\vdash_{\mathcal R}\widehat x:x
   [\varepsilon+\delta,c+d,\min(1,\eta+\theta);\mathcal E']
\]
whenever $\mathcal E'$ contains the original evidence.
\end{theorem}

\begin{proof}
The original error ball is contained in the enlarged error ball.  Increasing cost or failure allowance cannot invalidate a statement that was already valid under the smaller bounds.
\end{proof}

A refinement may be exact, while a coarsening usually carries a reconstruction residual.

\begin{theorem}[Resolution transport]
Suppose a transition $r:\mathcal R\to\mathcal S$ has Lipschitz constant $L_r$, representation defect $\delta_r$, cost $c_r$, and failure probability $\eta_r$.  Then
\[
   \widehat x\Vdash_{\mathcal R}x[\varepsilon,c,\eta]
\]
implies
\[
   r\widehat x\Vdash_{\mathcal S}rx
   [L_r\varepsilon+\delta_r,c+c_r,\min(1,\eta+\eta_r)].
\]
\end{theorem}

\begin{proof}
Write $\widetilde r$ for the implemented transition and assume
\[
   \|\widetilde r(z)-r(z)\|\le\delta_r
\]
on the certified ball.  Then
\[
\begin{aligned}
   \|\widetilde r(\widehat x)-r(x)\|
   &\le
   \|\widetilde r(\widehat x)-r(\widehat x)\|
   +\|r(\widehat x)-r(x)\|\\
   &\le
   \delta_r+L_r\varepsilon.
\end{aligned}
\]
The grade update follows from sequential composition.
\end{proof}

Exact refinement has $L_r=1$ and $\delta_r=0$.  A compression map has a named residual $\delta_r$.  This makes the logic faithful to the no-free-refinement principle: a finer description is not automatically cheaper, and a coarser one is not automatically harmless.

\begin{figure}[H]
\centering
\includegraphics[width=.78\linewidth]{figures/logic/03_refinement_frontier.png}
\caption{The useful question is not whether finer is better in isolation.  It is where the error decrease is worth the additional cost for the purpose at hand.}
\end{figure}

\subsection{Linear, multilinear, and tensor rules}

The linear rule is familiar.

\begin{theorem}[Linear transport]
If $A:X\to Y$ is bounded and
\[
   \widehat x\Vdash x[\varepsilon],
\]
then
\[
   A\widehat x\Vdash Ax[\|A\|\varepsilon].
\]
\end{theorem}

\begin{proof}
This is the operator norm inequality
\[
   \|A\widehat x-Ax\|
   \le
   \|A\|\,\|\widehat x-x\|.
\]
\end{proof}

The nonlinear calculus needs the corresponding multilinear rule.

\begin{theorem}[Multilinear resolution rule]
Let $M:X_1\times\cdots\times X_m\to Y$ be bounded and $m$-linear.  Suppose
\[
   \|x_j-\widehat x_j\|\le\varepsilon_j.
\]
Then
\[
\begin{aligned}
&\|M(x_1,\ldots,x_m)-M(\widehat x_1,\ldots,\widehat x_m)\|\\
&\qquad\le
\|M\|
\left[
\prod_{j=1}^m(\|\widehat x_j\|+\varepsilon_j)
-
\prod_{j=1}^m\|\widehat x_j\|
\right].
\end{aligned}
\]
\end{theorem}

\begin{proof}
Write $x_j=\widehat x_j+e_j$.  Multilinearity expands
\[
   M(x_1,\ldots,x_m)
\]
into one retained term $M(\widehat x_1,\ldots,\widehat x_m)$ and one term for every nonempty subset of error slots.  The norm of the term indexed by $S\subseteq\{1,\ldots,m\}$ is bounded by
\[
   \|M\|
   \prod_{j\in S}\varepsilon_j
   \prod_{j\notin S}\|\widehat x_j\|.
\]
Summing over nonempty $S$ gives the difference between the two products in the theorem.
\end{proof}

For $m=2$, this becomes
\[
   \|M\|
   \bigl(
   \|\widehat x\|\eta
   +\|\widehat y\|\varepsilon
   +\varepsilon\eta
   \bigr),
\]
which recovers the original bilinear rule.

Tensor products obey the same geometry:
\[
\|x\otimes y-\widehat x\otimes\widehat y\|
\le
\|\widehat x\|\eta
+
\|\widehat y\|\varepsilon
+
\varepsilon\eta.
\]
The tensor rule therefore transports both the values and their named resolutions into the product resolution $\mathcal R\otimes\mathcal S$.

\begin{figure}[H]
\centering
\includegraphics[width=.80\linewidth]{figures/logic/02_soundness_sweep.png}
\caption{A randomized finite-dimensional sweep compares observed errors with their certified linear and bilinear bounds.  Every ratio remains below one, as soundness requires.}
\end{figure}

\subsection{Nonlinear maps and majorant reasoning}

A general nonlinear map cannot be controlled by its value at one point.  It needs a local regularity statement.  The simplest is a Lipschitz bound.

\begin{theorem}[Lipschitz rule]
If $F:X\to Y$ is $L$-Lipschitz on the certified ball around $\widehat x$, then
\[
   \widehat x\Vdash x[\varepsilon]
   \quad\Longrightarrow\quad
   F(\widehat x)\Vdash F(x)[L\varepsilon].
\]
\end{theorem}

\begin{proof}
By definition,
\[
   \|F(\widehat x)-F(x)\|
   \le
   L\|\widehat x-x\|
   \le L\varepsilon.
\]
\end{proof}

For an analytic packet map
\[
   F(x)=\sum_{m\ge0}A_m[x,\ldots,x],
\]
a scalar majorant gives a computable local Lipschitz constant.  Let
\[
   \mathfrak m_F(t)=\sum_{m\ge0}\|A_m\|t^m.
\]
On the ball of radius $R$,
\[
   L_F(R)
   \le
   \mathfrak m_F'(R)
   =
   \sum_{m\ge1}m\|A_m\|R^{m-1}.
\]
Therefore a certified input ball of radius $\varepsilon$ is transported to an output ball of radius at most $\mathfrak m_F'(R)\varepsilon$, provided the entire input ball lies within radius $R$.

The packetwise version is sharper.  Each nonlinear hyperedge carries its own multilinear norm and input radii.  The total output error is obtained by summing the certified contributions of the active hyperedges and the residual bound of the omitted ones.

\begin{figure}[H]
\centering
\includegraphics[width=.80\linewidth]{figures/logic/11_hyperedge_logic.png}
\caption{A nonlinear transfer hyperedge is not merely an interaction coefficient.  It carries an error rule determined by its multilinear norm and the certified radii of its input packets.}
\end{figure}

\subsection{Nonlinear refinement and interaction trees}

The nonlinear transfer section proves that jets compose through packet-resolved Fa\`a di Bruno trees.  Resolution Logic adds the certification layer.  Every leaf of an interaction tree carries an input judgment.  Every internal node carries a multilinear or analytic transport rule.  The root carries the final error, cost, failure probability, and evidence ledger.

\begin{theorem}[Soundness of finite nonlinear derivation trees]
Consider a finite rooted tree whose leaves are sound certified approximations.  Suppose every internal node applies one of the sound rules proved above: linear transport, multilinear transport, tensoring, certified chart change, or weakening.  Then the root certificate is sound.
\end{theorem}

\begin{proof}
Proceed by induction on the height of the tree.  A leaf is sound by hypothesis.  Suppose every premise of an internal node is sound.  The theorem attached to the node gives a valid bound for the exact outputs represented by those premises, so the node conclusion is sound.  Finite induction reaches the root.
\end{proof}

This is the logical counterpart of operadic composition.  Values compose by substitution of interaction trees; certificates compose by the same tree with a quantitative rule at every vertex.

\begin{figure}[H]
\centering
\includegraphics[width=.90\linewidth]{figures/logic/08_derivation_tree.png}
\caption{A derivation tree may combine local chart certificates with a nonlinear rule.  The representation changes and the analytic error bound remain visible rather than disappearing into one final number.}
\end{figure}

\subsection{Pruning nonlinear hyperedges}

Sparse nonlinear computation omits weak interaction hyperedges.  A safe pruning rule must certify what those omissions can contribute on the current input ball.

Let $M_E$ denote the sum of retained multilinear hyperedges and $M_{E^c}$ the omitted part.  If
\[
   \|M_{E^c}\|\le\delta
\]
and the $m$ input packets satisfy
\[
   \|x_j\|\le r_j,
\]
then
\[
   \|M_{E^c}(x_1,\ldots,x_m)\|
   \le
   \delta\prod_{j=1}^m r_j.
\]
The radii may be replaced by certified upper bounds
\[
   r_j=\|\widehat x_j\|+\varepsilon_j.
\]
Thus pruning is a logical rule with an explicit residual, not an informal declaration that a coefficient looked small.

\begin{theorem}[Certified hyperedge pruning]
If the omitted $m$-linear transfer has norm at most $\delta$, then the retained nonlinear approximation has additional deterministic error at most
\[
   \delta
   \prod_{j=1}^m
   (\|\widehat x_j\|+\varepsilon_j).
\]
\end{theorem}

\begin{proof}
Every exact input lies in its certified ball, hence has norm at most $\|\widehat x_j\|+\varepsilon_j$.  Apply the multilinear operator norm bound to the omitted transfer.
\end{proof}

\subsection{Fixed points and implicit branches}

Many nonlinear problems are not evaluations but equations.  A useful logic should certify a computed solution without requiring the exact solution to be known.

\begin{theorem}[A posteriori contraction certificate]
Let $F$ be a contraction with constant $0\le q<1$ on a complete invariant set, and let $x_*$ be its fixed point.  For any approximate point $\widehat x$ in that set,
\[
   \|\widehat x-x_*\|
   \le
   \frac{\|F(\widehat x)-\widehat x\|}{1-q}.
\]
\end{theorem}

\begin{proof}
Using $F(x_*)=x_*$,
\[
\begin{aligned}
\|\widehat x-x_*\|
&\le
\|\widehat x-F(\widehat x)\|
+
\|F(\widehat x)-F(x_*)\|\\
&\le
\|\widehat x-F(\widehat x)\|
+q\|\widehat x-x_*\|.
\end{aligned}
\]
Move the final term to the left.
\end{proof}

The denominator records the approach to instability.  A small residual is not enough when $q$ is close to one.

\begin{figure}[H]
\centering
\includegraphics[width=.78\linewidth]{figures/logic/06_fixed_point_certificate.png}
\caption{The affine control saturates the a posteriori fixed-point bound.  In nonlinear problems the same formula remains valid whenever the contraction hypothesis is proved on the invariant region.}
\end{figure}

For an implicit equation
\[
   G(x,p)=0,
\]
assume the partial derivative $D_xG$ is invertible with inverse norm at most $M$, the parameter derivative has norm at most $K$, the parameter is known to error $\varepsilon_p$, and the computed state has residual $\rho$.  If the nonlinear feedback in the Newton correction is bounded by $\theta<1$, then the local state error obeys
\[
   \|x-\widehat x\|
   \le
   \frac{M(K\varepsilon_p+\rho)}{1-\theta}.
\]
This is the implicit-function analogue of the contraction certificate.  It tells the logic exactly which assumptions are carrying the conclusion.

\subsection{Predicates that may remain undecided}

A quantitative certificate need not decide every Boolean predicate.  This is a feature rather than a defect.

Let $\phi:X\to\mathbb R$ be $L_\phi$-Lipschitz and suppose
\[
   \widehat x\Vdash x[\varepsilon].
\]
Then
\[
   \phi(x)
   \in
   [\phi(\widehat x)-L_\phi\varepsilon,
    \phi(\widehat x)+L_\phi\varepsilon].
\]
For the predicate $\phi(x)\le t$ there are three cases:
\begin{enumerate}[label=\arabic*.]
\item if the upper endpoint is at most $t$, the predicate is certified true;
\item if the lower endpoint exceeds $t$, it is certified false;
\item otherwise the current resolution does not decide it.
\end{enumerate}

The undecided case is precisely where refinement is justified.  A logic that forces a binary answer here would be inventing information that the representation does not contain.

\subsection{Evidence does not collapse into one label}

A large mathematical project often contains several kinds of support for one proposition.  There may be a theorem under hypotheses, a finite exact computation, a numerical experiment, and a conjectural extrapolation.  These forms of support should be combined without being confused.

Define an evidence ledger
\[
   \mathcal E
   =
   (\mathcal P,\mathcal A,\mathcal C,\mathcal Q),
\]
where
\begin{itemize}
\item $\mathcal P$ lists proof steps;
\item $\mathcal A$ lists assumptions;
\item $\mathcal C$ lists computations;
\item $\mathcal Q$ lists conjectural inputs.
\end{itemize}
Composition takes unions in each component.  A computation never migrates into $\mathcal P$ merely because it agrees with a theorem.  A conditional theorem never loses $\mathcal A$ after being used downstream.

\begin{figure}[H]
\centering
\includegraphics[width=.88\linewidth]{figures/logic/09_status_ledger.png}
\caption{One proposition may have several kinds of support.  The ledger combines them without promoting one kind into another.}
\end{figure}

\begin{proposition}[Evidence preservation]
If two judgments are composed, the evidence ledger of the conclusion contains the componentwise union of the premise ledgers together with the new rule.  In particular, composition cannot turn a conjectural premise into an unconditional theorem or a finite computation into a proof.
\end{proposition}

\begin{proof}
This is built into the ledger operation.  Each evidence class is stored in a separate component and composition performs no coercion between components.
\end{proof}

\subsection{Resolution atlases and local truth}

Noncommuting projector systems may not admit one exact common refinement.  Resolution Logic therefore cannot assume that every claim lives in one global chart.  Instead it uses an atlas.

Let $\mathcal R_a$ be local resolutions with transition maps
\[
   g_{ba}:\mathcal R_a\to\mathcal R_b.
\]
A local judgment in chart $a$ can be transported to chart $b$ with a Lipschitz factor and a transition defect.  Local representatives may then be glued using a partition of unity or a chosen reference chart.

\begin{theorem}[Finite certified gluing]
Suppose local certificates $\widehat x_a\Vdash x_a[\varepsilon_a]$ have been transported into one reference chart, and let nonnegative weights $w_a$ sum to one.  If the overlap inconsistency is at most $\delta$, then
\[
   \widehat x=\sum_aw_a\widehat x_a
\]
satisfies
\[
   \|x-\widehat x\|
   \le
   \sum_aw_a\varepsilon_a+\delta.
\]
\end{theorem}

\begin{proof}
Choose compatible exact local representatives $x_a$ of the global object.  Their weighted sum is $x$ up to the overlap defect.  Then
\[
\begin{aligned}
\left\|x-\sum_aw_a\widehat x_a\right\|
&\le
\left\|x-\sum_aw_ax_a\right\|
+
\left\|\sum_aw_a(x_a-\widehat x_a)\right\|\\
&\le
\delta+\sum_aw_a\varepsilon_a.
\end{aligned}
\]
\end{proof}

\begin{figure}[H]
\centering
\includegraphics[width=.82\linewidth]{figures/logic/12_atlas_semantics.png}
\caption{Local truth is not second-class truth.  It becomes global when the chart transitions and overlap residual satisfy a gluing theorem.}
\end{figure}

\begin{figure}[H]
\centering
\includegraphics[width=.80\linewidth]{figures/logic/07_gluing_certificate.png}
\caption{The numerical gluing sweep keeps the observed global error below the propagated local certificate.}
\end{figure}

\subsection{Modal notation}

It is sometimes useful to package the judgments in modal notation.  Write
\[
   [\mathcal R]_{\varepsilon}\phi
\]
for a statement certified in resolution $\mathcal R$ to deterministic tolerance $\varepsilon$.  For a refinement arrow $r:\mathcal R\to\mathcal S$, write
\[
   \langle r\rangle_{\delta}\phi
\]
when $\phi$ becomes visible after transport along $r$ with additional defect at most $\delta$.  Compression may be written
\[
   C_{\mathcal R\to\mathcal S}^{\delta}\phi.
\]

These symbols do not introduce new theorems by themselves.  Their purpose is to make the representation dependence of a statement explicit.  The proved semantics in the finite theory is supplied by the pullback and quantifier adjunctions, the normed transport rules, and the atlas gluing theorem.

\subsection{A planner for resolution-dependent reasoning}

A calculation rarely wants the finest possible resolution.  It wants a resolution sufficient for a task at acceptable cost.  Let each candidate judgment $J$ carry error $\varepsilon(J)$, cost $c(J)$, failure probability $\eta(J)$, and an evidence penalty $s(J)$ chosen by the application.  Define
\[
   \mathcal J(J)
   =
   \lambda_\varepsilon\varepsilon(J)
   +\lambda_cc(J)
   +\lambda_\eta\eta(J)
   +s(J).
\]
The planner selects a sound judgment minimizing $\mathcal J$.

There is no universal choice of the weights.  A proof audit places nearly infinite weight on logical status and a strict tolerance on residual error.  An exploratory calculation may tolerate a larger residual in exchange for lower cost.  The point is not to reduce mathematics to one scalar objective.  It is to expose the choice that was already being made implicitly.

\begin{figure}[H]
\centering
\includegraphics[width=.82\linewidth]{figures/logic/04_path_planner.png}
\caption{A direct chart translation is preferred when cost dominates.  The two-stage refinement path is preferred once error receives enough weight.}
\end{figure}

A local refinement policy follows immediately.  If packet $i$ has residual $r_i$, refinement cost $c_i$, and penalty parameter $\lambda$, open the packet when
\[
   r_i>\lambda c_i.
\]
More elaborate policies can include bridge value, instability, or uncertainty, but even this simple rule prevents the vague instruction to ``refine wherever necessary.''

\subsection{A worked nonlinear example}

Consider the quadratic map $F:\mathbb R^2\to\mathbb R^2$,
\[
F(x_1,x_2)
=
\begin{pmatrix}
0.4x_1+0.08x_1x_2\\
-0.3x_2+0.05x_1^2
\end{pmatrix}.
\]
Resolve the source into the coordinate packets $P_1x=(x_1,0)$ and $P_2x=(0,x_2)$.  The first component has a linear self-edge $1\to1$ and a quadratic hyperedge $(1,2)\to1$.  The second has a linear self-edge $2\to2$ and a quadratic hyperedge $(1,1)\to2$.

Let
\[
   \widehat x=(0.4,-0.25)+e,
   \qquad
   \|e\|\le\varepsilon.
\]
On the ball $\|x\|\le R$, one may bound the Jacobian norm by $L(R)$.  The Lipschitz rule gives
\[
   \|F(\widehat x)-F(x)\|
   \le
   L(R)\varepsilon.
\]
A packetwise calculation can improve this by using separate radii for $x_1$ and $x_2$ and the exact multilinear hyperedge norms.

The computational sweep uses a conservative global $L$ and random perturbation directions.  The ratio of observed error to certified error remains below one throughout the tested range.

\begin{figure}[H]
\centering
\includegraphics[width=.80\linewidth]{figures/logic/05_nonlinear_majorant.png}
\caption{The nonlinear certificate is conservative, as intended, but remains valid across several orders of magnitude in the input radius.}
\end{figure}

\subsection{Relationship to the nonlinear transfer operad}

The nonlinear transfer operad and Resolution Logic answer different questions.

The operad asks:
\begin{quote}
What are the interaction tensors, how do they compose, and how do they refine?
\end{quote}

Resolution Logic asks:
\begin{quote}
What may be concluded from an approximate representation of the inputs, what residual does each interaction create, and which assumptions or computations support the conclusion?
\end{quote}

The two structures fit together without identification.  The operadic tree gives the substitution pattern.  The logical derivation tree gives the certificate propagation pattern.  They have the same combinatorial skeleton but different labels.

This distinction matters in practice.  A nonlinear packet map may be algebraically exact while the input data are uncertain.  Conversely, the input may be exact while a sparse hyperedge approximation introduces a residual.  The complete judgment records both sources of error separately before combining them.

\subsection{Categorical form of the finite theory}

Let $\mathbf{Res}$ be the finite category whose objects are packet resolutions and whose arrows are refinement maps.  Assign to each resolution $\mathcal R$ its Boolean algebra of predicates $\mathbf{Pred}(\mathcal R)$.  A refinement arrow $r:\mathcal R_f\to\mathcal R_c$ induces the reindexing map
\[
   r^*:\mathbf{Pred}(\mathcal R_c)
   \to
   \mathbf{Pred}(\mathcal R_f),
\]
with left and right adjoints $\exists_r$ and $\forall_r$.

Thus the finite Boolean fragment is an indexed Boolean doctrine.  The quantitative judgments form an additional graded layer over the same base.  Linear, multilinear, and chart-transition rules act as sound morphisms between graded fibers.

\begin{theorem}[Finite indexed semantics]
Finite packet resolutions, refinement maps, Boolean packet predicates, pullback, and the two refinement quantifiers define an indexed Boolean doctrine satisfying the finite Beck--Chevalley law for pullback squares.
\end{theorem}

\begin{proof}
Functoriality of pullback follows from composition of functions.  The Galois adjunction theorem supplies the two quantifiers.  Boolean operations are computed pointwise.  The finite Beck--Chevalley theorem was proved above.  These are precisely the required finite indexed structures.
\end{proof}

The theorem is deliberately finite.  Extending it to infinite resolutions, measurable predicates, sheaf semantics, or a genuine resolution topos requires additional hypotheses on admissible covers, limits, and continuity.

\subsection{Soundness theorem for the implemented calculus}

The Python implementation mirrors the finite theorems rather than inventing a second informal semantics.

\begin{theorem}[Soundness of the implemented finite rules]
Assume every input certificate satisfies its declared deterministic error bound.  Then the outputs produced by the implemented exact, weakening, linear, multilinear, tensor, Lipschitz, chart-change, compression, gluing, fixed-point, and hyperedge-pruning rules satisfy their declared deterministic error bounds.
\end{theorem}

\begin{proof}
The exact and weakening rules are immediate.  The linear, multilinear, tensor, and Lipschitz rules are the theorems proved above.  A chart change is a linear rule plus its declared transition defect.  Compression computes the residual directly.  Gluing uses the certified convex-combination theorem.  The fixed-point rule is the a posteriori contraction theorem.  Hyperedge pruning is the omitted multilinear norm estimate.  Composition of the rules remains sound by induction on the derivation tree.
\end{proof}

The randomized verification sweep is not the proof of the theorem.  It is a check that the implementation follows the theorem in representative finite cases.

\subsection{What this logic does not yet claim}

The finite theory is substantial, but its boundary should remain visible.

\begin{enumerate}[label=\arabic*.]
\item It does not yet provide a complete internal language for arbitrary infinite-dimensional projector categories.
\item It does not prove that every useful mathematical representation belongs to one resolution category.
\item It does not turn numerical evidence into proof.
\item It does not decide a predicate when the certified interval crosses its threshold.
\item It does not provide a topos until a site of admissible resolutions and covers is specified and the sheaf axioms are proved.
\item It does not make nonlinear global inversion automatic; the local contraction, majorant, or Kantorovich hypotheses remain essential.
\item It does not erase the distinction between incompatible local charts.  It transports and glues them only when transition data have been supplied.
\end{enumerate}

The next natural theorem is a representation and completeness theorem for a carefully chosen infinite category of resolutions.  The finite indexed doctrine, the normed certificate calculus, and the nonlinear operad now provide a concrete base from which that theorem can be stated without pretending it has already been proved.

\subsection{Closing perspective}

Resolution Theory began with the claim that representation should remain inside the mathematics.  The nonlinear extension shows that this remains true when several packets interact at once.  Resolution Logic adds the final missing discipline: a representation must carry not only a value but also the conditions under which that value may be trusted.

The resulting picture has three synchronized structures:
\[
\boxed{
\begin{aligned}
\text{resolution geometry}
&:\text{ how representations refine, change, and glue},\\
\text{nonlinear transfer operad}
&:\text{ how packet interactions compose},\\
\text{resolution logic}
&:\text{ how truth, error, cost, and evidence move through those operations}.
\end{aligned}}
\]

None of the three replaces the others.  Together they make it possible to ask a sharper question than whether one formula is correct: at what resolution is it correct, which interactions carry it, what was omitted, how large is the residual, and what kind of evidence supports the conclusion?


\section{Philosophy and methodology: proofs as compressed discovery fields}
\label{sec:proof-path-methodology}

Resolution Theory would be incomplete if it described only mathematical objects and ignored the way mathematical work itself is represented.  A proof, a program, a numerical experiment, and a research archive are also resolved objects.  Each has a coarse public form and a larger internal field of examples, failed routes, temporary definitions, computations, and questions.  The section below develops that idea carefully.  It does not weaken proof.  It asks what a proof leaves behind when it becomes short enough to publish.

\subsection{Why proof methodology belongs inside Resolution Theory}

A completed proof is usually presented as a clean implication
\[
        A\vdash B.
\]
That form is indispensable.  It is the certificate.  But it is not always the whole mathematical event.  The path by which a theorem was found may have been closer to
\[
        A\to S_1\to S_2\to S_3\to B,
\]
with side branches, failed analogies, examples, computational traces, boundary cases, and definitions that were later compressed away.  The final proof may be shorter and better for verification while being poorer as a map of future possibilities.

The thesis of this paper is:
\[
\boxed{\begin{minipage}{0.86\linewidth}
Proof compression can preserve verification value while erasing discovery value.  A mature account of proof should distinguish the final certificate from the larger discovery field that produced it.
\end{minipage}}
\]

This distinction is not an argument against compression.  Mathematics needs compression.  Good notation compresses.  Good proofs compress.  The point is that compression has a cost.  It may remove steps that are no longer necessary for the theorem but remain valuable as bridges to future mathematics.

\begin{figure}[H]
\centering
\includegraphics[width=0.92\linewidth]{figures/methodology/fig_architecture.png}
\caption{The proposed theory has five connected layers: proof-path objects, compression and loss, values and optimization, collective reading, and the growth of imagination, proof, and meaning.}
\end{figure}

\subsection{Basic objects}

\begin{definition}[Expanded proof path]
An expanded proof path is a finite or locally finite directed sequence
\[
        \pi_{\mathrm{exp}}=(S_0,S_1,\ldots,S_k),
        \qquad S_0=A,
        \qquad S_k=B.
\]
The states $S_i$ may include formal derivation steps, lemmas, examples, computational reductions, failed branches, boundary tests, diagrams, or temporary representations used during discovery.
\end{definition}

\begin{definition}[Compressed proof]
A compressed proof is a shorter representation
\[
        \pi_{\mathrm{cmp}}=C(\pi_{\mathrm{exp}})
\]
that preserves the endpoint implication $A\vdash B$ while omitting some parts of the expanded path.
\end{definition}

\begin{definition}[Lost-step set]
The lost-step set of a compression is
\[
        L_C(\pi)=\{S_i:S_i\in\pi_{\mathrm{exp}},\ S_i\notin C(\pi_{\mathrm{exp}})\}.
\]
It contains the steps, examples, failed moves, computations, and side branches removed from the final certificate.
\end{definition}

\begin{definition}[Latent bridge graph]
The latent bridge graph is a graph
\[
        G_L(\pi)=(V_L,E_L)
\]
whose vertices include lost steps and whose edges represent possible transfers from those steps to other results, fields, examples, or conjectures.
\end{definition}

These definitions separate two questions.  The first is:
\[
        \text{Does the proof certify }A\vdash B?
\]
The second is:
\[
        \text{What did the path reveal while reaching }B?
\]
Both questions matter.

\begin{figure}[H]
\centering
\includegraphics[width=0.95\linewidth]{figures/methodology/fig_expanded_vs_compressed.png}
\caption{The expanded path may contain examples, failures, and boundary tests that vanish from the compressed certificate but remain useful for future mathematics.}
\end{figure}

\subsection{Values of proof representations}

A proof representation can be evaluated along several axes:
\[
        V(\pi)=\text{verification value},
\]
\[
        D(\pi)=\text{discovery value},
\]
\[
        M(\pi)=\text{meaning or explanatory value},
\]
\[
        A(\pi)=\text{auditability},
\]
\[
        F(\pi)=\text{future-bridge value},
\]
\[
        N(\pi)=\text{noise or interpretive instability},
\]
\[
        C_0(\pi)=\text{cost, length, or cognitive load}.
\]
These values should not be collapsed too quickly.  A proof can be excellent as a certificate and weak as a teaching object.  A research archive can be rich in future bridges and noisy as a final paper.  A formally checked proof can be correct and still hard to interpret.

\begin{proposition}[Compression can preserve validity while losing bridge value]
There exist expanded paths $\pi$ and compression maps $C$ such that $C(\pi)$ proves the same endpoint theorem as $\pi$, but $L_C(\pi)$ contains a step with a bridge to another result that is not visible in $C(\pi)$.
\end{proposition}

\begin{proof}
Let $\pi=(A,S,B)$, where $A\to S\to B$ is an expanded proof route.  Suppose a later compression produces a direct certificate $C(\pi)=(A,B)$.  If $S$ also has a bridge $S\to Y$ to a future result, then the compressed proof still verifies $A\vdash B$, but the visible node $S$ has been removed.  The verification value is preserved while the bridge to $Y$ is no longer present in the compressed object.
\end{proof}

This proposition is simple, but it marks the central distinction.  Verification value and discovery value are not the same property.

\begin{figure}[H]
\centering
\includegraphics[width=0.72\linewidth]{figures/methodology/fig_value_frontier.png}
\caption{Proof forms occupy different positions in the verification-discovery plane.  No single representation is best for every purpose.}
\end{figure}


\subsection{Axioms of proof-path theory}

The previous definitions become useful only if they are governed by simple structural principles.  The following axioms are not meant to be final foundations.  They are a compact way to state what the theory is trying to preserve.

\paragraph{Axiom 1: endpoint preservation.}
A valid compression must preserve the endpoint implication:
\[
        \pi_{\mathrm{exp}}:A\leadsto B
        \quad\Longrightarrow\quad
        C(\pi_{\mathrm{exp}}):A\vdash B.
\]
Compression is not allowed to change the theorem.

\paragraph{Axiom 2: path non-uniqueness.}
The same implication may have many expanded paths:
\[
        \Pi(A,B)=\{\pi:A\leadsto B\}.
\]
There is no reason to expect one canonical discovery path.  Different paths can expose different bridges.

\paragraph{Axiom 3: value non-collapse.}
The values
\[
        V,D,M,A,F,N,C_0
\]
are not reducible to one another in general.  A path can improve one value while reducing another.

\paragraph{Axiom 4: lost-step persistence.}
A removed step is not automatically meaningless.  Formally, there may exist
\[
        S\in L_C(\pi)
\]
with positive future bridge value
\[
        b(S)>0.
\]

\paragraph{Axiom 5: audit plurality.}
A mathematical candidate should be evaluated by at least three audits:
\[
        \text{formal audit},\qquad \text{meaning audit},\qquad \text{future-bridge audit}.
\]
No one audit exhausts the mathematical status of the object.

\begin{proposition}[No universal best compression]
If two goals assign different nonproportional weight vectors to the values of a proof representation, then there need not be a single presentation that is optimal for both goals.
\end{proposition}

\begin{proof}
Let $\pi_1$ be high in verification and low in discovery, and let $\pi_2$ be lower in verification but higher in discovery.  A certification goal can weight verification and cost heavily, while a discovery goal can weight future-bridge value heavily.  The induced utilities may satisfy
\[
        U_{\mathrm{cert}}(\pi_1)>U_{\mathrm{cert}}(\pi_2)
\]
and
\[
        U_{\mathrm{disc}}(\pi_2)>U_{\mathrm{disc}}(\pi_1).
\]
Thus the same presentation need not optimize both goals.
\end{proof}

\subsection{Depth and the recursive proof cube}

A proof step is rarely atomic.  A step may open into definitions, examples, sublemmas, computations, background theorems, or conventions.  This gives a recursive depth function.

\begin{definition}[Expansion depth]
Let $Q$ be a class of accepted primitives.  The depth of a proof state $S$ relative to $Q$ is
\[
        \Depth_Q(S)=
        \min\{d:S\text{ can be expanded to primitives in }Q\text{ in }d\text{ refinement layers}\}.
\]
The depth of a path is
\[
        \Depth_Q(\pi)=\max_i\Depth_Q(S_i).
\]
\end{definition}

Depth is not only a measure of difficulty.  It is a measure of how much scaffolding is hidden.  A proof can be short and deep.  Another proof can be long and shallow.  These are different objects.

\begin{definition}[Recursive packet expansion]
A recursive packet expansion of a proof step $S$ is a decomposition
\[
        S=\sum_j B_jS+R,
\]
where the $B_jS$ are substeps, examples, definitions, or computational packets, and $R$ is the unresolved residual of the step.
\end{definition}

The recursive cube metaphor is the informal version of this definition.  A visible block opens into another block field.  The process stops only when the purpose of the proof no longer requires deeper expansion.

\begin{remark}
The right depth depends on the goal.  A journal proof may stop earlier than a teaching proof.  A formal proof may expand some steps much deeper than a philosophical explanation.  A discovery archive may preserve shallow failed steps because their future bridge value is high.
\end{remark}

\subsection{Compression operators}

Compression is not one operation.  It comes in several forms.

\paragraph{Notation compression.}
A repeated pattern is replaced by notation.  This reduces length but can hide the examples that suggested the pattern.

\paragraph{Lemma compression.}
A chain of local moves is replaced by a lemma.  This improves readability but may hide boundary cases.

\paragraph{Reference compression.}
A proof says that a statement follows from a known theorem.  This improves verification for experts but can reduce meaning for readers outside the subfield.

\paragraph{Formal compression.}
A machine-checkable derivation records every primitive step but can become opaque at the human level.  This is the reverse of ordinary compression: it expands formal detail while compressing human explanation.

\paragraph{Conceptual compression.}
A large family of arguments is replaced by one structural principle.  This is often the most valuable compression, but it can erase the path by which the principle was discovered.

Let
\[
        C=C_{\mathrm{not}}C_{\mathrm{lem}}C_{\mathrm{ref}}C_{\mathrm{form}}C_{\mathrm{conc}}
\]
represent a composite compression.  Each factor has a different loss profile.  The lost-step set should therefore be indexed by compression type:
\[
        L_C(\pi)=L_{\mathrm{not}}(\pi)\cup L_{\mathrm{lem}}(\pi)\cup L_{\mathrm{ref}}(\pi)\cup L_{\mathrm{form}}(\pi)\cup L_{\mathrm{conc}}(\pi).
\]

\begin{proposition}[Different compressions lose different information]
Two compression maps $C_1$ and $C_2$ may preserve the same endpoint theorem while producing different lost-step sets.
\end{proposition}

\begin{proof}
A lemma compression and a reference compression may both prove $A\vdash B$.  The lemma compression can preserve the internal lemma statement while removing examples; the reference compression can remove both the lemma and the examples.  The endpoint is the same, but the removed material differs.
\end{proof}

\subsection{Mathematical meaning as a separate value}

Meaning value is not the same as verification value.  A result can be verified and still not yet understood.  Conversely, a compelling explanation can precede a complete proof.  Let
\[
        M(\pi)=M_{\mathrm{concept}}(\pi)+M_{\mathrm{example}}(\pi)+M_{\mathrm{transfer}}(\pi)+M_{\mathrm{memory}}(\pi).
\]
Here:
\begin{itemize}
\item $M_{\mathrm{concept}}$ measures the clarity of the central idea;
\item $M_{\mathrm{example}}$ measures the examples preserved;
\item $M_{\mathrm{transfer}}$ measures how well the proof reveals connections to other objects;
\item $M_{\mathrm{memory}}$ measures whether the path can be reconstructed later.
\end{itemize}

A formal certificate may have high $V$ and low $M$.  A teaching proof may have lower cost efficiency and higher $M$.  A philosophical account may have high $M$ and low $V$ unless paired with a formal argument.  These are not contradictions.  They are different coordinates of the same proof field.

\subsection{Goal-dependent modeling choices}

The formalism can be modeled in several ways.  The correct model depends on the question being asked.

\subsubsection{Graph model}
A proof field can be represented as a directed graph
\[
        G=(V,E).
\]
Vertices are statements, examples, computations, representations, or failures.  Edges are derivations, analogies, reductions, transfers, or attempted transfers.  Compression is graph contraction.  Lost steps are removed vertices.  Future-bridge value is reachability from the removed vertices.

\subsubsection{Information model}
A proof can be treated as a code.  The compressed proof minimizes description length.  The expanded proof preserves side information.  The question becomes:
\[
        \text{how much side information can be removed before future tasks suffer?}
\]

\subsubsection{Optimization model}
A proof presentation is a point in a feasible set.  Different goals define different objectives.  The Pareto frontier of proof presentations is the set of forms that cannot improve one value without reducing another.

\subsubsection{Social model}
A community reading is a noisy measurement process.  Individual errors may cancel; shared bias may remain.  A dissenting reading is not automatically correct, but it may reveal a packet outside the current consensus lens.

\subsubsection{Developmental model}
Mathematical understanding changes over time.  The same proof can move from discovery object to teaching object to formal-library object.  Its best representation can change as the field matures.

\subsection{Simulation interpretation}

The simulations in this paper are not empirical claims about a historical community.  They are models of possible dynamics.  The aim is to show that the variables can move independently.

The idea-growth simulation separates imagination, proof, meaning, bridge density, and question quality.  This separation is important because a tool or method can increase the number of candidate connections faster than it increases interpretation.  In such a regime, the limiting factor is no longer generation.  It is audit and meaning.

The collective-noise simulation separates individual variance from shared bias.  It shows why consensus can become sharper without becoming unbiased.  The lesson is not anti-consensus.  The lesson is that consensus needs modeling.

The bridge-preservation simulation shows how forgotten steps may increase future reach.  It is a formal version of a familiar research experience: old failed arguments can become useful after the field changes.

\subsection{A storage model for expanded proofs}

A practical implementation of proof-path theory would store at least the following data:
\[
        \mathcal A(\pi)=
        (\pi_{\mathrm{exp}}, C(\pi), L_C(\pi), G_L(\pi), \mathcal E, \mathcal B, \mathcal R).
\]
Here $\mathcal E$ is the set of examples, $\mathcal B$ is the set of boundary cases, and $\mathcal R$ is the set of residual questions.  A proof repository could then allow searches such as:
\begin{itemize}
\item find all lost steps connected to a given concept;
\item find failed proof attempts that later became boundary cases;
\item find examples removed by compression;
\item find proof paths with high discovery value but low verification value;
\item find compressed proofs whose lost neighborhoods have high bridge potential.
\end{itemize}

This would turn the archive of mathematical work into a searchable discovery field, not only a library of final theorems.

\subsection{Case studies in miniature}

\subsubsection{A standard theorem}
A textbook proof may state that a result follows from a standard compactness argument.  The compressed proof is efficient.  The expanded path might include the finite examples that explain why compactness is the right packet.  Removing those examples does not weaken the proof, but it can weaken the reader's ability to transfer the idea.

\subsubsection{A computational proof}
A computational proof may produce a certificate and a log.  The certificate verifies the result.  The log may contain failed branches that reveal how the search space was structured.  A proof-path archive would keep both.

\subsubsection{A wrong proof}
A wrong proof can have high discovery value if the failure identifies a boundary.  The formal audit rejects the proof.  The future-bridge audit may preserve the failed step because it points toward a correct generalization.

\subsubsection{A mature theory}
When a theory matures, many formerly difficult steps become standard.  The field benefits from compression, but future researchers may lose sight of the old intermediate objects.  Proof-path theory suggests preserving some of those objects as bridge metadata.

\subsection{Open problems}

\begin{enumerate}[label=(\arabic*)]
\item Define a rigorous category of proof paths and compression maps.
\item Develop invariants of lost-step sets.
\item Measure bridge value in historical or formal proof corpora.
\item Build proof archives that preserve compressed proofs and expanded discovery graphs.
\item Develop proof compilers that output different representations for different goals.
\item Model collective reading with separate individual-noise and shared-bias components.
\item Study meaning-preserving compression.
\item Formalize the relation between imagination reach, proof reach, and meaning reach.
\end{enumerate}

These are not merely philosophical questions.  They are design questions for future mathematical infrastructure.

\subsection{Optimization of proof presentation}

Let a goal be represented by weights
\[
        \theta=(\theta_V,\theta_D,\theta_M,\theta_A,\theta_F,\theta_N,\theta_C).
\]
Define
\[
\begin{aligned}
        U_\theta(\pi)
        ={}&\theta_VV(\pi)+\theta_DD(\pi)+\theta_MM(\pi)
        +\theta_AA(\pi)+\theta_FF(\pi)\\
        &-\theta_NN(\pi)-\theta_CC_0(\pi).
\end{aligned}
\]
Then the goal-dependent representation is
\[
        \pi_\theta^*=
        \operatorname*{arg\,max}_{\pi\in\Pi(A,B)}U_\theta(\pi).
\]
The formula does not decide whether the theorem is true.  Correctness remains necessary.  The formula decides which correct representation best serves a goal.

\subsubsection{Goal-specific objectives}

\paragraph{Certificate objective.}
\[
        \pi^*_{\mathrm{cert}}
        =\operatorname*{arg\,max}_\pi
        \left[V(\pi)+\alpha A(\pi)-\lambda C_0(\pi)-\mu N(\pi)\right].
\]
This objective favors short, auditable, low-noise proofs.

\paragraph{Discovery objective.}
\[
        \pi^*_{\mathrm{disc}}
        =\operatorname*{arg\,max}_\pi
        \left[D(\pi)+\alpha F(\pi)+\beta\Bridge(L_C(\pi))-\lambda N(\pi)\right].
\]
This objective preserves intermediate steps that may become future bridges.

\paragraph{Teaching objective.}
\[
        \pi^*_{\mathrm{teach}}
        =\operatorname*{arg\,max}_\pi
        \left[M(\pi)+\alpha D(\pi)-\lambda C_0(\pi)-\mu N(\pi)\right].
\]
This objective values recoverability of the path.

\paragraph{Archive objective.}
\[
        \pi^*_{\mathrm{arch}}
        =\operatorname*{arg\,max}_\pi
        \left[D(\pi)+F(\pi)+M(\pi)-\lambda N(\pi)\right].
\]
This objective accepts length in order to preserve future use.

\paragraph{Collective review objective.}
\[
        \pi^*_{\mathrm{collab}}
        =\operatorname*{arg\,max}_\pi
        \left[V(\pi)+M(\pi)+D(\pi)-\lambda N_{\mathrm{ind}}(\pi)-\kappa N_{\mathrm{coll}}(\pi)\right].
\]
This objective recognizes both individual noise and shared collective bias.

\begin{figure}[H]
\centering
\includegraphics[width=0.92\linewidth]{figures/methodology/fig_goal_weight_profiles.png}
\caption{Different goals assign different weights to verification, discovery, meaning, auditability, future-bridge value, noise, and cost.}
\end{figure}

\subsection{Toy calculation: representation choice by goal}

The next calculation is schematic.  It assigns normalized values to six proof presentations and evaluates them under six goals.  The point is not the particular numbers, but the dependence of the best proof form on the goal.

\begin{table}[H]
\centering
\scriptsize
\begin{tabular}{lrrrrrrr}
\toprule
representation & V & D & M & A & F & N & cost\\
\midrule
formal\_certificate & 0.96 & 0.22 & 0.44 & 0.95 & 0.18 & 0.08 & 0.141 \\
journal\_proof & 0.94 & 0.54 & 0.7 & 0.82 & 0.42 & 0.11 & 0.318 \\
teaching\_path & 0.88 & 0.66 & 0.92 & 0.76 & 0.58 & 0.13 & 0.482 \\
discovery\_archive & 0.68 & 0.96 & 0.64 & 0.5 & 0.91 & 0.31 & 1 \\
balanced\_paper & 0.93 & 0.78 & 0.84 & 0.86 & 0.7 & 0.12 & 0.424 \\
candidate\_map & 0.56 & 0.9 & 0.47 & 0.38 & 0.86 & 0.48 & 0.659 \\
\bottomrule
\end{tabular}
\caption{Normalized values assigned to six proof representations.}
\end{table}

\begin{table}[H]
\centering
\small
\begin{tabular}{llr}
\toprule
goal & best representation & utility\\
\midrule
certificate & formal\_certificate & 0.736 \\
journal & balanced\_paper & 0.795 \\
teaching & balanced\_paper & 0.752 \\
discovery & balanced\_paper & 0.732 \\
archive & balanced\_paper & 0.797 \\
collective\_review & balanced\_paper & 0.762 \\
\bottomrule
\end{tabular}
\caption{Goal-dependent winners in the toy optimization.}
\end{table}

\begin{figure}[H]
\centering
\includegraphics[width=0.92\linewidth]{figures/methodology/fig_goal_score_heatmap.png}
\caption{The same proof material is scored differently under different goals.}
\end{figure}

\begin{figure}[H]
\centering
\includegraphics[width=0.96\linewidth]{figures/methodology/fig_representation_tradeoff.png}
\caption{Different proof forms preserve different kinds of value.}
\end{figure}

The calculation supports a practical conclusion.  It is not enough to ask for the proof.  One should ask which proof presentation is being optimized: certificate, journal article, teaching path, discovery archive, formal library, or collaborative review.

\subsection{Lost steps and future bridge value}

For a lost step $S\in L_C(\pi)$, define a bridge value
\[
        b(S)=\sum_{Y\in\mathcal Y}w_Y\,\mathbf 1\{S\leadsto Y\},
\]
where $\mathcal Y$ is a set of possible future targets and $w_Y$ is the importance of target $Y$.  The total bridge value of the lost neighborhood is
\[
        B(L_C(\pi))=\sum_{S\in L_C(\pi)}b(S).
\]
A compression can therefore be measured by two losses:
\[
        \Delta_D(\pi)=D(\pi)-D(C(\pi)),
\]
\[
        \Delta_B(\pi)=B(L_C(\pi)).
\]
If $\Delta_B$ is large, the compressed proof may be excellent for certification and poor for future discovery.

\begin{figure}[H]
\centering
\includegraphics[width=0.82\linewidth]{figures/methodology/fig_forgotten_bridge_reach.png}
\caption{A simulation in which preserving intermediate steps increases future bridge reach.}
\end{figure}

\subsection{Collective noise and interpretation}

Every observer produces noise.  Let the reading of observer $i$ be
\[
        R_i=O+C+L_i+\eta_i,
\]
where $O$ is the object, $C$ is collective bias, $L_i$ is the observer's lens, and $\eta_i$ is individual mean-zero noise.  A simplified version is
\[
        R_i=O+C+\eta_i,
\]
with $\eta_i$ independent and $\operatorname{Var}(\eta_i)=\sigma^2$.  Then
\[
        \overline R_m=O+C+\frac1m\sum_{i=1}^m\eta_i,
\]
so
\[
        \mathbb E\,\overline R_m=O+C,
\]
\[
        \operatorname{Var}(\overline R_m)=\frac{\sigma^2}{m}.
\]
Individual noise averages down.  Shared bias remains.

\begin{proposition}[Consensus has a bias floor]
In the model above, increasing the number of independent readers drives the individual-noise variance to zero, but it does not remove the collective bias term $C$.
\end{proposition}

\begin{proof}
The expectation of the average reading is $O+C$.  The variance of the average of independent mean-zero errors is $\sigma^2/m$, which tends to zero.  The bias $C$ is not random individual noise, so it remains.
\end{proof}

\begin{table}[H]
\centering
\small
\begin{tabular}{rrrrr}
\toprule
readers & individual variance & shared bias & lens bias & total MSE\\
\midrule
1 & 1 & 0.38 & 0.24 & 1.2 \\
2 & 0.5 & 0.38 & 0.24 & 0.702 \\
5 & 0.2 & 0.38 & 0.24 & 0.402 \\
10 & 0.1 & 0.38 & 0.24 & 0.302 \\
20 & 0.05 & 0.38 & 0.24 & 0.252 \\
40 & 0.025 & 0.38 & 0.24 & 0.227 \\
80 & 0.0125 & 0.38 & 0.24 & 0.214 \\
120 & 0.00833 & 0.38 & 0.24 & 0.21 \\
\bottomrule
\end{tabular}
\caption{Collective-noise calculation.  More readers reduce individual variance, but shared bias and lens bias remain.}
\end{table}

\begin{figure}[H]
\centering
\includegraphics[width=0.82\linewidth]{figures/methodology/fig_collective_noise.png}
\caption{Consensus reduces individual noise but can preserve a bias floor.}
\end{figure}

Thus consensus is evidence, not verdict.  It may be strong evidence.  It is still a reading packet, not the object itself.

\subsection{Imagination, proof, meaning, and question quality}

Let
\[
        I_t=\text{imagination reach at time }t,
\]
\[
        P_t=\text{formal proof reach},
\]
\[
        M_t=\text{meaning or interpretive reach},
\]
\[
        B_t=\text{bridge density},
\]
\[
        Q_t=\text{quality of questions available to the community}.
\]
These quantities need not grow together.  A system can produce many candidate ideas without enough audit.  Formal proof can grow faster than meaning.  Meaning can stabilize slowly.  Question quality may be the missing variable, because better questions can convert noise into search direction.

A schematic update is
\[
        I_{t+1}=I_t+a_tI_t\left(1-\frac{I_t}{K_I}\right)+bB_t+\epsilon_t,
\]
\[
        P_{t+1}=P_t+c_tI_t\left(1-\frac{P_t}{K_P}\right)+dB_t+\eta_t,
\]
\[
        M_{t+1}=M_t+e_tP_t\left(1-\frac{M_t}{K_M}\right)+fQ_t+\zeta_t.
\]
The model is not meant as a law of history.  It is a way of expressing a possibility: tools may expand imagination, proof, and meaning at different rates.

\begin{table}[H]
\centering
\small
\begin{tabular}{rrrrrrr}
\toprule
time & imagination & proof & meaning & bridges & question quality & noise\\
\midrule
0 & 7 & 5 & 4 & 2 & 3 & 1.5 \\
20 & 27.7 & 6.88 & 5.2 & 4.82 & 6.89 & 1.17 \\
46 & 96.3 & 15.1 & 7.08 & 13.8 & 11.2 & 2.15 \\
60 & 202 & 27.7 & 9.57 & 29.9 & 17.4 & 3.85 \\
85 & 231 & 63.6 & 18.1 & 59.1 & 40.6 & 4.17 \\
109 & 235 & 96 & 33.1 & 81.2 & 81.7 & 3.71 \\
\bottomrule
\end{tabular}
\caption{Selected states from the idea-growth simulation.}
\end{table}

\begin{figure}[H]
\centering
\includegraphics[width=0.94\linewidth]{figures/methodology/fig_growth_curves.png}
\caption{Simulation of imagination reach, proof reach, meaning reach, bridge density, and question quality.}
\end{figure}

\begin{figure}[H]
\centering
\includegraphics[width=0.88\linewidth]{figures/methodology/fig_meaning_lag.png}
\caption{Meaning can lag behind imagination, proof, and bridge creation.}
\end{figure}

\subsection{Phase diagram for interpretation}

The next simulation varies candidate-generation rate and audit/interpretation rate.  The stability index is high when interpretation can keep pace with new candidate structures and low when candidate production outruns audit.

\begin{table}[H]
\centering
\small
\begin{tabular}{rrr}
\toprule
candidate rate & audit rate & stability\\
\midrule
0.04 & 0.02 & 0.446 \\
0.0864 & 0.0606 & 0.306 \\
0.133 & 0.101 & 0.21 \\
0.179 & 0.142 & 0.144 \\
0.2 & 0.16 & 0.122 \\
\bottomrule
\end{tabular}
\caption{Sample values from the interpretation-stability phase diagram.}
\end{table}

\begin{figure}[H]
\centering
\includegraphics[width=0.76\linewidth]{figures/methodology/fig_phase_diagram_meaning.png}
\caption{A schematic phase diagram.  If candidate generation outruns audit and interpretation, stability drops.}
\end{figure}

The philosophical conclusion is not that candidate generation should be slowed.  It is that audit systems, archives, and meaning-preserving representations must scale with it.

\subsection{Discovery graphs}

The proof-path theory treats discovery as a graph.  Local clusters form first.  Weak long-range bridges often appear later.  These weak bridges may be more important than the strongest local edges, because they connect regions of the graph that otherwise remain separate.

\begin{figure}[H]
\centering
\includegraphics[width=0.98\linewidth]{figures/methodology/fig_discovery_graph_evolution.png}
\caption{A discovery graph grows through local clustering and weak long-range bridges.}
\end{figure}

A forgotten step may appear irrelevant inside the final proof and still be essential in the discovery graph.  The graph is larger than the certificate.

\subsection{The three-audit model}

Candidate mathematical connections should pass three audits:
\[
        \text{formal audit},
        \qquad
        \text{meaning audit},
        \qquad
        \text{future-bridge audit}.
\]
The formal audit asks whether the claim is derivable.  The meaning audit asks whether the claim can be interpreted.  The future-bridge audit asks what the claim opens.

\begin{figure}[H]
\centering
\includegraphics[width=0.88\linewidth]{figures/methodology/fig_audit_model.png}
\caption{A candidate bridge should be tested formally, interpretively, and by its future-bridge value.}
\end{figure}

This model prevents two opposite errors.  It prevents exploration from replacing proof, and it prevents formal proof from erasing the discovery field too early.

\subsection{How the theory can be modeled in different ways}

The same formal objects can be modeled differently depending on the goal.

\subsubsection{As a graph theory problem}
A discovery field can be modeled as a directed graph
\[
        G=(V,E),
\]
where vertices are statements, examples, computations, or representations, and edges are transfers.  Compression is graph contraction.  Lost steps are removed vertices.  Bridge value is future reachability.

\subsubsection{As an information theory problem}
A proof can be modeled as a code.  Compression reduces description length.  The cost is that side information may be removed.  Discovery value is then the value of side information for future tasks.

\subsubsection{As an optimization problem}
A proof presentation can be chosen by maximizing a utility function
\[
        U_\theta(\pi).
\]
Different goals yield different optimal presentations.

\subsubsection{As a learning problem}
The archive of proof paths can be treated as training data for future mathematical search.  Lost steps are not merely noise; they may be features for future theorem discovery.

\subsubsection{As a social epistemology problem}
Readers are noisy measurement systems.  Consensus reduces some noise and preserves other noise.  Dissent is not automatically correct, but it may contain a bridge that the consensus lens cannot see.

\subsection{Examples}

\subsubsection{Compressed proof versus teaching path}
A short proof may say that a result follows by a standard theorem.  A teaching path may show the examples, failed special cases, and boundary behavior that make the theorem feel inevitable.  Both are valid representations.  They optimize different values.

\subsubsection{Formal library versus research archive}
A formal library wants minimal ambiguity and high auditability.  A research archive may preserve failed attempts and exploratory reductions.  The library certifies.  The archive remembers.

\subsubsection{Collective review}
A group of readers may agree that a proof is uninteresting because it does not match the current field lens.  That agreement may reflect truth, or it may reflect shared bias.  The proof-path theory treats the collective reading as data, not as final reality.

\subsubsection{Candidate generation and meaning}
A candidate-generation tool can expand imagination by producing bridges quickly.  Those bridges are not proofs.  They are entries in a candidate field.  The key question is whether formal audit and meaning audit can keep pace with the growth of the field.

\subsection{Refinement program}

The theory can be developed further in several directions.

\paragraph{Proof archives with lost-step metadata.}
Formal repositories could store not only final derivations but expanded discovery fields:
\[
        \pi_{\mathrm{exp}},\quad C(\pi),\quad L_C(\pi),\quad G_L(\pi).
\]

\paragraph{Proof compilers with multiple outputs.}
A proof compiler could generate several versions:
\[
        \pi_{\mathrm{cert}},
        \quad
        \pi_{\mathrm{teach}},
        \quad
        \pi_{\mathrm{disc}},
        \quad
        \pi_{\mathrm{arch}}.
\]

\paragraph{Bridge-value estimation.}
One can study which lost steps later become useful in other fields.  This would give empirical content to $F(\pi)$ and $B(L_C(\pi))$.

\paragraph{Collective-noise models.}
Review systems could distinguish independent disagreement from shared bias.  Consensus would remain important, but it would be modeled rather than absolutized.

\paragraph{Meaning-preserving compression.}
The most important practical problem is to compress proofs without destroying too much interpretive and future-bridge value.  This is not only a writing problem.  It is a mathematical-design problem.


\subsection{Operational refinement}

The theory should not end as a fixed philosophical vocabulary.  It should give a working protocol for choosing how much of a proof path to preserve.  The most practical object is not a single proof representation but a representation portfolio:
\[
        \mathcal R(\pi)=
        \{\pi_{\mathrm{cert}},\pi_{\mathrm{teach}},\pi_{\mathrm{disc}},\pi_{\mathrm{arch}},\pi_{\mathrm{form}}\}.
\]
Here $\pi_{\mathrm{cert}}$ is optimized for verification, $\pi_{\mathrm{teach}}$ for meaning, $\pi_{\mathrm{disc}}$ for future bridge discovery, $\pi_{\mathrm{arch}}$ for preserving the search field, and $\pi_{\mathrm{form}}$ for formal audit.

For a goal vector $\theta$, define
\[
        U_\theta(\pi)
        =\theta_VV(\pi)+\theta_DD(\pi)+\theta_MM(\pi)+\theta_AA(\pi)+\theta_FF(\pi)
        -\theta_NN(\pi)-\theta_CC_0(\pi).
\]
The selected representation is
\[
        \pi^*(\theta)=\arg\max_{\pi\in\mathcal R(\pi_{\mathrm{exp}})}U_\theta(\pi).
\]
This formula makes the main methodological point explicit: there is no reason to expect the same proof form to be best for every goal.

\subsubsection{Refinement operator}

Let $\mathcal T$ be a current representation of a proof field.  A refinement step opens one packet, records what is gained, and measures the new residual:
\[
        \mathcal T_{k+1}=\operatorname{Refine}(\mathcal T_k,B_j)
        =\mathcal T_k\cup\{B_{j,1},\ldots,B_{j,m},R_j\}.
\]
The refinement is kept if it improves at least one required value without making the intended goal unusable:
\[
        \Delta U_\theta
        =U_\theta(\mathcal T_{k+1})-U_\theta(\mathcal T_k)>0.
\]
For archival goals, one may keep a refinement even when $\Delta U_\theta<0$ for the final paper, provided the bridge value is positive:
\[
        F(B_{j,r})>0.
\]
This is the formal version of the forgotten-step principle.  A step may be bad for the final text and good for the future search.

\subsubsection{Stopping rules}

A proof path can always be expanded further.  The practical question is where to stop.  The stopping rule should depend on purpose.

For certification, stop when the formal audit is stable:
\[
        V(\pi)\ge \tau_V
        \quad\text{and}\quad
        A(\pi)\ge \tau_A.
\]
For teaching, stop when meaning value no longer increases enough to justify the cost:
\[
        \Delta M < \lambda\Delta C_0.
\]
For discovery, stop later.  A high-cost expansion may be justified if it creates a high-value bridge:
\[
        \Delta F > \mu\Delta C_0.
\]
For archives, the stopping rule is weakest: preserve the object whenever its future value is uncertain but nontrivial.

The deeper principle is that expansion depth is not an absolute virtue.  It is a goal-dependent variable.  Mathematics advances by compression, but it also advances by knowing when not to compress yet.

\subsection{Final synthesis}

The refined theory can be summarized as a sequence of separations:
\[
        \text{proof}\ne\text{discovery path},
\]
\[
        \text{verification}\ne\text{meaning},
\]
\[
        \text{consensus}\ne\text{truth},
\]
\[
        \text{compression}\ne\text{losslessness},
\]
\[
        \text{failed step}\ne\text{worthless step}.
\]
These separations are not meant to weaken proof.  They are meant to protect everything proof does not record.  The final certificate remains essential.  The discovery field around the certificate may also contain mathematics that the certificate no longer needs.

A final proof is therefore best understood as one projection of a larger object.  The larger object is the proof path, its compression history, its lost-step set, its bridge graph, its interpretation history, and its future questions.  That larger object is what this paper calls a compressed discovery field.

\subsection{Conclusion}

The theory developed here does not replace proof.  It expands what counts as part of a proof's mathematical life.  A proof certifies an implication, but the discovery field around the proof may contain future mathematics.  Compression is necessary, but compression has losses.  Collective reading is necessary, but collective reading has noise.  Candidate generation is powerful, but it must be matched by audit and meaning.

The final thesis is:
\[
\boxed{\begin{minipage}{0.88\linewidth}
A mature philosophy of mathematical proof should study proofs as compressed discovery fields.  It should distinguish verification value from discovery value, track the lost neighborhoods created by compression, model collective noise, and design proof representations according to the goal they are meant to serve.
\end{minipage}}
\]


\subsection{Reproducibility of the methodology models}
The figures and tables in this section were generated by the accompanying Python program \texttt{proof\_path\_methodology\_figures.py}.  They are finite models of representation choice, proof compression, collective noise, bridge retention, and audit capacity.  The simulations are illustrations of the formal distinctions, not empirical laws of mathematical history.

\clearpage
\section{A final combined view: what mathematics, computer science, statistics, and technology mean}
\label{sec:final-combined-view}

The opening question can now be answered without reducing the four fields to slogans.  They remain different disciplines, with different objects and standards, but the theory developed in this appendix reveals a common architecture beneath their practice.

\begin{figure}[H]
\centering
\includegraphics[width=.90\linewidth]{figures/synthesis/07_combined_architecture.png}
\caption{The combined architecture.  The lower layers do not replace the upper ones.  They explain how representation, interaction, justification, and human judgment fit together.}
\end{figure}

\subsection{Mathematics}

Mathematics is the study of structures, transformations, invariants, and justified conclusions.  Resolution Theory adds one clause: these structures are encountered through representations, and the relation between representations is itself mathematical.  A mathematical object is therefore studied together with its admissible resolutions, the transfer laws visible in those resolutions, the residual left by compression, and the conditions under which reconstruction is exact.

In this view, a theorem does not become relative merely because its representation is recorded.  Exact truth remains exact.  What becomes explicit is the route by which truth is made visible.  Two resolutions may prove the same theorem while exposing different invariants, different algorithms, or different future questions.

\subsection{Computer science}

Computer science is the study of executable representations, algorithms, information flow, and resource bounds.  The packet calculus makes the representation-algorithm connection explicit.  A dense operator, a sparse transfer graph, a low-rank factorization, and a recursive packet tree can denote the same transformation while requiring radically different storage, arithmetic, communication, and verification cost.

The central computer-science question becomes:
\[
\text{Which resolution makes the required transfer exact or certified at the lowest total cost?}
\]
That question includes compilation, scheduling, memory movement, approximation, and hardware realization.  It also includes the cost of discovering the representation in the first place.

\subsection{Statistics}

Statistics is the study of inference from incomplete and noisy evidence.  Resolution Logic adds that every inference is made at a stated resolution.  A partition can reveal a conditional pattern or conceal it.  A coarse model can make a question undecidable that becomes decidable after refinement.  A posterior, confidence statement, or model comparison should therefore carry its observational resolution, approximation residual, failure probability, and evidence type.

This does not make inference arbitrary.  It makes its conditions visible.  A sound statistical conclusion is not only a number.  It is a number together with the experiment, model, resolution, uncertainty, and decision rule that justify its use.

\subsection{Technology}

Technology is the physical embodiment of a representation and a decision.  Sensors choose observables.  Data formats choose coordinates.  Processors choose instruction sets.  Interfaces choose what a person can see and control.  Every technological artifact therefore freezes some resolution into matter.

The engineering question is not merely whether a device works.  It is whether the chosen representation preserves the structure needed by the task, whether its residuals are measured, whether its failures are legible, and whether another resolution would provide a better balance of accuracy, cost, energy, and human control.

\subsection{A resolved practice object}

The four disciplines can be joined without identifying them.  Let a resolved practice object be
\[
\mathfrak P
=
(X,\mathcal R,\mathcal T,\mathcal J,\mathcal C,\mathcal E,\mathcal A),
\]
where
\begin{itemize}
\item $X$ is the object or problem;
\item $\mathcal R$ is the packet atlas or family of representations;
\item $\mathcal T$ is the linear transfer field;
\item $\mathcal J$ is the nonlinear jet hypergraph;
\item $\mathcal C$ is the certificate field carrying error, cost, and failure probability;
\item $\mathcal E$ is the typed evidence ledger;
\item $\mathcal A$ is the expanded archive of proof paths, computations, examples, and unresolved questions.
\end{itemize}
The purpose of this enlarged object is not to preserve every detail forever.  It is to make compression a named operation rather than an invisible disappearance.

\begin{definition}[Practice-coherent transformation]
A transformation $F:\mathfrak P\to\mathfrak Q$ is practice coherent when it transports the mathematical value, the relevant resolution data, the certificate, and the evidence ledger without changing their types.  A numerical experiment may support a proposition, but it does not become a proof under transport.  A conjecture may guide refinement, but it does not become an assumption unless that change is recorded.
\end{definition}

\subsection{The portfolio principle}

A single proof form cannot generally optimize verification, teaching, discovery, archival value, and cost at once.  The correct response is not to abandon compression.  It is to preserve a portfolio of representations.

\begin{theorem}[Portfolio dominance]
Let $\mathcal P=\{p_1,\ldots,p_m\}$ be a finite family of representations, and let $U_g(p)$ be the utility of representation $p$ for goal $g$.  Define
\[
U_g(\mathcal P)=\max_{p\in\mathcal P}U_g(p).
\]
Then for every $p_j\in\mathcal P$ and every goal $g$,
\[
U_g(\mathcal P)\ge U_g(p_j).
\]
If two goals have different unique maximizers, no single representation attains the portfolio utility for both goals.
\end{theorem}

\begin{proof}
The first statement is the defining property of a maximum.  For the second, let $p_a$ uniquely maximize $U_{g_1}$ and $p_b\ne p_a$ uniquely maximize $U_{g_2}$.  Any single representation differs from at least one of these two maximizers and therefore has strictly smaller utility for at least one goal.
\end{proof}

\begin{figure}[H]
\centering
\includegraphics[width=.94\linewidth]{figures/synthesis/02_goal_dependent_representation.png}
\caption{Different goals reward different proof representations.  The portfolio principle keeps several views of one mathematical object rather than pretending that one view serves every purpose.}
\end{figure}

\subsection{Typed evidence and compositional trust}

The evidence ledger prevents a familiar category error.  Proof, assumption, exact computation, numerical evidence, and conjecture may interact, but they should not collapse into one confidence score.

Suppose a certified stage carries a grade
\[
g=(\varepsilon,c,\eta),
\]
where $\varepsilon$ is deterministic error, $c$ is cost, and $\eta$ is failure probability.  If a second stage has grade $g'=(\varepsilon',c',\eta')$ and is $L$-Lipschitz in the output of the first, sequential composition gives
\[
\boxed{
 g'\circ_L g
 =
 (\varepsilon'+L\varepsilon,
  c+c',
  \min\{1,\eta+\eta'\}).
}
\]
The formula is elementary, but it is enough to make the audit trail compositional.

\begin{proposition}[Status preservation]
If two evidence ledgers are merged by taking the disjoint union inside each evidence class, then no proof, computation, numerical observation, assumption, or conjecture changes type under the merge.
\end{proposition}

\begin{proof}
The merge is performed componentwise in the coproduct
\[
\mathcal E
=
\mathcal E_{\mathrm{proof}}
\sqcup
\mathcal E_{\mathrm{assumption}}
\sqcup
\mathcal E_{\mathrm{exact}}
\sqcup
\mathcal E_{\mathrm{numerical}}
\sqcup
\mathcal E_{\mathrm{conjecture}}.
\]
No inclusion map identifies two distinct components, so evidence type is preserved.
\end{proof}

\begin{figure}[H]
\centering
\includegraphics[width=.88\linewidth]{figures/synthesis/05_evidence_ledger.png}
\caption{A theorem paper, a computational study, and a research archive can contain different mixtures of evidence.  The ledger records those mixtures without pretending they are interchangeable.}
\end{figure}

\subsection{Compression debt and forgotten bridges}

Compression is necessary.  The question is what it removes.  Let a discovery field contain steps with nonnegative future-bridge values $b_i$.  If $S$ is the retained set, define
\[
R(S)=\frac{\sum_{i\in S}b_i}{\sum_i b_i}
\]
and the normalized compression debt
\[
D(S)=1-R(S)+C_{\mathrm{reconstruct}}(S).
\]
This is a model, not an intrinsic universal invariant.  Its value is that it separates two different losses: the bridge value discarded now and the work required to recover context later.

\begin{figure}[H]
\centering
\includegraphics[width=.88\linewidth]{figures/synthesis/04_compression_debt.png}
\caption{Retaining the same number of steps can preserve very different amounts of future bridge value.  Compression should be guided by structure rather than length alone.}
\end{figure}

\subsection{Collective reading and shared noise}

Independent review is powerful because independent errors average down.  Correlated review is weaker than its head count suggests.  For $m$ equally weighted reviewers whose error correlation matrix is $C$, the variance of the mean error is
\[
\frac{\mathbf 1^TC\mathbf 1}{m^2}.
\]
The equivalent number of independent reviewers is therefore
\[
\boxed{
 m_{\mathrm{eff}}
 =
 \frac{m^2}{\mathbf 1^TC\mathbf 1}.
}
\]
For equicorrelation $\rho$,
\[
 m_{\mathrm{eff}}
 =
 \frac{m}{1+(m-1)\rho}.
\]
The formula does not diminish expert consensus.  It explains why diversity of method and representation matters as much as the number of readers.

\begin{figure}[H]
\centering
\includegraphics[width=.86\linewidth]{figures/synthesis/03_effective_reviewer_count.png}
\caption{Nominal review capacity can grow while independent review capacity saturates.  Different representations and methods reduce shared error more effectively than repeated inspection from one viewpoint.}
\end{figure}

\subsection{The no-universal-view principle}

Resolution Theory should not be misread as a search for one perfect decomposition.  In many problems, no single resolution simultaneously diagonalizes every operator, preserves every bridge, minimizes every cost, and answers every question.  Noncommuting projector atlases are the mathematical version of this fact.  Proof portfolios are the methodological version.  Multiple models in statistics and multiple intermediate representations in a compiler are further instances.

The correct general principle is:
\[
\boxed{
\text{seek coherent transitions among useful views, not one view that erases all others.}
}
\]

\subsection{A human protocol}

A resolution-aware workflow can be stated without specialized vocabulary.
\begin{enumerate}[label=\arabic*.]
\item Begin with a real question, not a favored representation.
\item State how the object is currently being divided or observed.
\item Record which transfers or interactions the representation makes visible.
\item Separate exact statements from approximations and conjectures.
\item Measure the residual rather than assuming omitted structure is harmless.
\item Refine where the residual, uncertainty, or future value justifies the cost.
\item Compress for the intended audience, but preserve a richer archive.
\item Invite review from genuinely different methods and representations.
\item Publish code, data, assumptions, and failure cases beside the claim they support.
\item Reopen the representation when a new question makes an old compression inadequate.
\end{enumerate}

\begin{figure}[H]
\centering
\includegraphics[width=.95\linewidth]{figures/synthesis/06_resolution_aware_practice_loop.png}
\caption{Mathematical work is not a straight line from question to theorem.  It is a cycle of representation, calculation, residual measurement, refinement, verification, compression, and new questions.}
\end{figure}

\subsection{What this combined theory does not say}

The theory does not say that truth is subjective, that every representation is equally useful, or that a long archive is automatically better than a short proof.  It does not turn numerical evidence into proof or formal proof into physical meaning.  It does not imply that every problem has a sparse or efficient resolution.  It does not remove the need for taste, judgment, or domain knowledge.

It says something more restrained: representation has consequences, and those consequences can often be modeled.  Refinement has a cost.  Compression has a residual.  Evidence has a type.  Review has a covariance structure.  Meaning and future bridge value are not captured completely by formal validity, even though formal validity remains indispensable.

\subsection{The enlarged closure theorem}

\begin{theorem}[Resolution-aware practice closure]
Let a finite mathematical construction be generated from:
\begin{enumerate}[label=\textup{(\roman*)}]
\item exact linear packet transfers;
\item resolved nonlinear jets and interaction trees;
\item compatible refinement and chart changes;
\item certified compression and pruning;
\item finite Resolution Logic rules;
\item typed evidence-ledger merges;
\item proof-path compression maps with named residual data.
\end{enumerate}
Then the construction has a well-defined resolved semantic value, a finite transfer or interaction tree, a composite quantitative grade, a typed evidence ledger, and a traceable compression history.
\end{theorem}

\begin{proof}
Proceed by induction on the construction tree.  Exact leaves have a semantic value, zero residual, and a typed evidence entry.  Linear, multilinear, tensor, refinement, and chart-change nodes are handled by the closure theorems proved earlier.  Compression and pruning add their named residuals.  Sequential and parallel certificate rules produce the composite grade.  Evidence ledgers merge componentwise, so their types are preserved.  Compression histories compose as ordinary maps between proof-path representations.  Every generating node therefore preserves the five claimed structures, and the same holds at the root.
\end{proof}

The theorem is finite.  The ordinary resolved $C^*$-algebraic representation problem has now been settled, but the theorem does not provide an infinite topos semantics, a monoidal fiber functor in complete generality, a higher-categorical realization, or a machine-independent measure of meaning.  It states exactly what the present theory can close.

\subsection{Final answer to the opening question}

What does mathematics mean?  It is the search for structures and conclusions that survive justified changes of representation.

What does computer science mean?  It is the art and science of making representations executable under finite resources.

What does statistics mean?  It is the discipline of learning and deciding while keeping uncertainty, evidence, and observational resolution visible.

What does technology mean?  It is the moment those representations and decisions are embodied in instruments, software, institutions, and physical systems.

The common thread is not that the four fields are identical.  It is that each turns an inaccessible whole into a workable representation.  The responsibility of a mature practice is to know what that representation preserves, what it hides, what it costs, and how to change it when the question changes.

The deepest claim of Resolution Theory is therefore not that one decomposition explains everything.  It is that the act of resolving an object---and the discipline of moving among resolutions---deserves to be studied as mathematics in its own right.
\clearpage
\section{Reference implementations}
\label{sec:resolution-reference-code}

The source is included in full because the computational layer is part of the audit trail.  The modules are deliberately direct.  They favor a close correspondence with the theorems over aggressive optimization.

\subsection{Linear Resolution Theory}
\lstinputlisting[language=Python,basicstyle=\ttfamily\tiny,breaklines=true]{code/resolution_theory.py}

\clearpage
\subsection{Universal representation theorem}
\lstinputlisting[language=Python,basicstyle=\ttfamily\tiny,breaklines=true]{code/universal_representation.py}

\clearpage
\subsection{Nonlinear Resolution Theory}
\lstinputlisting[language=Python,basicstyle=\ttfamily\tiny,breaklines=true]{code/nonlinear_resolution.py}

\clearpage
\subsection{Resolution Logic}
\lstinputlisting[language=Python,basicstyle=\ttfamily\tiny,breaklines=true]{code/resolution_logic.py}
\clearpage
\subsection{Resolution-aware practice and evidence}
\lstinputlisting[language=Python,basicstyle=\ttfamily\tiny,breaklines=true]{code/resolution_practice.py}


\clearpage
\section{Automated tests}
\label{sec:resolution-tests}

The tests are not substitutes for the proofs.  They check that the finite implementations satisfy the identities and certificate inequalities used by the examples.

\subsection{Linear theory tests}
\lstinputlisting[language=Python,basicstyle=\ttfamily\tiny,breaklines=true]{tests/test_resolution_theory.py}

\clearpage
\subsection{Universal representation tests}
\lstinputlisting[language=Python,basicstyle=\ttfamily\tiny,breaklines=true]{tests/test_universal_representation.py}

\clearpage
\subsection{Nonlinear theory tests}
\lstinputlisting[language=Python,basicstyle=\ttfamily\tiny,breaklines=true]{tests/test_nonlinear_resolution.py}

\clearpage
\subsection{Resolution Logic tests}
\lstinputlisting[language=Python,basicstyle=\ttfamily\tiny,breaklines=true]{tests/test_resolution_logic.py}

\clearpage
\subsection{Resolution-aware practice tests}
\lstinputlisting[language=Python,basicstyle=\ttfamily\tiny,breaklines=true]{tests/test_resolution_practice.py}


\clearpage
\section{Experiment and visual-generation source}
\lstinputlisting[language=Python,basicstyle=\ttfamily\tiny,breaklines=true]{resolution_logic_experiments.py}


\clearpage
\subsection{Final synthesis experiments}
\lstinputlisting[language=Python,basicstyle=\ttfamily\tiny,breaklines=true]{resolution_practice_experiments.py}

\clearpage
\subsection{Universal representation experiments}
\lstinputlisting[language=Python,basicstyle=\ttfamily\tiny,breaklines=true]{universal_representation_experiments.py}

\clearpage
\begin{thebibliography}{9}
\bibitem{GLR1985}
P.~Ghez, R.~Lima, and J.~E.~Roberts,
``W$^*$-categories,''
\emph{Pacific Journal of Mathematics} 120 (1985), 79--109.

\bibitem{Heunen2009}
C.~Heunen,
``An embedding theorem for Hilbert categories,''
\emph{Theory and Applications of Categories} 22 (2009), 321--344.

\bibitem{FritzWesterbaan2020}
T.~Fritz and B.~Westerbaan,
``The universal property of infinite direct sums in $C^*$-categories and $W^*$-categories,''
\emph{Applied Categorical Structures} 28 (2020), 355--365.
\end{thebibliography}
\end{document}
